\documentclass[11pt,a4paper]{article}
% ---------------------------------------------------------------------------
%  FunctionalDiagrams -- reference manual
%  The examples are computed by make_examples.wl (examples/NAME.wl: the code;
%  NAME.tex: the formula; NAME_pictures.tex: the diagrams; NAME.txt: other
%  output); build_manual.sh typesets this file.
% ---------------------------------------------------------------------------
\usepackage[T1]{fontenc}
\usepackage[sc]{mathpazo}
\usepackage[scaled=0.88]{helvet}
\renewcommand{\ttdefault}{txtt}
\usepackage{amsmath,amssymb}
\usepackage{slashed}
\allowdisplaybreaks[2]
\usepackage[a4paper,margin=2.2cm,top=2.4cm,bottom=2.4cm,headheight=22pt]{geometry}
\usepackage{microtype}
\usepackage[dvipsnames]{xcolor}
\usepackage{graphicx}
\usepackage{tikz}
\usetikzlibrary{decorations.pathmorphing, decorations.markings, arrows.meta, calc}
\usepackage{listings}
\usepackage[most]{tcolorbox}
\usepackage{booktabs,array,tabularx,longtable}
\usepackage{titlesec}
\usepackage{enumitem}
\usepackage{fancyhdr}
\usepackage{needspace}
\usepackage{etoolbox}
\usepackage[colorlinks=true,linkcolor=accent,urlcolor=accent]{hyperref}
\hypersetup{pdftitle={FunctionalDiagrams: reference manual},pdfauthor={Nazzareno Africani}}

\definecolor{accent}{HTML}{1F4E79}
\definecolor{accentlight}{HTML}{DCE6F2}
\definecolor{accentpale}{HTML}{EEF3F9}
\definecolor{codebg}{HTML}{F5F7FA}
\definecolor{codeframe}{HTML}{C9D3DF}
\definecolor{keyword}{HTML}{1F4E79}
\definecolor{comment}{HTML}{6A7B8C}
\definecolor{checkgreen}{HTML}{2E7D32}

\titleformat{\section}{\sffamily\Large\bfseries\color{accent}}{\thesection}{0.8em}{}
\titleformat{\subsection}{\sffamily\large\bfseries\color{accent}}{\thesubsection}{0.7em}{}
\titleformat{\subsubsection}{\sffamily\normalsize\bfseries\color{accent}}{\thesubsubsection}{0.6em}{}
\titleformat{\paragraph}[runin]{\sffamily\bfseries\color{accent}}{}{0pt}{}[.]
\titlespacing*{\section}{0pt}{2.6ex plus 1ex minus 0.4ex}{1.2ex}
\titlespacing*{\subsection}{0pt}{2.2ex plus 0.6ex minus 0.3ex}{0.8ex}
\titlespacing*{\subsubsection}{0pt}{1.8ex plus 0.5ex minus 0.3ex}{0.6ex}
\titlespacing*{\paragraph}{0pt}{1.2ex plus 0.4ex}{0.6em}
% a heading is never left alone at the bottom of a page (a subsection directly after its section heading is covered by
% the space asked by the section: a \needspace there would break the page after the section heading)
\makeatletter
\pretocmd{\section}{\needspace{8\baselineskip}}{}{}
\pretocmd{\subsection}{\if@nobreak\else\needspace{6\baselineskip}\fi}{}{}
\makeatother
\setcounter{secnumdepth}{3}
\setcounter{tocdepth}{2}
\setlist{itemsep=1pt, topsep=3pt}
\setlength{\emergencystretch}{1.5em}
\setlength{\parindent}{0pt}
\setlength{\parskip}{4pt plus 1.5pt minus 1pt}

\pagestyle{fancy}
\fancyhf{}
\fancyhead[L]{\sffamily\small\color{accent}FunctionalDiagrams \textperiodcentered{} reference manual}
\fancyhead[R]{\sffamily\small\color{comment}\nouppercase{\rightmark}}
\fancyfoot[C]{\sffamily\small\thepage}
\renewcommand{\sectionmark}[1]{\markboth{\thesection\enspace #1}{\thesection\enspace #1}}
\renewcommand{\headrulewidth}{0.4pt}
\renewcommand{\headrule}{\hbox to\headwidth{\color{codeframe}\leaders\hrule height \headrulewidth\hfill}}

\lstdefinelanguage{WL}{
  morekeywords={IntegralOver,DeltaFunction,GrassmannProduct,DummyVariable,DeclareSource,DeclareGrassmannSource,
    DeclareGrassmannOdd,DeclareGrassmannGraded,DeclareSymmetric,DeclareAntisymmetric,DeclareFunctionalDerivative,
    DeclareDiscreteIndexValues,DeclareField,DeclareLineStyle,DeclareDirectedPropagator,DeclareInteractionLine,DeclareCoupling,
    DeclareDisplayName,DeclareLocalVertex,DeclareBlob,ResetDeclarations,FunctionalDerivative,SimplifyDiagrams,DiagramTerms,SetSourcesToZero,
    FieldDerivative,FieldExpectationValue,GaussianExpectationValue,PerturbativeCorrelationFunction,
    ConnectedCorrelationFunction,ConnectedVacuumDiagrams,OrderPart,ConnectedDiagrams,DropVacuumBubbles,
    ParticleIrreducibleQ,ParticleIrreducibleDiagrams,SkeletonDiagrams,ClosedLoopCount,DropClosedLoops,
    AmputateExternalLegs,LuttingerWardFunctional,SelfEnergy,IrreducibleVertex,PerturbativeFullVertex,
    TwoParticleReducibleQ,ReducibleDiagrams,IrreducibleDiagrams,ParquetDecomposition,BetheSalpeterLadder,
    CheckParquetEquations,DysonSeries,ConnectedFunctionFromVertices,FullPropagator,VertexFunction,MeanField,
    ExternalSource,DrawDiagram,DiagramPlot,DiagramToTikZ,ExpressionToLaTeX,PerturbativeDensityVertex,LabelFunction,
    ExactEquation,CheckEquation,FermiFunction,BoseFunction,DeclareMatsubaraFrequency,MatsubaraSum,LoopMomentum,
    LoopFrequency,MomentumSum,DiagramIntegrand,EvaluateMatsubara,AnalyticContinuation,RetardedImaginaryPart,EnergyDelta,
    CuttingRuleImaginaryPart,CheckCuttingRules,ZeroTemperatureLimit,HubbardStratonovich,IntegrateOutField,
    DiagramLoopIntegral,LoopIntegral,ScalarPropagator,MomentumSquared,RenormalizationScale,IntegrateLoops,TadpoleIntegral,
    GFunction,BubbleIntegral,EpsilonExpand,PolePart,FeynmanParameterIntegral,FeynmanParameter,InfraredRearrangedIntegral,
    MassInsertion,BetaFunctionFromZ,AnomalousDimensionFromZ,GammaMatrix,SlashedMomentum,DiracProduct,DiracTrace,SpinorChain,
    LorentzMetric,FourVector,ScalarProduct,SpacetimeDimension,DiracSimplify,DiracOrder,QEDIntegrand,FeynmanDenominator,
    MinkowskiLoopIntegral,IntegrateOneLoop,SimplexIntegral,ParameterIntegral,ReduceParameterIntegrals,
    CombineParameterIntegrals,ExpandParameterIntegrals,DifferentiateUnderIntegral,VertexFormFactors},
  sensitive=true,
  morecomment=[n]{(*}{*)},
  morestring=[b]",
}
\lstset{language=WL, basicstyle=\ttfamily\small, keywordstyle=\color{keyword}\bfseries, commentstyle=\color{comment}\itshape,
  stringstyle=\color{BrickRed}, columns=fullflexible, keepspaces=true, breaklines=true, breakatwhitespace=true,
  breakindent=2em, showstringspaces=false, upquote=true}
\newtcblisting{code}{listing only, colback=codebg, colframe=codeframe, boxrule=0.5pt, arc=2pt, left=6pt, right=6pt,
  top=3pt, bottom=3pt, listing options={language=WL}}
\lstdefinestyle{plainoutput}{language={}, basicstyle=\ttfamily\small, breaklines=true, breakatwhitespace=false,
  columns=fullflexible, keepspaces=true}
\newcommand{\pkg}[1]{\texttt{#1}}
\newcommand{\ok}{\textcolor{checkgreen}{\ensuremath{\checkmark}}}
\newcommand{\phbar}{\overline{\mathrm{ph}}}
\newcommand{\ph}{\mathrm{ph}}
\newcommand{\pp}{\mathrm{pp}}

% ---------------------------------------------------------------- the examples
\newcommand{\exmissing}[1]{\textcolor{red}{[missing example: \detokenize{#1}]}}
% the code of an example, as it was evaluated, always followed by its output. The two stay together: after the code
% (with the space of "after skip=2pt") come \nobreak and the no-break flag of LaTeX, as after a heading, so that the
% output box, breakable, starts on the same page or, when it does not fit and fewer than "lines before break" lines are
% left, goes to the next page with the code (tcolorbox then clears the flag)
\makeatletter
\newcommand{\excodeafter}{\tcb@parfillskip@check\par\ifvmode\nobreak\@nobreaktrue\nointerlineskip
  \addvspace{\glueexpr2pt-\parskip}\fi}
\makeatother
% a sentence that introduces a code box stays on its page (\par\nobreak when the box follows text), and is never left as
% one line at the top of a page (the widow penalty, in a group, acts on that paragraph only)
\newcommand{\excode}[1]{\ifhmode{\widowpenalty=10000\par}\nobreak\fi\IfFileExists{examples/#1.wl}{\tcbinputlisting{listing only, colback=codebg, colframe=codeframe,
  boxrule=0.5pt, arc=2pt, left=6pt, right=6pt, top=3pt, bottom=3pt, before skip=4pt, after=\excodeafter, listing file={examples/#1.wl},
  listing options={language=WL}}}{\exmissing{#1}}}
% its output: a formula, the diagrams, or text
\newtcolorbox{exoutput}{colback=white, colframe=codeframe, boxrule=0.5pt, arc=2pt, left=24pt, right=6pt, top=2pt, bottom=2pt,
  before skip=2pt, after skip=8pt, breakable, lines before break=10, overlay unbroken and first={\node[anchor=north west, font=\scriptsize\sffamily\color{comment}] at ([xshift=4pt, yshift=-5pt]frame.north west) {out};}}
\newcommand{\exformulainput}[1]{\IfFileExists{examples/#1.tex}{\input{examples/#1.tex}}{\text{\exmissing{#1}}}}
% a formula on one line, or on several (examples/NAME_lines.tex, written by make_index.py when it is long), not broken
% across pages
\newcommand{\exformuladisplay}[1]{\IfFileExists{examples/#1_lines.tex}{\interdisplaylinepenalty=10000
  \begin{gather*}\input{examples/#1_lines.tex}\end{gather*}}{\[\exformulainput{#1}\]}}
\newcommand{\exformula}[1]{\begin{exoutput}\vspace{-6pt}\exformuladisplay{#1}\vspace{-4pt}\end{exoutput}}
\newcommand{\diagramgap}{\hskip 0.5em plus 0.6em minus 0.15em\relax}
\newcommand{\expicturesinput}[1]{\IfFileExists{examples/#1_pictures.tex}{{\rightskip=0pt plus 1fil \parfillskip=0pt plus 1fil
  \baselineskip=0pt \lineskiplimit=7pt \lineskip=7pt \noindent$\displaystyle\input{examples/#1_pictures.tex}$\par}}{\exmissing{#1}}}
\newcommand{\expictures}[1]{\begin{exoutput}\expicturesinput{#1}\end{exoutput}}
\newcommand{\exboth}[1]{\begin{exoutput}\vspace{-6pt}\exformuladisplay{#1}\vspace{-2pt}\expicturesinput{#1}\end{exoutput}}
\newcommand{\extext}[1]{\IfFileExists{examples/#1.txt}{\tcbinputlisting{listing only, colback=white, colframe=codeframe,
  boxrule=0.5pt, arc=2pt, left=24pt, right=6pt, top=2pt, bottom=2pt, before skip=2pt, after skip=8pt, listing file={examples/#1.txt},
  listing options={style=plainoutput}, breakable, lines before break=10, overlay unbroken and first={\node[anchor=north west, font=\scriptsize\sffamily\color{comment}] at ([xshift=4pt, yshift=-5pt]frame.north west) {out};}}}{\exmissing{#1}}}

% ---------------------------------------------------------------- reference entries
% \begin{function}{Name} syntax lines \end{function}: the head of the entry of a function, with a label for the index
\newtcolorbox{fnbox}{colback=accentpale, colframe=accentpale, arc=2pt, left=6pt, right=6pt, top=3pt, bottom=3pt,
  before skip=2pt, after skip=3pt, fontupper=\small\raggedright}
% (directly after a heading the entry stays with it; elsewhere it goes to the next page if less than four lines are left)
\makeatletter
\newenvironment{function}[1]{\par\if@nobreak\else\medskip\needspace{4\baselineskip}\fi\noindent\phantomsection\label{fn:#1}%
  \begin{fnbox}}{\end{fnbox}\par\nobreak}
\makeatother
\newcommand{\syntax}[1]{\texttt{#1}\par}
\newcommand{\fnref}[1]{\hyperref[fn:#1]{\pkg{#1}}}
\newcommand{\opt}[1]{\texttt{"#1"}}

\begin{document}

% ------------------------------------------------------------------ title
\thispagestyle{empty}
\begin{tikzpicture}[remember picture, overlay]
  \fill[accent] (current page.north west) rectangle ([yshift=-6.0cm]current page.north east);
  \node[anchor=north west, text=white, font=\sffamily\bfseries\fontsize{34}{38}\selectfont] at ([xshift=2.2cm, yshift=-1.45cm]current page.north west) {FunctionalDiagrams};
  \node[anchor=north west, text=white, font=\sffamily\Large, text width=12.2cm] at ([xshift=2.2cm, yshift=-2.9cm]current page.north west)
    {Reference manual\\[2pt]\large functional derivatives and Feynman diagrams in Mathematica};
  \node[anchor=north west, text=accentlight, font=\sffamily\normalsize] at ([xshift=2.2cm, yshift=-4.7cm]current page.north west)
    {Nazzareno Africani \textperiodcentered{} SISSA \textperiodcentered{} version 2.4.0 \textperiodcentered{} October 2026};
  \IfFileExists{../Logo/FunctionalDiagrams_logo_white.pdf}{\node[anchor=north east] at ([xshift=-2.0cm, yshift=-1.2cm]current page.north east)
    {\includegraphics[height=3.6cm]{../Logo/FunctionalDiagrams_logo_white.pdf}};}{}
\end{tikzpicture}
\vspace*{3.6cm}

\noindent This manual describes every function, option and convention of the package \textbf{FunctionalDiagrams}.
The examples were computed by the package itself (\pkg{Documentation/make\_examples.wl}): the code shown in a grey box
was evaluated, and the box marked \textsf{\scriptsize out} after it shows what the code gives, as the package writes it:
a sum of diagrams as \fnref{DiagramPlot} draws it, with the free labels of the output as the external labels, at the size
of this page (\fnref{DiagramToTikZ} with \pkg{"Scale" -> 0.6}, \pkg{"LineWidth" -> 0.55} and \pkg{"LabelFont" ->
"\textbackslash\textbackslash scriptsize"}, which only set the size and are left out of the code); a formula as
\fnref{ExpressionToLaTeX} writes it, with the usual names of the functions (the text gives them where they differ from
the names in the code) and $\varepsilon$ for \pkg{eps}; any other result in \pkg{InputForm}. The examples of a section
build on those before it, in one session. A guided tour of the package is \pkg{Examples.nb}; a presentation, with $\phi^4$ theory and the
electron gas worked out to third order at every step, is \pkg{Showcase/FunctionalDiagrams\_Showcase.pdf}.

\medskip
{\small\setlength{\parskip}{1pt}\tableofcontents}

\clearpage
% ------------------------------------------------------------------ 1
\section{Overview}\label{sec:overview}

\subsection{What the package does}

FunctionalDiagrams is a Mathematica package for the functional calculus of quantum field theory and many-body physics,
and for the Feynman diagrams that it generates. It has five layers.

\begin{itemize}
\item \textbf{An engine for functional derivatives} (sections~\ref{sec:expressions}--\ref{sec:derivatives}). Expressions
  are sums of integrals, over labels, of products of functions; some factors are Grassmann odd and their order matters.
  The engine differentiates such expressions with respect to commuting and anticommuting sources, with the product and
  chain rules, the Grassmann signs and rules declared for objects that depend on the sources implicitly; it integrates
  the delta functions that the derivatives produce, and it brings every term to a canonical form, in which terms that
  differ only by the names of their integration variables are recognized as equal and combined.
\item \textbf{Generating functionals and diagrams} (sections~\ref{sec:fields}--\ref{sec:classes}). Fields are
  derivatives with respect to their sources; expectation values follow by differentiating a generating functional, and
  for Gaussian functionals by Wick's theorem. Perturbation theory, connected and vacuum diagrams, irreducibility,
  skeletons and amputation are all built on the engine. The terms are diagrams: every integration variable is a vertex,
  every two-label factor a line.
\item \textbf{Functionals and their derivatives} (sections~\ref{sec:conventions}--\ref{sec:trees}). The functionals of
  many-body theory and field theory --- $W=\log Z$, the 1PI effective action $\Gamma[\bar\phi]$, the 2PI effective action
  $\Gamma[\bar\phi,G]$, the Luttinger--Ward functional $\Phi[G]$, the functional $\Phi[\mathcal G]$ of a superfield, Potthoff's
  functionals --- are objects of the package, written with inert traces of operators. Their derivatives with respect to
  sources, means and two-point objects are the self-energy and the irreducible vertex, the full vertex and the vertex
  functions, the fully irreducible vertex and the kernels of the two-particle channels; Bethe--Salpeter ladders, the
  parquet decomposition, the exact equations between full objects and the tree expansion of connected functions in
  vertex functions check and relate them.
\item \textbf{Pictures and formulas} (sections~\ref{sec:pictures}--\ref{sec:latex}). Every expression can be drawn, as
  Mathematica graphics or as TikZ code for \LaTeX{}, in a standard style computed from the diagram; formulas can be
  written as \LaTeX{}.
\item \textbf{Evaluation} (sections~\ref{sec:matsubara}--\ref{sec:qed}). For a translation-invariant system at finite
  temperature, the diagrams are evaluated in momentum and Matsubara frequency: momenta and frequencies are routed through
  the lines, the frequency sums are carried out by residues, the result is continued to real frequency, and its imaginary
  part is checked against the finite-temperature cutting rules. At zero temperature, in Euclidean space, the loop
  integrals of the diagrams are evaluated in dimensional regularization and expanded in $\varepsilon$, and the
  renormalization constants and renormalization-group functions follow (the $\overline{\rm MS}$ renormalization of
  $\phi^4$ theory to two loops). In QED at one loop, the diagrams with spinor and Lorentz indices become products and
  traces of Dirac matrices in $d$ dimensions, their loop integrals are done by Feynman parameters, and the electron
  self-energy, the vacuum polarization and the vertex give the renormalization constants, $\beta(e)$ and the form
  factors of the vertex.
\end{itemize}

The guiding principle is that everything is generated by functional derivatives. The user writes a functional --- $Z[J]$,
$W=\log Z$, $\Gamma[\bar\phi]$, $\Gamma[\bar\phi,G]$, $\Phi[G]$, $\Phi[\mathcal G]$ --- and obtains every object by
\fnref{FunctionalDerivative} with respect to sources, means and propagators (normal and anomalous); functions such as
\fnref{SelfEnergy}, \fnref{IrreducibleVertex}, \fnref{PerturbativeFullVertex}, \fnref{PerturbativeDensityVertex} or the option
\opt{Channel} are shortcuts for such derivatives, and this manual shows the derivative first and the shortcut as the
check. In particular, weights are never put in by hand: symmetry factors and fermionic signs come out of the algebra, and
the package only has to recognize when two terms are the same diagram.

\subsection{Requirements and installation}\label{sec:install}

The package is a single file, \pkg{FunctionalDiagrams.wl}, written for Mathematica~12.1 or later (for the
paclet manager that loads it).
The recommended installation is the paclet of a release, which also contains the showcase, this manual and the guided
tour:
\begin{code}
PacletInstall["/path/to/FunctionalDiagrams-2.4.0.paclet"]   (* once *)
Needs["FunctionalDiagrams`"]
\end{code}
(\pkg{PacletInstall} also takes the address of the paclet on the page of the release.) This manual describes
version~2.4.0, the latest release. A clone of the repository is a
paclet too, loaded by \pkg{PacletDirectoryLoad} of its folder and then \pkg{Needs}.
Without the paclet, load the file with
\begin{code}
Get["/path/to/FunctionalDiagrams/FunctionalDiagrams.wl"]
\end{code}
or copy it into the folder \pkg{Applications} of \pkg{\$UserBaseDirectory} and load it with
\begin{code}
Needs["FunctionalDiagrams`"]
\end{code}
The functions are in the context \pkg{FunctionalDiagrams`}; \pkg{?FunctionalDiagrams`*} lists them and \pkg{?Name}
shows their usage. \pkg{\$FunctionalDiagramsVersion} is the version, \pkg{"2.4.0"}. Loading the
package again in the same session starts afresh: the definitions of the package and its record of declarations are
cleared (attributes that the declarations gave to your own symbols, such as \pkg{Orderless} for a symmetric function,
remain).

When the package is loaded, it checks that none of its public names is already a symbol of the context \pkg{System`}
(a built-in function of a later version of Mathematica), which would be used in place of the package's function, and
issues \pkg{FunctionalDiagrams::clash} if one is. A name of the package that another package loaded in the same
session also defines, or that was typed before the package was loaded (a symbol in \pkg{Global`}), gives instead
Mathematica's own message \pkg{::shdw} (\pkg{DiagramPlot::shdw}, say): the function used is that of the context loaded
last, which comes first in \pkg{\$ContextPath}.
In a notebook the package then prints a banner two lines high: the logo, the attribution ``FunctionalDiagrams, by Nazzareno
Africani'' (which modified versions must keep), the version, the license, how to cite the package
(section~\ref{sec:citing}) and links to the repository and to this manual. In a script (\pkg{wolfram -script}) nothing
is printed.

In a notebook, loading the package also checks for a newer release, at most once a day: it asks GitHub
(\pkg{api.github.com}) for the latest release in the background, so that loading does not wait for the answer, and if
that release is newer than \pkg{\$FunctionalDiagramsVersion}, the banner gets a third line with its version, a link to
what is new and the \pkg{PacletInstall} line of its paclet, ready to copy. The date of the last check is kept in the
file \pkg{LastUpdateCheck.txt} of the folder \pkg{ApplicationData/FunctionalDiagrams} of \pkg{\$UserBaseDirectory}.
Every failure (no connection, a proxy, the rate limit of GitHub) is silent, and a script never checks, nor does a
session in which Mathematica may not use the Internet (\pkg{\$AllowInternet}). \pkg{\$FunctionalDiagramsUpdateCheck}
is \pkg{True}; set to \pkg{False} before the package is loaded, it turns the check off:
\begin{code}
$FunctionalDiagramsUpdateCheck = False;
Needs["FunctionalDiagrams`"]
\end{code}
(in the file \pkg{init.m} of the folder \pkg{Kernel} of \pkg{\$UserBaseDirectory}, for every session).

\begin{function}{FunctionalDiagramsUpdateCheck}
\syntax{FunctionalDiagramsUpdateCheck[]}
\end{function}
Asks GitHub for the latest release now, whatever the date of the last check and the value of
\pkg{\$FunctionalDiagramsUpdateCheck}, and waits for the answer (up to ten seconds). It gives
\pkg{<|"InstalledVersion" -> \dots, "LatestVersion" -> \dots, "InstallLine" -> \dots|>}: the version loaded, the
latest release and the line that installs its paclet, as a string; when the release is newer, it also prints that line.
When the release cannot be read it gives \pkg{\$Failed}, with the message \pkg{FunctionalDiagramsUpdateCheck::latest}.

\subsection{A first session}

The free fermion propagator $G_0$, a two-body interaction $V$, the Green function to first order, its self-energy and a
picture:
\begin{code}
DeclareGrassmannSource[eta, etabar];
DeclareField[psi, etabar]; DeclareField[psibar, eta, -1];
DeclareSymmetric[V]; DeclareDirectedPropagator[G0]; DeclareInteractionLine[V];
Zf = Exp[-IntegralOver[G0[x, y] etabar[x] ** eta[y], {x, y}]];
CoulombVertex = 1/2 IntegralOver[V[x, y] psibar[x] ** psibar[y] ** psi[y] ** psi[x], {x, y}];
greenFunction = -PerturbativeCorrelationFunction[Zf, {psi[x], psibar[y]}, CoulombVertex, 1]
DiagramPlot[greenFunction, {x, y}]
\end{code}
The output and the pictures of this session are in section~\ref{sec:perturbation}, where the same computation is
evaluated step by step.

\subsection{Scope and limitations}

\begin{itemize}
\item The package computes with diagrams in the space of labels. A label stands for whatever the problem needs ---
  position, imaginary time, spin, orbital --- and an integral is the integral or sum over it. Nothing is assumed about
  the propagators and interactions beyond their declared symmetries. Its results are exact symbolic expressions, sums of
  integrals of products of functions, which can be translated into any representation; for translation-invariant systems
  at finite temperature the package does it itself (section~\ref{sec:matsubara}), with the frequency sums carried out and
  the momentum integrals left symbolic; at zero temperature it evaluates the loop integrals that reduce to nested
  one-loop ones in dimensional regularization (section~\ref{sec:dimreg}), and the one-loop diagrams of QED, with their
  Dirac algebra (section~\ref{sec:qed}). It is not a program for the numerical evaluation of Feynman integrals.
\item Expectation values can be computed in any generating functional, by differentiation. Perturbation theory is
  fastest around a Gaussian free generating functional (any number of bosonic and fermionic fields), by Wick's theorem;
  around any other $Z_0$ it works by differentiation, more slowly. The interaction is polynomial in the
  fields, with any number of vertices of any kind and coupling constants as symbols.
\item The number of diagrams grows factorially with the order. With Mathematica the electron gas to third order --- 87
  diagrams of the Green function, 94 of the full vertex --- takes a few seconds (in the showcase, \pkg{values.txt}: about 0.1~s for $G$ and a second for $F$); higher orders are limited by time
  and memory only.
\item The canonical form is exact: two terms get the same canonical form if and only if they are the same diagram (up to
  the renaming of integration variables, the declared symmetries and the reordering of Grassmann factors). The algorithm
  can be exponential for very symmetric terms, but it is fast for diagrams.
\item The pictures are computed by heuristic rules, designed for the diagrams of many-body theory (fermion lines with
  density-type vertices) and for trees of vertex functions; other diagrams are drawn by a general embedding, which is
  readable but not always as in a textbook.
\end{itemize}

\subsection{Citing the package}\label{sec:citing}

The package is released under the GNU General Public License, version 3 or later, with two additional terms: the
attribution to the author must be kept, and modified versions must say that they are modified, by whom and when (the
file \pkg{NOTICE}, with the copyright and the terms; the text of the license is in \pkg{LICENSE}). If you use the package
in published work, please cite it, for instance as
\begin{quote}
N.~Africani, \emph{FunctionalDiagrams: functional derivatives and Feynman diagrams in Mathematica}, version 2.4.0 (2026).
\end{quote}
The file \pkg{CITATION.cff} in the repository has the same information in a machine-readable form.

\subsection{Notation used in this manual}

Function names are in \pkg{typewriter}; \pkg{f[x, y]} is Mathematica input, $f(x,y)$ the corresponding mathematical
object. In the formulas, $\int_{z_1z_2}$ is \pkg{IntegralOver[\dots, \{z1, z2\}]}, the integral or sum over everything a
label stands for (position, imaginary time, spin, \dots). A source is written $J$ (bosonic) or $\eta$, $\bar\eta$
(Grassmann), a field $\phi$, $\psi$, $\bar\psi$. In diagrams, an arrow from $y$ to $x$ is a propagator $G(x,y)$, a wavy
line an interaction $V$.

% ------------------------------------------------------------------ 2
\section{Expressions}\label{sec:expressions}

\subsection{Terms}

Every expression the package works with is a sum of \emph{terms}
\[
c\;\int_{z_1\cdots z_n}f\;\theta_1\theta_2\cdots\theta_k,
\]
in Mathematica \pkg{c IntegralOver[f GrassmannProduct[$\theta_1$, \dots, $\theta_k$], \{z1, \dots, zn\}]},
where $c$ does not depend on the integration variables, $f$ is a product of commuting factors (functions of labels,
symbols, numbers, functions of other expressions), and $\theta_1\cdots\theta_k$ is an ordered product of Grassmann-odd
factors. A term is a diagram: its two-label factors are lines between its labels, its integration variables are internal
vertices, its other labels are external points. Any expression built from these pieces with \pkg{Plus}, \pkg{Times},
\pkg{**}, powers and functions is accepted as input; the functions of the package expand it into terms internally and
return a canonical sum.

\subsection{Notation}

\begin{function}{IntegralOver}
\syntax{IntegralOver[f, \{x, y, \dots\}]}
\end{function}
The integral (or sum) of $f$ over the labels $x,y,\dots$, which stand for position, imaginary time, spin, orbital or any
other index. The integration variables must be symbols. Integrals can be nested, and the integrand can contain anything,
including other integrals in the arguments of functions. It is displayed as $\int_{x,y}f$. In results the integration
variables are canonical: \pkg{DummyVariable[1, k]}, displayed $z_k$, in the outermost integrals,
\pkg{DummyVariable[2, k]}, displayed $w_k$, in integrals inside them, and so on ($u$, $v$, $t$).
\excode{canonical_names}
\exformula{canonical_names}

\begin{function}{DeltaFunction}
\syntax{DeltaFunction[x, y]}
\end{function}
The delta function $\delta(x,y)=\delta(x-y)$ of two labels. It is symmetric (\pkg{Orderless}), and it is integrated
automatically against integration variables:
\excode{delta_integrated}
\exformula{delta_integrated}
Between two labels that are not integrated over it is used to simplify: $\delta(a,b)\,f(a,b)=\delta(a,b)\,f(a,a)$, the
labels of a group joined by delta functions being replaced by the first of them in canonical order. Diagrams with
external points joined by delta functions (the vertices of amputated functions) are thereby recognized as equal.
\excode{delta_external}
\exformula{delta_external}
With discrete values declared (\fnref{DeclareDiscreteIndexValues}), \pkg{DeltaFunction[up, up]} is 1 and
\pkg{DeltaFunction[up, down]} is 0.

\begin{function}{GrassmannProduct}
\syntax{GrassmannProduct[a, b, \dots]}
\syntax{a ** b ** \dots}
\end{function}
The ordered product of Grassmann-odd factors. Products entered with \pkg{**}, that is with
\pkg{NonCommutativeMultiply}, are converted to it. It is displayed as $a\cdot b\cdots$. Grassmann-odd factors are the Grassmann sources, the fields
attached to them, and the heads declared with \fnref{DeclareGrassmannOdd} or \fnref{DeclareGrassmannGraded}. A product of
two odd factors entered with \pkg{Times} is ambiguous, since \pkg{Times} is commutative: it gives a message
(\pkg{FunctionalDerivative::times}) and the evaluation is aborted; such products must be written with \pkg{**}. An odd
factor squared, \pkg{eta[x]\^{}2}, vanishes. Commuting factors can be multiplied
either way. In canonical form the odd factors are sorted, with the sign of the permutation, and a product with a
repeated odd factor vanishes.

\begin{function}{DummyVariable}
\syntax{DummyVariable[depth, k]}
\end{function}
The $k$-th canonical integration variable at nesting depth \pkg{depth}, displayed $z_k$, $w_k$, $u_k$, $v_k$, $t_k$ for
depths 1 to 5. These are the names given to the integration variables in results; they can be used in input like any
other symbol.
\excode{nested_integrals}
\exformula{nested_integrals}

\subsection{Display}

In the output, \pkg{IntegralOver}, \pkg{DeltaFunction}, \pkg{GrassmannProduct} and \pkg{DummyVariable} are displayed in a
compact notation that remains valid input (\pkg{Interpretation}); \fnref{DeclareDisplayName} does the same for functions of
your own. \fnref{ExpressionToLaTeX} writes any expression as \LaTeX{}, as in the boxes of this manual.

% ------------------------------------------------------------------ 3
\section{Declarations}\label{sec:declarations}

Declarations tell the package which heads are sources, fields, symmetric or Grassmann-odd functions, which objects
depend on the sources implicitly, and how lines are drawn. They are global, they accumulate, and
\fnref{ResetDeclarations} removes them all. The declarations of heads warn, with the message \pkg{FunctionalDiagrams::builtin},
if a name is a built-in symbol of the Wolfram Language (\pkg{C}, \pkg{D}, \pkg{E}, \pkg{I}, \pkg{K}, \pkg{N}, \pkg{O}, \pkg{Gamma}, \pkg{Beta},
\pkg{GreenFunction}, \dots): the definitions of built-in symbols interfere with those of the package, and other names
must be used, such as \pkg{BareGreenFunction} and \pkg{DressedGreenFunction}.

\subsection{Sources and fields}

\begin{function}{DeclareSource}
\syntax{DeclareSource[J1, J2, \dots]}
\end{function}
Declares commuting (bosonic) source heads: \pkg{J[x]} is then differentiated by \fnref{FunctionalDerivative}, and set to
zero by \fnref{SetSourcesToZero}. A source can have any number of arguments, \pkg{J[x]} or \pkg{J[x, s]}; the derivative
of \pkg{J[x, s]} with respect to \pkg{J[y, t]} is \pkg{DeltaFunction[x, y] DeltaFunction[s, t]}.

\begin{function}{DeclareGrassmannSource}
\syntax{DeclareGrassmannSource[eta, etabar, \dots]}
\end{function}
Declares anticommuting (Grassmann) source heads. Products of Grassmann sources must be written with \pkg{**}.

\begin{function}{DeclareField}
\syntax{DeclareField[phi, J]}
\syntax{DeclareField[psibar, eta, -1]}
\end{function}
Associates the field \pkg{phi} with the source \pkg{J}: \pkg{phi[x]} stands for $\delta/\delta J(x)$, and with the third
argument $s$ for $s\,\delta/\delta J(x)$. For the fermionic convention of this manual,
\pkg{DeclareField[psi, etabar]} and \pkg{DeclareField[psibar, eta, -1]}: $\psi\leftrightarrow\delta/\delta\bar\eta$,
$\bar\psi\leftrightarrow-\delta/\delta\eta$. Fields attached to Grassmann sources are Grassmann odd. Fields are used in
\fnref{FieldDerivative}, in expectation values and in interactions (section~\ref{sec:perturbation}).

\subsection{Symmetries and Grassmann parity}

\begin{function}{DeclareSymmetric}
\syntax{DeclareSymmetric[h1, h2, \dots]}
\end{function}
Makes the heads totally symmetric in their arguments (they get the attribute \pkg{Orderless}): bosonic propagators,
density--density interactions, totally symmetric vertices. The canonical form takes the symmetry into account.
\excode{symmetric}
\exformula{symmetric}

\begin{function}{DeclareAntisymmetric}
\syntax{DeclareAntisymmetric[h1, h2, \dots]}
\end{function}
Makes the heads totally antisymmetric: the arguments are sorted with the sign of the permutation, and a function with a
repeated argument vanishes. In a pattern or a rule write \pkg{h[\_\_]}, never \pkg{h[\_, \_]}, which evaluates to 0. The
derivative with respect to an antisymmetric kernel is antisymmetrized (\fnref{FunctionalDerivative}), and the kernel of the
inverse of an antisymmetric operator, \pkg{OperatorKernel[h\^{}-1]}, is antisymmetric: so is the propagator of a superfield,
the Nambu matrix of $(\psi,\bar\psi)$ (\fnref{SuperfieldComponents}).
\excode{antisymmetric}
\extext{antisymmetric}

\begin{function}{DeclareGrassmannOdd}
\syntax{DeclareGrassmannOdd[c1, c2, \dots]}
\end{function}
Declares heads of Grassmann-odd quantities that do not depend on the sources, such as fermion fields written explicitly
in an expression. They anticommute with each other and with the Grassmann sources.
\excode{grassmann_odd}
\exformula{grassmann_odd}

\begin{function}{DeclareGrassmannGraded}
\syntax{DeclareGrassmannGraded[h1, h2, \dots]}
\end{function}
Declares heads whose Grassmann parity is the number of their arguments modulo 2 and which depend on the sources: their
odd instances vanish at zero source. This is the grading of superfields; \fnref{ConnectedFunctionFromVertices} uses it for
the tree expansion of fermions (superfield rules).

\begin{function}{DeclareDiscreteIndexValues}
\syntax{DeclareDiscreteIndexValues[\{up, down, \dots\}]}
\end{function}
Declares symbols that are values of a discrete index rather than labels: \pkg{DeltaFunction[up, up]} is 1,
\pkg{DeltaFunction[up, down]} is 0, and the drawing and the cuts ignore them. Sum over such an index explicitly, with
\pkg{Sum}: for spin-$\frac12$ fermions, with the fields \pkg{psi[x, s]} and \pkg{psibar[x, s]},
\begin{code}
Zf = Exp[-Sum[IntegralOver[G0[x, y] etabar[x, s] ** eta[y, s], {x, y}],
    {s, {up, down}}]];
interaction = 1/2 Sum[IntegralOver[V[x, y] psibar[x, s] ** psibar[y, t] **
    psi[y, t] ** psi[x, s], {x, y}], {s, {up, down}}, {t, {up, down}}];
\end{code}
The external fields carry
explicit values, \pkg{psi[a, up]}, not a symbolic spin label. Do not integrate over a
spin label with \pkg{IntegralOver}: the closed loops would leave traces $\int_s\delta(s,s)$, which the classes of
diagrams do not count. Lines are functions of two labels: a propagator that depends on spin has one head per spin
(\pkg{Gup}, \pkg{Gdown}); a line head with more arguments is treated as a vertex, with the message
\pkg{FunctionalDiagrams::lineargs}. With one head for both spins (\pkg{G0 -> G} in \fnref{LuttingerWardFunctional}),
$\delta\Phi/\delta G$ is the sum of the derivatives with respect to $G_\uparrow$ and $G_\downarrow$, taken at
$G_\uparrow=G_\downarrow=G$: twice the self-energy of one spin, $(2s+1)\Sigma$. The self-energy of one spin, and the vertex
with given spins, are derivatives with respect to one head per spin, \pkg{\{G0up -> Gup, G0down -> Gdown\}}:
$\Sigma_\uparrow(x,y)=\delta\Phi/\delta G_\uparrow(y,x)$, $\Gamma^0_{\uparrow\downarrow}=\delta\Sigma_\uparrow/\delta G_\downarrow$ (the direct
interaction only). A field without the spin index of its source (\pkg{psi[x]} when $Z_0$ holds \pkg{eta[x, s]}) has vanishing
expectation values: the message \pkg{FunctionalDiagrams::fieldarity}.
\excode{discrete_values}
\extext{discrete_values}

\subsection{Implicit dependence on the sources}

\begin{function}{DeclareFunctionalDerivative}
\syntax{DeclareFunctionalDerivative[object[a\_, \dots], source[c\_], rhs]}
\end{function}
Declares $\delta\,\mathrm{object}/\delta\,\mathrm{source}=\mathrm{rhs}$ for an object that depends on the sources implicitly,
for instance a propagator, a mean field or a vertex function in the presence of a source. The first two arguments are
patterns; \pkg{rhs} is evaluated for every match. Integration variables in \pkg{rhs} are renamed automatically, so that
the rule can be used several times in one term. Once declared, the rule is used by \fnref{FunctionalDerivative} whenever
it meets the object, including inside other functions (chain rule); the object is treated as depending on the sources of
that head. Several rules can be declared for the same object and different sources. The rules are removed by
\fnref{ResetDeclarations}. Section~\ref{sec:derivatives} has an example; \fnref{ConnectedFunctionFromVertices} uses the same mechanism
internally.

\subsection{Pictures and display}

\begin{function}{DeclareLineStyle}
\syntax{DeclareLineStyle[h, style]}
\end{function}
Sets how a two-label factor \pkg{h[x, y]} is drawn. The styles are \opt{Plain} (the default), \opt{Arrow} (a propagator
running from \pkg{y} to \pkg{x}), \opt{DoubleArrow} (a dressed propagator), \opt{DottedArrow} (a ghost), \opt{Double},
\opt{Wavy}, \opt{DoubleWavy} (a dressed interaction, e.g.\ the screened $W$ beside $V$), \opt{Curly} (a gluon), \opt{Dashed},
\opt{Dotted}, \opt{Zigzag}, \opt{Impurity} (dashed, with a cross), and for the anomalous propagators of a superfield
\opt{DivergingArrows} ($-\langle T\psi\psi\rangle$, two arrowheads pointing to its ends) and \opt{ConvergingArrows}
($-\langle T\bar\psi\bar\psi\rangle$, pointing to its middle). The seven from \opt{Wavy} to \opt{Impurity} are bosonic or
interaction lines: they are not part of fermion lines, and the layout treats
them as such (section~\ref{sec:layout}). As in the textbooks, a skeleton expansion, in which every line is a dressed one,
is drawn with the plain styles; the double ones mark the dressed propagator where the bare one appears too, as in a Dyson
equation or in $W=V+VPW$. The options \opt{Directed}, \opt{Interactions} and \opt{InteractionStyle}
(section~\ref{sec:options}) change the style of a head in one picture only. An unknown style gives the message
\pkg{DeclareLineStyle::style}.

\begin{function}{DeclareDirectedPropagator}
\syntax{DeclareDirectedPropagator[G1, G2, \dots]}
\end{function}
The same as \pkg{DeclareLineStyle[G, "Arrow"]} for every head: \pkg{G[x, y]} is drawn with an arrow from \pkg{y} to \pkg{x}
(a fermion propagating from $y$ to $x$).

\begin{function}{DeclareInteractionLine}
\syntax{DeclareInteractionLine[V1, V2, \dots]}
\syntax{DeclareInteractionLine[V1, V2, \dots, style]}
\end{function}
Declares interaction lines, the kernels of the interaction: for the electron gas $V$, in
$S_{\rm int}=\frac12\int_{xy}V(x,y)\,\bar\psi(x)\bar\psi(y)\psi(y)\psi(x)$. They are drawn as wavy lines, as with \pkg{DeclareLineStyle[V, "Wavy"]}, or in the given style, e.g.\
\opt{DoubleWavy} for a screened interaction $W$ drawn beside $V$, and in the weight written with a picture each of them stands for $-V$: the minus sign comes
from expanding $e^{-S_{\rm int}}$ (section~\ref{sec:weights}). A head declared only with \fnref{DeclareLineStyle} is a
propagator, and stands for itself.

\begin{function}{DeclareCoupling}
\syntax{DeclareCoupling[c1, c2, \dots]}
\end{function}
Declares couplings of local interactions, such as $\lambda$ in $S_{\rm int}=\frac{\lambda}{4!}\int_z\phi^4$ or $g$ in
$g\int_z\varphi\bar\psi\psi$. In the weight written with a picture they are absorbed into the vertices, each of which stands
for minus its coupling: $\lambda^n$ in the coefficient of a term becomes $(-1)^n$ (section~\ref{sec:weights}). The
expressions are not changed.

\begin{function}{DeclareLocalVertex}
\syntax{DeclareLocalVertex[h1, h2, \dots]}
\end{function}
Declares heads of local vertices: kernels such as the quark--gluon vertex \pkg{h[x, y, z]} or the three- and four-gluon
vertices, whose labels carry different indices (Lorentz, colour, spinor) at one point of space-time. A picture draws
\pkg{h[x, y, \dots]} as one point at which its lines meet, instead of a grey blob; the cuts already treat any factor of
several labels as one vertex. In the weight written with a picture a local vertex stands for minus its kernel, $-h$, as
an interaction line does (section~\ref{sec:weights}). The one-loop gluon self-energy of QCD, with the four vertices of the theory
(\pkg{"Curly"} gluons, \pkg{"DottedArrow"} ghosts):
\excode{qcd_gluon_self_energy}
\expictures{qcd_gluon_self_energy}

\begin{function}{DeclareBlob}
\syntax{DeclareBlob[h, label, options]}
\end{function}
Draws a factor \pkg{h[x1, \dots, xn]} with $n\ge2$ labels --- also a self-energy \pkg{h[x, y]}, which would otherwise be a
line --- as a blob with \pkg{label} written inside: a LaTeX string such as \pkg{"\textbackslash\textbackslash Sigma"} (in
the graphics of Mathematica, plain text read from it: Greek letters, the scripts on the line, accents and bars above).
Options: \opt{Shape} -> \opt{Circle} (the default), \opt{Box}, \opt{Triangle} (equilateral,
three legs at its corners, two on the left and one on the right: the vertex of a one-body operator), \opt{TriangleUp} (the
same triangle, two legs at the bottom, one at the top) or \opt{Diamond} (a square turned by 45 degrees, legs at its four corners);
\opt{AxialPorts} -> \pkg{\{k1, \dots\}}: the lines of these legs leave horizontally, so that a self-energy is drawn with its
fermion line on an axis (default \pkg{\{1, 2\}} for \opt{TriangleUp});
\opt{Fill} -> \opt{Hatched} (the default), \opt{CrossHatched}, \opt{Dotted} (dots on a staggered grid, sparser than the
hatching), \opt{Grey} or \opt{White}; an unknown value gives the message \pkg{DeclareBlob::opt}. The legs leave the blob counterclockwise from the top left, from the corners of a square
for four legs, and its lines leave from these corners, as from those of a triangle, never from along a side: a point
that is twice an argument (\pkg{h[z, z, w, w]}, the density of a two-particle function) stands outside two corners,
joined to each by a line. A blob is as large as its label needs: the label sits at the centre, or near the incentre of a triangle,
on a hatched, cross-hatched or dotted blob in a white clearing with a band of the fill left around it, and a label with an
accent, a subscript or
a superscript ($\tilde\Lambda$, $\Gamma_{\rm ph}$, $\Gamma^{(4)}$) makes the blob a little larger, never the label
smaller; in one picture the blobs of one shape have one size, so that the lines between them are parallel, and so do
those of all the terms of a sum drawn together, so that a blob looks the same in every term. The size is
estimated from the LaTeX of the label for a label font up to 0.6 units of length (\pkg{"\textbackslash\textbackslash
footnotesize"} at \opt{Scale} -> \pkg{0.55}, the default \pkg{"\textbackslash\textbackslash small"} at \pkg{0.7}). Blobs
joined by lines are drawn side by side, on the side of the legs that join them, and a line that would
cross a blob, or pass close to it, is bent around it. The hatching is drawn as lines and the dots as small discs, so the
TikZ code still needs only \pkg{decorations.pathmorphing}. For the cuts a blob is a vertex, like any factor that is not a line. The exact equations
of many-body theory are drawn in this way, e.g. the right-hand side of the Bethe--Salpeter equation
$F=\Gamma^0+\Gamma^0L_0F$ ($L_0$, the pair of propagators of the channel, section~\ref{sec:conventions}):
\excode{blobs}
\expictures{blobs}
The documents of the package fill the blobs by one rule: a vertex function the more densely the more diagrams it contains
--- the fully irreducible vertex $\Gamma_{\rm irr}$ dotted, the irreducible parts ($\Sigma$, $\Delta$, $P$, $\Gamma^0$,
the $\Gamma_r$ irreducible in one channel, $\Gamma_{\rm xc}$, $\tilde\Lambda$) hatched, the full vertex functions ($F$,
$\Lambda$, $\Gamma^{(4)}$) cross-hatched --- and the correlation and response functions ($W^{(4)}$, $L$, $\chi$) white.
The three fills of the vertices in the vertex irreducible in the particle--particle channel,
$\Gamma_{\pp}=\Gamma_{\rm irr}+R_{\ph}+R_{\phbar}$ (section~\ref{sec:parquet}), with $F$ and the propagator above:
\excode{blobs_fills}
\expictures{blobs_fills}
A box joined by single lines from three of its corners to one vertex is turned by 45 degrees, so that the picture is
symmetric: the vertex lies on the diagonal through the fourth corner, the lines to the two other corners are arcs, and
the label stays upright (the diamond is as large as the square, or larger if the label needs it). The term
$\frac{\lambda}{6}\int G\,G\,G\,\Gamma^{(4)}$ of the self-energy of $\phi^4$ theory (the Schwinger--Dyson equation):
\excode{blobs_turned}
\expictures{blobs_turned}

\begin{function}{DeclareDisplayName}
\syntax{DeclareDisplayName[h, name]}
\end{function}
Displays \pkg{h[x, y, \dots]} as \pkg{name} with the subscript $xy\dots$. After
\begin{code}
DeclareDisplayName[BareGreenFunction, "G0"]
\end{code}
\pkg{BareGreenFunction[x, y]} is displayed $\mathrm{G0}_{xy}$: input keeps the explicit name, output is compact and remains
valid input. \fnref{ExpressionToLaTeX} uses the name too (\pkg{"G0"} becomes $G_{0}$).

\begin{function}{ResetDeclarations}
\syntax{ResetDeclarations[]}
\end{function}
Removes all declarations: sources, fields, Grassmann parities, symmetries, derivative rules, discrete values, line
styles, local vertices and display names. Symmetric heads lose the attribute \pkg{Orderless} and antisymmetric heads
their sign rules; the display formats remain (\pkg{ClearAll[h]} removes them).

% ------------------------------------------------------------------ 4
\section{Functional derivatives and canonical form}\label{sec:derivatives}

\subsection{The derivative}

\begin{function}{FunctionalDerivative}
\syntax{FunctionalDerivative[F, s1, s2, \dots, sn]}
\syntax{FunctionalDerivative[F, \{s1, s2, \dots, sn\}]}
\syntax{FunctionalDerivative[F, G[x, y]]}
\syntax{FunctionalDerivative[F, s1, \dots, sn, "MaximumOrder" -> \{x, n\}]}
\end{function}
The functional derivative
\[
\frac{\delta}{\delta s_1}\frac{\delta}{\delta s_2}\cdots\frac{\delta}{\delta s_n}F ,
\]
the rightmost derivative acting first. The $s_i$ are source atoms such as \pkg{J[x]} or \pkg{eta[x]}, with declared heads
(\fnref{DeclareSource}, \fnref{DeclareGrassmannSource}); their arguments can be any labels, including integration
variables of an enclosing expression. The result is in canonical form (section~\ref{sec:canonical}). The derivative is
computed term by term, with the following rules.
\begin{itemize}
\item \emph{Sources.} $\delta J(x)/\delta J(y)=\delta(x,y)$, with one delta function per argument; sources of other heads
  and objects that do not depend on the source give zero.
\item \emph{Product rule}, for the commuting factors and for the Grassmann-odd factors. Derivatives with respect to
  Grassmann sources are left derivatives: when the derivative reaches the $k$-th odd factor it has passed $k-1$ odd
  factors, and the term gets the sign $(-1)^{k-1}$.
\item \emph{Chain rule}, for powers, exponentials, logarithms and any other function: the derivative of
  \pkg{f[u, v, \dots]} is $\sum_i \partial_i f\,\delta u_i/\delta s$, with $\partial_i f$ written
  \pkg{Derivative[0, \dots, 1, \dots][f][u, v, \dots]} (displayed $f'$, $f''$, \dots{} for functions of one argument).
\item \emph{Integrals}, nested at any depth: the integrand is differentiated, the integration variables being renamed
  first so that they cannot clash with the labels of the source.
\item \emph{Declared rules} (\fnref{DeclareFunctionalDerivative}) for objects that depend on the sources implicitly.
\item \emph{Delta functions} produced by the derivatives are integrated at once against integration variables.
\end{itemize}
Successive derivatives are applied one at a time, each followed by the canonical form, so that equal terms are combined
early. Examples:
\excode{derivative_source}
\exformula{derivative_source}
\excode{derivative_power}
\exformula{derivative_power}
\excode{derivative_exp}
\exformula{derivative_exp}
\excode{derivative_function}
\exformula{derivative_function}
\excode{derivative_nested}
\exformula{derivative_nested}
\excode{derivative_list}
\exformula{derivative_list}
Grassmann sources: the two orders of the derivatives differ by a sign.
\excode{grassmann_sources}
\extext{grassmann_sources}
A declared rule, $\delta G(a,b)/\delta J(c)=G(a,c)\,G(c,b)$ for $G=(G_0^{-1}-J)^{-1}$ (\pkg{PropagatorOfJ}) inside an
integral (the two terms of the second derivative are one diagram, $V$ being symmetric):
\excode{implicit_rule}
\exformula{implicit_rule}
A source can be any function, with any number of arguments (the two-point source $K(x,y)$ of the 2PI action, for
instance).

\paragraph{Derivatives with respect to two-index objects} The propagator need not be declared as a source:
\pkg{FunctionalDerivative[F, G[x, y]]}, with \pkg{G} any kernel of two labels that is not a source --- a propagator, a
self-energy, an interaction $V$, the propagator in a background --- treats every $G(a,b)$ in \pkg{F} as an independent
variable ($\delta G(a,b)/\delta G(x,y)=\delta(a,x)\,\delta(b,y)$, symmetrized for a symmetric \pkg{G}, antisymmetrized,
$\frac12\bigl(\delta(a,x)\,\delta(b,y)-\delta(a,y)\,\delta(b,x)\bigr)$, for an antisymmetric one). So
$\Sigma(x,y)=\delta\Phi/\delta G(y,x)$ and $\Gamma^0(x_1,x_2;x_3,x_4)=\delta^2\Phi/\delta G(x_4,x_3)\,\delta G(x_2,x_1)$ are
plain derivatives, and $\delta\Phi/\delta V$ opens an interaction line. \pkg{F} may contain the traces and kernels of
operators of section~\ref{sec:effective} (\fnref{TraceLog}, \fnref{OperatorTrace}, \fnref{OperatorKernel},
\fnref{LegendreTransform}), differentiated by $\delta\,{\rm Tr}\ln A={\rm Tr}(A^{-1}\delta A)$ and $\delta G^{-1}=-G^{-1}\,
\delta G\,G^{-1}$, and kernels whose dependence on a source is declared (\fnref{DeclareFunctionalDerivative}). An expression
free of \pkg{G} gives 0 when \pkg{G} is a known kernel (a declared line, a symmetric head); an undeclared head gives the
message \pkg{nosource}. A variation that is not known (a non-integer power of an operator) gives the message
\pkg{opvariation}. The derivative of a functional of the propagator opens its lines: the Hartree--Fock functional gives
the Hartree--Fock self-energy $\Sigma(x,y)=\delta\Phi/\delta G(y,x)$, and its derivative with respect to $V$ opens the
interaction line:
\excode{derivative_propagator}
\exformula{derivative_propagator}
\excode{derivative_two_index}
\exformula{derivative_two_index}

\paragraph{Up to an order} The option \opt{MaximumOrder} \pkg{-> \{x, n\}}, written after the sources, keeps the
terms of order 0 to $n$ in \pkg{x}, a coupling or the head of a line, of the whole derivative (as \fnref{OrderPart} with
\pkg{Range[0, n]} would). The terms of higher order are dropped
after each derivative, since a derivative cannot lower their order; a term whose order cannot be read (\pkg{x} in an
exponential) is kept until the end. A logarithm of a sum polynomial in a coupling is first replaced by its series up to
$x^n$, $\ln(p_0+d)=\ln p_0+\sum_k(-1)^{k+1}(d/p_0)^k/k$, with a common exponential factor taken out,
$\ln(e^X p)=X+\ln p$: so $W=\ln Z$ of $Z=Z_0+Z_1+Z_2/2+\dots$ is differentiated order by order, as in the guided tour.
When \pkg{x} is a source, or a head whose dependence on a source is declared (\fnref{DeclareFunctionalDerivative}), the
derivatives lower its order: the whole derivative is taken, then cut.
\excode{derivative_maximum_order}
\exformula{derivative_maximum_order}

\subsection{Canonical form}\label{sec:canonical}

Two terms are the same if one is obtained from the other by renaming its integration variables, by permuting the
arguments of symmetric functions, and by reordering its Grassmann factors with the corresponding sign. The canonical
form of a term is the representative of its class chosen by a fixed rule (section~\ref{sec:algorithms}); in it the
integration variables are \pkg{DummyVariable[depth, k]}, the odd factors are sorted, and the numerical coefficient is in
front. A term that changes sign under a renaming that maps it to itself --- a symmetry of the diagram with an odd
permutation of Grassmann factors, or an antisymmetric function --- is zero. Terms with the same canonical form are
combined, and their weights added.

\begin{function}{SimplifyDiagrams}
\syntax{SimplifyDiagrams[expr]}
\end{function}
Brings \pkg{expr} to canonical form: the delta functions are integrated, the delta functions between external labels are
used to simplify, the integration variables are renamed canonically and equal terms are combined. Two expressions are
equal if \pkg{SimplifyDiagrams[a - b]} is 0; this is how the identities of the showcase are checked.

\begin{function}{DiagramTerms}
\syntax{DiagramTerms[expr]}
\end{function}
The list of the canonical terms of \pkg{expr}, each with its weight; \pkg{Length[DiagramTerms[expr]]} is the number of
distinct diagrams.
\excode{diagram_terms}
\extext{diagram_terms}

\begin{function}{SetSourcesToZero}
\syntax{SetSourcesToZero[expr]}
\end{function}
Sets all the declared sources to zero: the terms with a Grassmann source, or with an odd graded object
(\fnref{DeclareGrassmannGraded}), vanish, and the commuting sources are replaced by zero.
\excode{sources_to_zero}
\exformula{sources_to_zero}
\pkg{SetSourcesToZero} holds its argument (attribute \pkg{HoldFirst}). A derivative written in it,
\pkg{SetSourcesToZero[FunctionalDerivative[F, s1, \dots, sn]]} (also nested, or with the option \opt{MaximumOrder}), is
taken one source at a time, and after each derivative the terms with more source factors than derivatives left are
dropped: a derivative removes at most one of them, so they vanish at zero sources. The result is the same as setting the
sources to zero in the whole derivative, and the correlation functions of section~\ref{sec:fields}, written as
derivatives at zero source, are computed this way.

% ------------------------------------------------------------------ 5
\section{Generating functionals and Wick's theorem}\label{sec:fields}

\subsection{Fields}

A generating functional $Z[J]$ gives the correlation functions of the fields as derivatives at zero source,
$\langle\phi(a)\phi(b)\cdots\rangle=Z[0]^{-1}\,\delta^n Z/\delta J(a)\,\delta J(b)\cdots\bigr|_{J=0}$. The fields are attached to
their sources by \fnref{DeclareField}.

\begin{function}{FieldDerivative}
\syntax{FieldDerivative[Z, \{phi[a], psibar[b], \dots\}]}
\end{function}
Applies to \pkg{Z} the source derivatives that correspond to the fields, with their signs, the rightmost first; the
sources are not set to zero.
\excode{free_boson}
\exformula{free_boson}

\begin{function}{FieldExpectationValue}
\syntax{FieldExpectationValue[Z, \{phi[a], phi[b], \dots\}]}
\end{function}
The expectation value $\langle\phi(a)\phi(b)\cdots\rangle$ in the generating functional \pkg{Z}: the field derivatives of
\pkg{Z} at zero source divided by \pkg{Z} at zero source. It is computed by literal differentiation, and \pkg{Z} can be any
generating functional, Gaussian or not.
\excode{boson_two_point}
\exformula{boson_two_point}
\excode{free_fermion}
\exformula{free_fermion}
A non-Gaussian functional: the second moment at one point, with a local fourth cumulant (it does not contribute, being
of fourth order in the source):
\excode{non_gaussian}
\exformula{non_gaussian}
The derivatives of $\log Z$ at zero source are the connected functions (the cumulants of the distribution):
\excode{cumulants}
\exformula{cumulants}
For an interacting theory, $W=\log Z$ is best expanded in the coupling before it is differentiated: every derivative of
\pkg{Log[Z]} brings a power of $1/Z$ with all its orders, so that the four-point function of $\phi^4$ theory takes minutes
this way. With $Z=Z_0(1+z)$, where \pkg{z = SimplifyDiagrams[(Z - Z0)/Z0]} is a polynomial in the source (the division
cancels the Gaussian factor), \pkg{W = Log[Z0] + Normal[Series[Log[1 + z], \{lambda, 0, n\}]]} is differentiated in seconds
(the guided tour, section~4); \fnref{ConnectedCorrelationFunction} gives the same functions directly.

\subsection{Gaussian functionals}

\begin{function}{GaussianExpectationValue}
\syntax{GaussianExpectationValue[Z0, \{f1, f2, \dots\}]}
\syntax{GaussianExpectationValue[Z0, \{f1, f2, \dots\}, "Method" -> method]}
\end{function}
The expectation value of a product of fields in a Gaussian generating functional $Z_0=\exp(\text{quadratic form in the
sources})$, by Wick's theorem. The option \opt{Method} can be
\begin{itemize}
\item \pkg{Automatic} (the default): \opt{WickPairings} if \pkg{Z0} is the exponential of a form quadratic in the sources,
  \opt{Differentiation} otherwise;
\item \opt{WickPairings}: the sum over the complete pairings of the fields, every pair being replaced by its contraction
  $\langle f_if_j\rangle_0$ (computed once from \pkg{Z0}, by differentiation), with the sign of the permutation for
  Grassmann fields;
\item \opt{Differentiation}: the literal differentiation of \pkg{Z0}, the terms that cannot survive at zero source being
  dropped along the way; it works for any \pkg{Z0}.
\end{itemize}
The two methods give the same result; the first is much faster for many fields. Asking for \opt{WickPairings} with a
non-Gaussian \pkg{Z0} gives the message \pkg{GaussianExpectationValue::method} and uses differentiation.
\excode{boson_four_point}
\exboth{boson_four_point}
\excode{fermion_four_point}
\exformula{fermion_four_point}

% ------------------------------------------------------------------ 6
\section{Perturbation theory}\label{sec:perturbation}

The interaction is an expression written with the fields, for instance
\begin{code}
Phi4Vertex = lambda/24 IntegralOver[phi[z]^4, {z}];
CoulombVertex = 1/2 IntegralOver[V[x, y] psibar[x] ** psibar[y] ** psi[y] ** psi[x], {x, y}];
\end{code}
for $\phi^4$ theory and for a two-body interaction. The weight of a configuration is $e^{-S_{\rm int}}$. Bosonic fields can appear with powers, Grassmann fields in
\pkg{**} products; the interaction can be a sum of several vertices, with coupling constants as symbols. The order
\pkg{n} of the expansion can be an integer (the orders $0,\dots,n$), \pkg{\{k\}} (order $k$ only) or \pkg{\{kmin, kmax\}}.
All the functions of this section take the option \opt{Method} of \fnref{GaussianExpectationValue}: with
\opt{WickPairings} (the default for Gaussian $Z_0$) interchangeable vertices and identical fields are enumerated once,
with their multiplicity, and the components that cannot be part of the result are discarded while the pairings are
built; with \opt{Differentiation} the full expansion is computed and the unwanted diagrams are dropped at the end.

\begin{function}{PerturbativeCorrelationFunction}
\syntax{PerturbativeCorrelationFunction[Z0, \{f1, f2, \dots\}, interaction, n]}
\end{function}
The correlation function
\[
\langle f_1f_2\cdots\rangle=\frac{\langle f_1f_2\cdots e^{-S_{\rm int}}\rangle_0}{\langle e^{-S_{\rm int}}\rangle_0}
\]
to order $n$ in the interaction: all the diagrams without vacuum bubbles (components without external points). For
fermions with $G_0=-\langle\psi\bar\psi\rangle_0$, the Green function is minus the correlation function of
\pkg{\{psi[x], psibar[y]\}}.
\excode{phi4_propagator}
\exboth{phi4_propagator}
The electron gas to first order (the first session of section~\ref{sec:overview}): the Hartree and Fock diagrams.
\excode{electron_gas_first_order}
\exboth{electron_gas_first_order}

\begin{function}{ConnectedCorrelationFunction}
\syntax{ConnectedCorrelationFunction[Z0, \{f1, f2, \dots\}, interaction, n]}
\end{function}
The connected correlation function (cumulant) to order $n$: the diagrams in which all the external points are connected,
the derivatives of $W=\log Z$ with respect to the sources at zero source (example \pkg{cumulants} above), computed directly
order by order. Fields at the same point form a composite operator. The density correlation function is the response of
the density to a potential $U$ coupled to it, $\langle\rho_a\rho_b\rangle_c=-\delta\langle\rho(a)\rangle/\delta
U(b)=-\delta G(a,a)/\delta U(b)$: with $G_0^{-1}\to G_0^{-1}-U$, the derivative of the bare line is declared once
(\fnref{DeclareFunctionalDerivative}), and \fnref{FunctionalDerivative} inserts the potential into every line of $G$; at
first order:
\excode{density_correlation}
\expictures{density_correlation}
The shortcut is the connected four-point function with the fields taken pairwise at the same point,
\pkg{ConnectedCorrelationFunction[Z0, \{psibar[a], psi[a], psibar[b], psi[b]\}, interaction, n]}:
\excode{density_correlation_derivative}
\extext{density_correlation_derivative}

\paragraph{Only one order} These two functions and \fnref{ConnectedVacuumDiagrams} (below) have the option
\opt{CouplingOrder} \pkg{-> \{g, k\}}, with which they give only the terms of order $g^k$, the same as
\pkg{OrderPart[\dots, g, k]}, without building the combinations of vertices of other orders: for theories with several
vertices (QCD) this is much faster.

\begin{function}{ConnectedVacuumDiagrams}
\syntax{ConnectedVacuumDiagrams[Z0, interaction, n]}
\end{function}
$\log(Z/Z_0)$ to order $n$: the connected vacuum diagrams with their weights. A disconnected vacuum diagram is the product
of its connected pieces, with an extra $1/k!$ for $k$ identical pieces, so $Z/Z_0$ is the exponential of the sum of the
connected ones (linked-cluster theorem); for fermions the sign of each closed loop belongs to its piece. The result is the
same as \pkg{ConnectedDiagrams[}$z$\pkg{, \{\}]} of the series $z=Z/Z_0-1$ computed by derivatives. The grand
potential is $\Omega-\Omega_0=-T\log(Z/Z_0)$.
\excode{phi4_vacuum}
\exboth{phi4_vacuum}

\begin{function}{OrderPart}
\syntax{OrderPart[expr, x, k]}
\end{function}
The terms of \pkg{expr} of order $k$: if \pkg{x} is a symbol, the terms proportional to $x^k$; if \pkg{x} is the head of
the interaction lines, the terms with $k$ factors \pkg{x[\dots]}.
\excode{phi4_order_part}
\exformula{phi4_order_part}
\excode{order_by_head}
\exformula{order_by_head}

% ------------------------------------------------------------------ 7
\section{Classes of diagrams}\label{sec:classes}

\subsection{The graph of a diagram}\label{sec:graph}

The functions of this section decide the class of a diagram on a graph built from its term. The \emph{points} are the
external labels (the second argument of the functions) and the integration variables. Every factor that is not one of
the \emph{lines} --- the factors with the heads given as last argument, typically the propagators --- glues its points
into one \emph{node}: an interaction line, a vertex function, a delta function. The lines are the edges of the graph.
An external label that appears in a single factor, a line joining it to an integration variable, is an \emph{external
leg}: the leg is not an edge, and the node at its other end is marked as external. Any other external label (one that
sits on a vertex, as in amputated functions) is a point of the graph, and its node is external. Lines from a node to
itself (a tadpole $G(z,z)$) are never cut.

\subsection{Connectivity}

\begin{function}{ConnectedDiagrams}
\syntax{ConnectedDiagrams[expr, \{x1, \dots, xn\}]}
\end{function}
The terms whose diagram is connected; the external labels \pkg{x1, \dots} and the integration variables are its points,
and every factor joins the points it contains.

\begin{function}{DropVacuumBubbles}
\syntax{DropVacuumBubbles[expr, \{x1, \dots, xn\}]}
\end{function}
Drops the terms with a vacuum bubble, a connected component without external labels.

\subsection{Irreducibility and skeletons}

\begin{function}{ParticleIrreducibleQ}
\syntax{ParticleIrreducibleQ[term, n, \{x1, \dots\}, \{G, \dots\}]}
\end{function}
\pkg{True} if the diagram of the single term \pkg{term} is connected and stays connected after cutting any $n$ or fewer
lines with the heads \pkg{G, \dots} (section~\ref{sec:graph}).

\begin{function}{ParticleIrreducibleDiagrams}
\syntax{ParticleIrreducibleDiagrams[expr, n, \{x1, \dots\}, \{G, \dots\}]}
\end{function}
The $n$-particle-irreducible terms of \pkg{expr}. With $n=1$ and the propagators as lines, the one-particle-irreducible
(1PI) diagrams; \pkg{ParticleIrreducibleDiagrams[vacuum, 2, \{\}, \{G\}]} gives the two-particle-irreducible vacuum
diagrams. Since the external legs are not lines, the 1PI part of a Green function is the sum of its diagrams that do not
fall apart when one internal line is cut.
\excode{phi4_second_order}
\expictures{phi4_second_order}
\excode{one_particle_irreducible}
\expictures{one_particle_irreducible}

\begin{function}{SkeletonDiagrams}
\syntax{SkeletonDiagrams[expr, \{x1, \dots\}, \{G, \dots\}]}
\end{function}
The skeleton diagrams: one-particle irreducible, and such that no cut of two lines \pkg{G} separates a part without
external points (a self-energy insertion). For a two-point function with external legs these are the diagrams that close
into two-particle-irreducible vacuum diagrams; for amputated vertex functions, the usual skeleton diagrams.
\excode{electron_gas_second_order}
\expictures{electron_gas_second_order}
\excode{skeleton}
\expictures{skeleton}

\subsection{Amputation}

\begin{function}{AmputateExternalLegs}
\syntax{AmputateExternalLegs[expr, \{x1, \dots\}, \{G, \dots\}]}
\end{function}
Removes the external legs: in every term, the line \pkg{G} that joins the external label $x_i$ to an integration variable
$z$ is deleted and $z$ is renamed $x_i$; when two external labels end on the same vertex, the second is joined to the
first by a \pkg{DeltaFunction}. Terms in which an external label is not joined to an integration variable by a single
line are dropped. The self-energy is the amputated 1PI part of the Green function:
\excode{amputated}
\exboth{amputated}

\subsection{Cuts between groups of external points}

\begin{function}{ReducibleDiagrams}
\syntax{ReducibleDiagrams[expr, \{\{x1, \dots\}, \{y1, \dots\}\}, n, \{G, \dots\}]}
\end{function}
The terms in which cutting $n$ or fewer lines \pkg{G} separates the external labels $x_1,\dots$ from $y_1,\dots$ (each group
staying connected). For instance the diagrams of a response function that fall apart when a single interaction line is
cut --- the improper (interaction-reducible) diagrams, of which the chain of bubbles of the random-phase approximation is
the simplest:
\excode{density_reducible}
\expictures{density_reducible}

\begin{function}{IrreducibleDiagrams}
\syntax{IrreducibleDiagrams[expr, \{\{x1, \dots\}, \{y1, \dots\}\}, n, \{G, \dots\}]}
\end{function}
The other terms: no cut of $n$ or fewer lines separates the two groups. The proper part of the density
correlation, minus the polarization $P$ of Hedin's equations (section~\ref{sec:exact}), is
\pkg{IrreducibleDiagrams[chi, \{\{a\}, \{b\}\}, 1, \{V\}]}.

\begin{function}{TwoParticleReducibleQ}
\syntax{TwoParticleReducibleQ[term, \{\{x1, x2\}, \{x3, x4\}\}, \{G, \dots\}]}
\end{function}
\pkg{True} if cutting two or fewer lines \pkg{G} separates $x_1,x_2$ from $x_3,x_4$ in the diagram of the single term \pkg{term}.

\subsection{Closed loops}

\begin{function}{ClosedLoopCount}
\syntax{ClosedLoopCount[term, \{G, \dots\}]}
\end{function}
The number of closed loops made of the directed lines \pkg{G} --- fermion loops --- in the diagram of a single term:
the connected components of these lines in which every point has one line in and one out.
\excode{closed_loops}
\extext{closed_loops}

\begin{function}{DropClosedLoops}
\syntax{DropClosedLoops[expr, \{G, \dots\}]}
\end{function}
Drops the terms with closed loops of the lines \pkg{G}, for instance the fermion loops that vanish in a quenched average
over disorder (replica limit).

% ------------------------------------------------------------------ 8
\section{Conventions for many-body problems}\label{sec:conventions}

\begin{itemize}
\item \textbf{Fermions.} $G_0(x,y)=-\langle T\psi(x)\bar\psi(y)\rangle_0$, with the free generating functional and the
  fields
\begin{code}
DeclareGrassmannSource[eta, etabar]; DeclareField[psi, etabar]; DeclareField[psibar, eta, -1];
Z0 = Exp[-IntegralOver[G0[x, y] etabar[x] ** eta[y], {x, y}]];
\end{code}
  The Green function, $G=G_0+G_0\Sigma G$, is \pkg{-PerturbativeCorrelationFunction[Z0, \{psi[x], psibar[y]\}, \dots]}, and
  \pkg{G[x, y]} is drawn as an arrow from $y$ to $x$.
\item \textbf{Two-body interaction.}
\begin{code}
DeclareSymmetric[V];
interaction = 1/2 IntegralOver[V[x, y] psibar[x] ** psibar[y] ** psi[y] ** psi[x], {x, y}];
\end{code}
  The labels contain imaginary time (an instantaneous interaction carries
  $\delta(\tau_x-\tau_y)$) and spin: for a spin-independent interaction each closed fermion loop carries the factor $2s+1$
  on evaluation, or spin is summed explicitly (\fnref{DeclareDiscreteIndexValues}). An interaction line stands for $-V$ and a vertex for minus its coupling: the minus
  sign comes from expanding $e^{-S_{\rm int}}$ (section~\ref{sec:weights}). The Hartree self-energy is
  $+\delta(x,y)\int_zV(x,z)G(z,z)$, drawn with the weight $-1$, the Fock self-energy $-G(x,y)V(x,y)$, drawn with the weight
  $+1$: a diagram of $G$ or $\Sigma$ with $n_L$ closed fermion loops has the weight $(-1)^{n_L}$, as in the textbooks,
  and its coefficient is $(-1)^{n+n_L}$ at order $n$.
\item \textbf{Closed fermion loops} carry their factor $-1$ automatically, from the Grassmann algebra; so do the exchange
  terms.
\item \textbf{Vertices.} $\Sigma(x,y)=\delta\Phi/\delta G(y,x)$, $\Gamma^0(x_1,x_2;x_3,x_4)=\delta\Sigma(x_1,x_2)/\delta G(x_4,x_3)$; the
  fermions leave at $x_1$, $x_3$ and enter at $x_2$, $x_4$;
  $\langle\psi_{1'}\bar\psi_{2'}\bar\psi_{4'}\psi_{3'}\rangle_c=G_{1'1}G_{22'}F_{12;34}G_{44'}G_{3'3}$. A pair of
  propagators joining two vertices in a channel is written $L_0$ (never $GG$): in the particle--hole channel
  $L_0(x_1,x_2;x_3,x_4)=G(x_1,x_4)\,G(x_3,x_2)$, so that $F=\Gamma^0+\Gamma^0L_0F$.
\item \textbf{Real bosons.} $G=\langle\phi\phi\rangle$, $Z_0=\exp\bigl(\frac12\int_{xy}J(x)D_0(x,y)J(y)\bigr)$, $G=G_0+G_0\Sigma G$,
  $\Sigma=-2\,\delta\Phi/\delta G$ ($\Phi$ the 2PI part $\Gamma_2$ of the effective action), $F=\Gamma^{(4)}$.
\item \textbf{Superfields.} The Nambu field $\boldsymbol\psi_a=(\psi,\bar\psi)$, a label $a=(x,s)$ with a component: one
  Grassmann source, $Z_0=\exp\bigl(\frac12\int_{ab}J_a\mathcal G_{0,ab}J_b\bigr)$, the antisymmetric propagator $\mathcal G_{ab}=-\langle
  T\boldsymbol\psi_a\boldsymbol\psi_b\rangle$ with the blocks $\mathcal G_{\psi\psi}$, $G=\mathcal G_{\psi\bar\psi}$, $-G^T$, $\mathcal G_{\bar\psi\bar\psi}$, and
  $S_{\rm int}=\frac1{4!}\int\mathcal V_{abcd}\boldsymbol\psi_a\boldsymbol\psi_b\boldsymbol\psi_c\boldsymbol\psi_d$; $\boldsymbol\Sigma=-2\,\delta\Phi/\delta\mathcal G$,
  with the gap functions $\Delta(x,y)=\boldsymbol\Sigma_{x\bar\psi,y\bar\psi}$ and $\bar\Delta(x,y)=\boldsymbol\Sigma_{x\psi,y\psi}$, and
  $\Gamma^0=4\,\delta^2\Phi/\delta\mathcal G\,\delta\mathcal G$, the kernel of the channel $s$. With the legs $(x_1\bar\psi,x_2\psi,x_3\bar\psi,x_4\psi)$
  the channels $s$, $t$, $u$ are $\ph$, $\pp$, $\phbar$.
\item \textbf{Electrons and phonons, disorder.} Further Gaussian fields multiply the free generating functional: a
  phonon field $\varphi$ with propagator $D_0$ and the vertex $g\int_z\varphi(z)\bar\psi(z)\psi(z)$ (\pkg{DeclareCoupling[g]}: in the
  pictures each vertex stands for $-g$); a Gaussian random
  potential $u$ with $\overline{u(x)u(y)}=\gamma(x,y)$ in the same way, the quenched average dropping the closed fermion
  loops (\fnref{DropClosedLoops}). The showcase has the code of both.
\item \textbf{Momentum and frequency.} $G_0(k,i\omega_n)=1/(i\omega_n-\xi_k)$; the density correlation function of the
  package is minus the Lindhard function, and its boson propagator $D_0=\langle\phi\phi\rangle_0$ minus the
  $\mathcal D_0=-\langle T\phi\phi\rangle_0$ of the textbooks (section~\ref{sec:matsubara-conventions}).
\item \textbf{Names.} The same objects have different names in condensed matter, field theory, nuclear physics, quantum
  chemistry and statistical physics: the self-energy is also the mass operator, $\Phi$ the 2PI part $\Gamma_2$ of the
  effective action, $\Gamma^0$ the Bethe--Salpeter kernel, the bubble chain the RPA, and so on. The table \emph{Other names} at the
  beginning of the showcase collects them.
\end{itemize}

% ------------------------------------------------------------------ 9
\section{Self-energy, Luttinger--Ward functional and vertices}\label{sec:functionals}

\subsection{Conventions}

\begin{itemize}
\item Fermions: $G=-\langle T\psi\bar\psi\rangle$, $G=G_0+G_0\Sigma G$. The Luttinger--Ward functional
  $\Phi[G]$~\cite{LuttingerWard1960,Baym1962} is minus the sum of the two-particle-irreducible connected vacuum diagrams
  with $G_0\to G$; $\Sigma(x,y)=\delta\Phi/\delta G(y,x)$ and
  $\Gamma^0(x_1,x_2;x_3,x_4)=\delta\Sigma(x_1,x_2)/\delta G(x_4,x_3)$. The fermions leave the vertex at $x_1$, $x_3$ and enter
  at $x_2$, $x_4$.
\item Real bosons (option \pkg{"Statistics" -> "Boson"}): $G=\langle\phi\phi\rangle$, $G=G_0+G_0\Sigma G$, $\Phi$ is the
  two-particle-irreducible part $\Gamma_2$ of the effective action, $\Sigma(x,y)=-2\,\delta\Phi/\delta G(x,y)$ and
  $\Gamma^0(x_1,x_2,x_3,x_4)=-2\,\delta\Sigma(x_1,x_2)/\delta G(x_3,x_4)$, the derivatives with respect to the symmetric $G$
  being symmetrized.
\item Complex bosons (option \pkg{"Statistics" -> "ComplexBoson"}): the propagator $G=\langle\phi\phi^*\rangle$ is
  directed, $G=G_0+G_0\Sigma G$, $\Sigma(x,y)=-\delta\Phi/\delta G(y,x)$ and
  $\Gamma^0(x_1,x_2;x_3,x_4)=\delta\Sigma(x_1,x_2)/\delta G(x_4,x_3)$. The channels and the parquet functions are for
  fermions and real bosons only.
\item The derivatives with respect to $G$ are functional derivatives in which \pkg{G} is treated as a source:
  differentiating a diagram opens one of its lines, in all possible ways.
\end{itemize}
Every object of this section is a derivative of $\Phi$, and \fnref{FunctionalDerivative} computes it as one, with no
declaration of \pkg{G} as a source; \fnref{SelfEnergy} and \fnref{IrreducibleVertex} are shortcuts for these derivatives,
with the factors of each statistics. The Luttinger--Ward functional of the electron gas to second order
(\fnref{LuttingerWardFunctional}):
\excode{luttinger_ward}
\exboth{luttinger_ward}
its self-energy, $\Sigma(x,y)=\delta\Phi/\delta G(y,x)$:
\excode{self_energy}
\exboth{self_energy}
and the irreducible vertex to first order, $\Gamma^0(x_1,x_2;x_3,x_4)=\delta^2\Phi/\delta G(x_4,x_3)\,\delta G(x_2,x_1)$ (the
rightmost derivative is taken first), the direct and the exchange interaction:
\excode{irreducible_vertex}
\exboth{irreducible_vertex}
For real bosons $\Sigma(x,y)=-2\,\delta\Phi/\delta G(x,y)$ and $\Gamma_s(a,b,c,d)=4\,\delta^2\Phi/\delta G(c,d)\,\delta G(a,b)$, in code
\pkg{-2 FunctionalDerivative[Phi, G[x, y]]} and \pkg{4 FunctionalDerivative[Phi, G[c, d], G[a, b]]}.

\subsection{Functions}

\begin{function}{LuttingerWardFunctional}
\syntax{LuttingerWardFunctional[Z0, interaction, n, G0 -> G]}
\end{function}
The Luttinger--Ward functional to order $n$: minus the sum of the two-particle-irreducible connected vacuum diagrams,
with the bare propagator \pkg{G0} replaced by the dressed one \pkg{G}. Takes the option \opt{Method}. The functional of
the electron gas to second order is the first example of this section.

\begin{function}{SelfEnergy}
\syntax{SelfEnergy[Phi, G, \{x, y\}]}
\syntax{SelfEnergy[Phi, G, \{x, y\}, "Statistics" -> "Boson"]}
\end{function}
The self-energy $\Sigma(x,y)$ from the functional \pkg{Phi}. For fermions (the default) it is $\delta\Phi/\delta G(y,x)$;
for real bosons $-2\,\delta\Phi/\delta G(x,y)$, symmetrized; for complex bosons, with the option
\pkg{"Statistics" -> "ComplexBoson"}, $-\delta\Phi/\delta G(y,x)$; for a superfield (\pkg{"Statistics" -> "Superfield"}, an
antisymmetric propagator, see \fnref{SuperfieldComponents}) $-2\,\delta\Phi/\delta\mathcal G(x,y)$, antisymmetrized.
\pkg{"Statistics" -> "Boson"} with a propagator that is not
declared symmetric gives a message (section~\ref{sec:troubleshooting}), and so does the default on a propagator declared
symmetric or antisymmetric (\pkg{statisticsdefault}: the fermionic formula on the $\Gamma_2$ of a real boson gives
$-\frac12$ of its self-energy). It is the derivative of the examples above, $\delta\Phi/\delta G(y,x)$ for fermions:
\excode{self_energy_shortcut}
\extext{self_energy_shortcut}

\begin{function}{IrreducibleVertex}
\syntax{IrreducibleVertex[Phi, G, \{x1, x2, x3, x4\}]}
\syntax{IrreducibleVertex[Phi, G, \{x1, x2, x3, x4\}, "Statistics" -> "Boson"]}
\syntax{IrreducibleVertex[Phi, G, \{a, b, c, d\}, "Statistics" -> "Superfield", "Channel" -> All]}
\end{function}
The vertex irreducible in the particle--hole channel $(x_1x_2|x_3x_4)$, the kernel of the Bethe--Salpeter equation:
$\Gamma^0(x_1,x_2;x_3,x_4)=\delta\Sigma(x_1,x_2)/\delta G(x_4,x_3)=\delta^2\Phi/\delta G(x_4,x_3)\,\delta G(x_2,x_1)$ for fermions,
$-2\,\delta\Sigma(x_1,x_2)/\delta G(x_3,x_4)$ for real bosons and superfields (antisymmetrized). With \opt{Channel} the kernels
of the other channels come from the same functional: \pkg{"s"}, \pkg{"t"}, \pkg{"u"} or \pkg{All} for real bosons and
superfields, $\Gamma_t(a,b,c,d)=\pm\Gamma_s(a,c,b,d)$, $\Gamma_u(a,b,c,d)=\Gamma_s(a,d,b,c)$ (the sign $-$ for superfields);
\pkg{"TransverseParticleHole"} for fermions, $-\Gamma^0(x_3,x_2;x_1,x_4)$; the particle--particle kernel of fermions is a
derivative with respect to the anomalous propagators, the channel \pkg{"t"} of the superfield. For fermions $\Gamma^0$ is symmetric under
$(x_1x_2)\leftrightarrow(x_3x_4)$; the bare part of the first order is the direct and exchange interaction. As a derivative
of the components $\Phi[G,\mathcal G_{\psi\psi},\mathcal G_{\bar\psi\bar\psi}]$ of the superfield (\fnref{SuperfieldComponents}), the
particle--particle kernel is
\[
\Gamma_{pp}(x_1,x_2;x_3,x_4)=-4\,\frac{\delta^2\Phi}{\delta\mathcal G_{\psi\psi}(x_2,x_4)\,\delta\mathcal G_{\bar\psi\bar\psi}(x_1,x_3)}
\Big|_{\mathcal G_{\psi\psi}=\mathcal G_{\bar\psi\bar\psi}=0},
\]
next to $\Gamma^0=\delta^2\Phi/\delta G(x_4,x_3)\,\delta G(x_2,x_1)$, the kernel of the particle--hole channel. A list of labels
that is not of length four gives the message \pkg{vertexlabels} and \pkg{\$Failed}. The shortcut and the derivative of
\fnref{FunctionalDerivative} above:
\excode{irreducible_vertex_shortcut}
\extext{irreducible_vertex_shortcut}

\begin{function}{DysonSeries}
\syntax{DysonSeries[sigma, \{s1, s2\}, G0, \{x, y\}, m]}
\syntax{DysonSeries[sigma, \{s1, s2\}, G0, \{x, y\}, m, "MaximumOrder" -> \{V, n\}]}
\end{function}
$G_0(x,y)+G_0\Sigma G_0+G_0\Sigma G_0\Sigma G_0+\dots$ with up to $m$ self-energy insertions, where \pkg{sigma} is an
expression for $\Sigma(s_1,s_2)$. With \opt{MaximumOrder} \pkg{-> \{V, n\}}, only the chains of total order $n$ or less in
\pkg{V} (a coupling symbol or the head of the interaction lines) are built. The chains of the skeleton self-energy
$\Sigma[G_0]$ with up to two insertions, kept to second order in $V$:
\excode{dyson_series}
\expictures{dyson_series}

\begin{function}{PerturbativeFullVertex}
\syntax{PerturbativeFullVertex[Z0, interaction, n, \{psi, psibar\}, \{x1, x2, x3, x4\}, G0 -> G]}
\syntax{PerturbativeFullVertex[Z0, interaction, n, \{phi\}, \{x1, x2, x3, x4\}, G0 -> G]}
\end{function}
The full four-point vertex $F$ to order $n$, as a skeleton expansion in the dressed propagator \pkg{G}; with
\pkg{\{k\}} or \pkg{\{kmin, kmax\}} only those orders are computed. For fermions
(fields \pkg{\{psi, psibar\}}),
\[
\langle\psi(1')\bar\psi(2')\bar\psi(4')\psi(3')\rangle_c=G(1',1)\,G(2,2')\,F(1,2;3,4)\,G(4,4')\,G(3',3):
\]
the connected four-point function is computed to order $n$, its legs are removed and its skeleton diagrams are kept,
with $G_0\to G$;
$F$ is one-particle irreducible and $F=\Gamma^0+O(V^2)$. In code, $F$ is the connected function
\begin{code}
ConnectedCorrelationFunction[Z0, {psi[x1], psibar[x2], psibar[x4], psi[x3]}, interaction, n]
\end{code}
amputated with \fnref{AmputateExternalLegs}, its skeleton diagrams kept with \fnref{SkeletonDiagrams}, and \pkg{G0 -> G}; with bare lines it
is the 1PI part, \fnref{ParticleIrreducibleDiagrams} with $n=1$, of the amputated function (\fnref{AmputateExternalLegs} alone
keeps the self-energy insertions on the legs). With spin the fields carry the index of their source, and $F$ with given
spins is written in this way, \pkg{psi[x1, up]}, \dots; the bare propagators made full at once are a list of rules,
\pkg{\{G0up -> Gup, G0down -> Gdown\}}. For real bosons (field \pkg{\{phi\}}), $F$ is the 1PI vertex
$\Gamma^{(4)}=-(\text{amputated connected four-point function})$, so that $F=\lambda+O(\lambda^2)$ in $\phi^4$ theory; the
same for a superfield (one Grassmann field, \pkg{\{NambuField\}}), whose $F$ is $\mathcal V+O(\mathcal V^2)$.
Takes the option \opt{Method}. The vertex is the fourth derivative of the 1PI action (\fnref{EffectiveAction1PI},
section~\ref{sec:effective}) at zero mean field, and the result is its skeleton diagrams with $G_0\to G$ (for fermions, the
derivative $\delta^4\Gamma/\delta\bar\Psi\,\delta\Psi\,\delta\bar\Psi\,\delta\Psi$ of $\Gamma[\bar\Psi,\Psi]$); for $\phi^4$ theory,
with the dressed line $D$ (\pkg{DressedPropagator}):
\excode{full_vertex_phi4}
\exboth{full_vertex_phi4}
This function computes the same from the connected function, the other side of the Legendre transform:
\excode{full_vertex_derivative}
\extext{full_vertex_derivative}

% ------------------------------------------------------------------ 10
\section{Effective actions and operator traces}\label{sec:effective}

The effective actions are objects of the package, written with inert traces of operators and differentiated with respect
to the propagator (section~\ref{sec:derivatives}), to the mean fields and to one-body sources. An operator is a kernel ---
a symbol \pkg{h} standing for $h(x,y)$, or a \fnref{LabelFunction} --- a number (times the identity), a sum, a coefficient
times an operator, a product \pkg{A . B} (\pkg{Dot}), an integer power or an inverse \pkg{A\^{}-1}. The kernel of
\pkg{A . B} is $\int A(x,u)\,B(u,y)$, and ${\rm Tr}(A\,B)=\int A(x,y)\,B(y,x)$: with this trace ${\rm Tr}(\Sigma G)$ is the
$\int\Sigma(x,y)\,G(y,x)$ of \fnref{SelfEnergy}. A single symbol is a kernel; in a product with \pkg{Times} the factor
that contains a declared kernel (\fnref{DeclareDirectedPropagator}, \fnref{DeclareSymmetric}, a source) is the operator
and the others are coefficients, and when this cannot be told (\pkg{lambda S}, neither declared) the message
\pkg{opscalar} is given and the result is \pkg{\$Failed}. A chemical potential is therefore
\pkg{LabelFunction[mu DeltaFunction[a, b], \{a, b\}]}. \pkg{Inverse[A]} is read as \pkg{A\^{}-1}. Write a product with an
inverse as \pkg{G0\^{}(-1) . G}: Mathematica reads \pkg{G0\^{}-1 . G} as \pkg{G0\^{}(-(1 . G))}, a power whose exponent is a
product, which the package takes as \pkg{G0\^{}(-1) . G} with the message \pkg{opdotpower}. A source enters an operator
through a kernel, a \fnref{LabelFunction} or a head whose derivative is declared (\fnref{DeclareFunctionalDerivative}), and
the derivatives of traces and kernels follow it there; a source outside any kernel (\pkg{D0\^{}-1 + lambda/2 phibar\^{}2},
without labels) gives the message \pkg{opsource}, and the derivative is taken as 0.

\subsection{Operators and traces}

\begin{function}{TraceLog}
\syntax{TraceLog[A]}
\end{function}
${\rm Tr}\ln A$, inert. \fnref{FunctionalDerivative} differentiates it, $\delta\,{\rm Tr}\ln A={\rm Tr}(A^{-1}\delta A)$;
\fnref{ExpandTrace} expands it around a bare propagator, \fnref{EvaluateMatsubara} evaluates it in momentum and frequency,
and \fnref{ZeroDimensionalValue} in zero dimensions.
\excode{derivative_trace_log}
\exformula{derivative_trace_log}
The one-loop action of $\phi^4$ theory written by hand, $\frac12{\rm Tr}\ln S''[\bar\phi]$ with $S''=D_0^{-1}+\frac\lambda2\bar\phi^2$:
its fourth derivative at $\bar\phi=0$ is the bubble in the three channels, the one-loop part of $\Gamma^{(4)}$:
\excode{trace_log_label_function}
\exformula{trace_log_label_function}

\begin{function}{OperatorTrace}
\syntax{OperatorTrace[A]}
\end{function}
The trace ${\rm Tr}\,A$, inert; \pkg{OperatorTrace[1]} is the trace of the identity. \fnref{ExpandTrace} writes the trace
of a product of kernels as a ring of diagrams.

\begin{function}{OperatorKernel}
\syntax{OperatorKernel[A][x, y]}
\end{function}
The kernel $A(x,y)$ of an operator that is not a single kernel, e.g.\ $G^{-1}(x,y)$ (symmetric when $G$ is).
\fnref{ExpandTrace} writes it as diagrams and contracts $\int G^{-1}(a,u)\,G(u,b)$ to $\delta(a,b)$:
\excode{operator_kernel}
\exformula{operator_kernel}

\begin{function}{ExpandTrace}
\syntax{ExpandTrace[expr]}
\syntax{ExpandTrace[expr, G0, n]}
\syntax{ExpandTrace[expr, G0, n, "MaximumOrder" -> \{V, m\}]}
\end{function}
Writes the traces, the kernels of operators and Potthoff's functional \fnref{LegendreTransform} (as
$\Phi[G]+\kappa\,{\rm Tr}(S\,G)$) in \pkg{expr} as diagrams, contracts the products of inverse kernels and brings the result
to canonical form. With \pkg{G0} and \pkg{n} it also expands around the bare propagator, to \pkg{n} insertions:
${\rm Tr}\ln\bigl(c(G_0^{-1}-X)\bigr)={\rm Tr}\ln(cG_0^{-1})-\sum_k{\rm Tr}\bigl((G_0X)^k\bigr)/k$, a series of rings, and
$\bigl(c(G_0^{-1}-X)\bigr)^{-1}=c^{-1}\sum_k(G_0X)^kG_0$, the Dyson series; a propagator replaced by an operator, e.g.\
$G\to(G_0^{-1}-S)^{-1}$ with $S$ a \fnref{LabelFunction} of a self-energy, is expanded inside the diagrams too. With
\opt{MaximumOrder} \pkg{-> \{V, m\}} only the terms of order $m$ or less in \pkg{V} (a coupling or the head of the
interaction lines) are kept, the products being truncated as they are formed. The rings of the self-energy $\Sigma$ (\pkg{Sg}):
\excode{trace_log_rings}
\exboth{trace_log_rings}

\subsection{The 2PI and 1PI effective actions}

\begin{function}{EffectiveAction2PI}
\syntax{EffectiveAction2PI[Phi, G0 -> G]}
\syntax{EffectiveAction2PI[Phi, G0 -> G, "Statistics" -> "Boson"]}
\syntax{EffectiveAction2PI[Phi, G0 -> G, "Statistics" -> "Superfield"]}
\syntax{EffectiveAction2PI[Z0, interaction, n, G0 -> G, options]}
\end{function}
The two-particle-irreducible effective action $\Gamma[G]$ built on the Luttinger--Ward functional \pkg{Phi} (minus the 2PI
vacuum diagrams, \fnref{LuttingerWardFunctional}):
\[
\Gamma[G]=\Phi[G]-{\rm Tr}\ln(-G^{-1})-{\rm Tr}(G_0^{-1}G-1)
\]
for fermions ($G=-\langle\psi\bar\psi\rangle$), $\Gamma_2[G]+\frac12{\rm Tr}\ln G^{-1}+\frac12{\rm Tr}(D_0^{-1}G-1)$ for real
bosons, $\Phi+{\rm Tr}\ln G^{-1}+{\rm Tr}(G_0^{-1}G-1)$ for complex bosons (\opt{Statistics} \pkg{-> "ComplexBoson"}) and
\[
\Gamma[\mathcal G]=\Phi[\mathcal G]-\tfrac12{\rm Tr}\ln(-\mathcal G^{-1})-\tfrac12{\rm Tr}(\mathcal G_0^{-1}\mathcal G-1)
\]
for a superfield (\opt{Statistics} \pkg{-> "Superfield"}), the functional of Nambu and Gor'kov, whose variable is the Nambu
propagator $\mathcal G=-\langle T\boldsymbol\psi\boldsymbol\psi\rangle$ of $\boldsymbol\psi=(\psi,\bar\psi)$ with its anomalous blocks $-\langle
T\psi\psi\rangle$, $-\langle T\bar\psi\bar\psi\rangle$: the trace terms are half those of fermions (the Gaussian integral over one
Grassmann field is a Pfaffian), at zero anomalous blocks it is the fermionic $\Gamma[G]$ term by term
(\fnref{SuperfieldComponents}), and its stationarity at first order gives the Hartree--Fock--Bogoliubov equations, with the gap
equation $\Delta=V\mathcal G_{\psi\psi}$ of BCS theory. It
is the Legendre transform of $-\log Z$ ($W=\log Z$) with respect to the two-point source $K$ (and, with \opt{MeanFields},
the one-point source $J$), stationary where $G^{-1}=G_0^{-1}-\Sigma[G]$ (the
Dyson equation) and equal to $-\log Z$, the grand potential over $T$, there. It is also called the Luttinger--Ward or
Baym--Kadanoff functional $\Omega[G]$~\cite{LuttingerWard1960,Baym1962}, and it is the $\Gamma[\bar\phi,G]$ of Cornwall, Jackiw
and Tomboulis at $\bar\phi=0$~\cite{CornwallJackiwTomboulis1974}; Almbladh, von Barth and van
Leeuwen~\cite{AlmbladhVonBarthVanLeeuwen1999} call this form the Klein functional and keep the name Luttinger--Ward for the
form with $\Sigma[G]$ inside the logarithm (the two agree at the stationary point). A truncated $\Phi$ gives a
$\Phi$-derivable, conserving approximation. With four arguments $\Phi$ is computed by \fnref{LuttingerWardFunctional}.
\excode{effective_action_2pi}
\exformula{effective_action_2pi}
\excode{effective_action_dyson}
\exformula{effective_action_dyson}
The objects of sections~\ref{sec:functionals} and \ref{sec:parquet} are derivatives of these functionals, at zero means and
fixed propagator, in the order of \fnref{FunctionalDerivative} (the rightmost derivative first, so that two Grassmann
derivatives in the other order change the sign; four do not):
\begin{center}\footnotesize
\begin{tabular}{@{}r@{\enspace}l@{\quad}l@{}}
\toprule
factor & arguments of \pkg{FunctionalDerivative} & object\\
\midrule
\multicolumn{3}{@{}l}{\emph{fermions}, $\Gamma[\bar\Psi,\Psi,G]$ (\pkg{"MeanFields" -> \{psi -> Psi, psibar -> Psibar\}})}\\
 & \pkg{Phi, G[y, x]} & $\Sigma(x,y)$\\
 & \pkg{Phi, G[x4, x3], G[x2, x1]} & $\Gamma^0(x_1,x_2;x_3,x_4)$, particle--hole\\
 & \pkg{Gamma, G[y, x]} & $(\Sigma+G^{-1}-G_0^{-1})(x,y)$, 0 solves Dyson\\
 & \pkg{Gamma, Psibar[x], Psi[y]} & $(G_0^{-1}-\Sigma[G])(x,y)$\\
 & \pkg{Gamma, G[x4, x3], Psibar[x1], Psi[x2]} & $-\Gamma^0(x_1,x_2;x_3,x_4)$\\
 & \pkg{Gamma, Psibar[x1], Psi[x2], Psibar[x3], Psi[x4]} & $\Gamma_{\rm irr}(x_1,x_2;x_3,x_4)$, fully irreducible\\
\multicolumn{3}{@{}l}{\emph{real bosons}, $\Gamma[\bar\phi,G]$ (\pkg{"Boson"}, \pkg{"MeanFields" -> \{phi -> phiMean\}})}\\
$-2$ & \pkg{Phi, G[x, y]} & $\Sigma(x,y)$\\
$4$ & \pkg{Phi, G[c, d], G[a, b]} & $\Gamma_s(a,b,c,d)$; $\Gamma_t$, $\Gamma_u$ by relabelling\\
 & \pkg{Gamma, phiMean[a], phiMean[b]} & $(D_0^{-1}-\Sigma[G])(a,b)$\\
 & \pkg{Gamma, phiMean[a], \dots, phiMean[d]} & $\Gamma_{\rm irr}(a,b,c,d)$\\
\multicolumn{3}{@{}l}{\emph{superfield}, $\Gamma[\boldsymbol\Psi,\mathcal G]$ (\pkg{"Superfield"}, \pkg{"MeanFields" -> \{NambuField -> X\}}, the Nambu propagator $\mathcal G$ written \pkg{calG})}\\
$-2$ & \pkg{Phi, calG[a, b]} & $\boldsymbol\Sigma(a,b)$, with the gap functions\\
$4$ & \pkg{Phi, calG[c, d], calG[a, b]} & $\Gamma_s(a,b,c,d)$; $\Gamma_t(a,b,c,d)=-\Gamma_s(a,c,b,d)$\\
 & \pkg{Gamma, calG[a, b]} & $-\frac12(\boldsymbol\Sigma+\mathcal G^{-1}-\mathcal G_0^{-1})(a,b)$\\
 & \pkg{Gamma, X[a], X[b]} & $(\mathcal G_0^{-1}-\boldsymbol\Sigma[\mathcal G])(a,b)$\\
 & \pkg{Gamma, X[a], X[b], X[c], X[d]} & $\Gamma_{\rm irr}(a,b,c,d)$\\
\multicolumn{3}{@{}l}{\emph{its components}, at $\mathcal G_{\psi\psi}=\mathcal G_{\bar\psi\bar\psi}=0$ (\pkg{Gpp}, \pkg{Gbb}; \fnref{SuperfieldComponents})}\\
$-4$ & \pkg{Phi, Gpp[x2, x4], Gbb[x1, x3]} & $\Gamma_{pp}(x_1,x_2;x_3,x_4)$, particle--particle\\
$-2$ & \pkg{Gamma, Gpp[x2, x4], Psibar[x1], Psibar[x3]} & $\Gamma_{pp}(x_1,x_2;x_3,x_4)$\\
 & \pkg{Gamma, Psi[x2], Psi[x4], Psibar[x1], Psibar[x3]} & $-\Gamma_{\rm irr}(x_1,x_2;x_3,x_4)$\\
\bottomrule
\end{tabular}
\end{center}
With $\Phi$ to second order, the second and third derivatives of the fermionic $\Gamma[\bar\Psi,\Psi,G]$ against the table:
\excode{effective_action_means_derivatives}
\extext{effective_action_means_derivatives}
Without the anomalous propagators, the fermionic $\Gamma[\bar\Psi,\Psi,G]$ at fixed $G$ holds them at zero, so that its
fourth derivative is the vertex irreducible in all three channels, $\Gamma_{\rm irr}$ (the bare vertex, then the envelope from the fourth
order), not $\Gamma_{\rm irr}+R_{pp}$.

Options:
\begin{itemize}
\item \opt{MeanFields} \pkg{-> \{phi -> phiMean\}} (one real boson) or \pkg{\{psi -> Psi, psibar -> Psibar\}} (fermions, with
  Grassmann-odd means $\Psi=\langle\psi\rangle$, $\bar\Psi=\langle\bar\psi\rangle$; complex bosons likewise) makes the expectation values of the fields variables too ($\bar\phi=\langle\phi\rangle$ for a real boson), declared as
  sources. For a real boson this is $\Gamma[\bar\phi,G]=S[\bar\phi]+\frac12{\rm Tr}\ln G^{-1}+\frac12{\rm Tr}(D^{-1}[\bar\phi]G-1)
  +\Gamma_2[\bar\phi,G]$, with $D^{-1}[\bar\phi]=S''[\bar\phi]=D_0^{-1}+S_{\rm int}''[\bar\phi]$ the bare inverse propagator in the
  background, whose part $\frac12{\rm Tr}(S_{\rm int}''[\bar\phi]G)$ closes the vertex of degree two with one line $G$, and $\Gamma_2$
  minus the 2PI vacuum diagrams with lines $G$ and the vertices of $S_{\rm int}[\bar\phi+\phi]$ of degree three or more in $\phi$. The
  stationarity in $G$ is the Dyson equation $G^{-1}=D^{-1}[\bar\phi]-\Sigma[\bar\phi,G]$, $\Sigma$ the self-energy of $\Gamma_2$: through
  $D^{-1}$ the self-energy gets its mean-field part $-S_{\rm int}''[\bar\phi]$ (for $\phi^4$ the tadpole $-\lambda\bar\phi^2/2$, which
  with $-\lambda G(x,x)/2$ from $\Gamma_2$ is the tadpole of $\langle\phi\phi\rangle=G+\bar\phi\bar\phi$), zero at $\bar\phi=0$. For fermions
  $\Gamma[\bar\Psi,\Psi,G]=S[\bar\Psi,\Psi]-{\rm Tr}\ln(-G^{-1})-{\rm Tr}(D^{-1}[\bar\Psi,\Psi]G-1)+\Gamma_2[\bar\Psi,\Psi,G]$, with
  $S[\bar\Psi,\Psi]=-\bar\Psi G_0^{-1}\Psi+S_{\rm int}[\bar\Psi,\Psi]$ and $D^{-1}[\bar\Psi,\Psi](x,y)=\delta^2S/\delta\bar\Psi(x)\,\delta\Psi(y)
  =G_0^{-1}(x,y)-\Sigma_{\rm HF}[-\Psi\bar\Psi](x,y)$ (\pkg{FunctionalDerivative[S, Psibar[x], Psi[y]]}, the normal block of the
  Nambu matrix $S''[\bar\Psi,\Psi]$ of the second derivatives, \fnref{EffectiveAction1PI}): $\Sigma_{\rm HF}[-\Psi\bar\Psi]$ is the first-order self-energy $\Sigma_{\rm HF}[G]$ with $-\Psi(x)\bar\Psi(y)$ in
  place of $G(x,y)$, the Hartree and Fock potentials of the means, which ${\rm Tr}(\Sigma_{\rm HF}[-\Psi\bar\Psi]\,G)$ closes with one
  line $G$; the Dyson equation $G^{-1}=D^{-1}[\bar\Psi,\Psi]-\Sigma[\bar\Psi,\Psi,G]$ has at first order $\Sigma_{\rm HF}[G-\Psi\bar\Psi]$,
  the Hartree and Fock potentials of the whole $-\langle T\psi\bar\psi\rangle=G-\Psi\bar\Psi$. For an interaction
  even in the fields the derivatives with respect to the means at zero mean field and fixed $G$ are skeleton vertex
  functions: the second is $D_0^{-1}-\Sigma[G]$, the fourth the fully irreducible vertex $\Gamma_{\rm irr}$ of
  \fnref{ParquetDecomposition}. At $\bar\phi\neq0$ the perturbative $\Sigma$ in the background is selected with
  \fnref{SkeletonDiagrams}, not as a 2PI open graph (the bubble with two vertices $\lambda\bar\phi$ is a skeleton but not 2PI).
  \opt{MeanFieldDegree} \pkg{-> \{k, \dots\}} computes only the terms of these degrees in the means. With an interaction,
  means of complex bosons are not available. For fermions the form is that of the superfield at zero anomalous propagator,
  exact at zero means in the normal state (no pairing source and no spontaneous pairing: in a paired state the anomalous
  propagator is finite at zero means too); at finite means the anomalous propagator is of order $V\Psi\Psi$, and the superfield,
  \pkg{\{NambuField -> X\}} with \pkg{"Statistics" -> "Superfield"} (one Grassmann-odd mean $\boldsymbol\Psi=(\Psi,\bar\Psi)$, the
  \pkg{X} of the code, the free part $-\frac12\boldsymbol\Psi\mathcal G_0^{-1}\boldsymbol\Psi$, the trace term
  $-\frac12{\rm Tr}(S''[\boldsymbol\Psi]\mathcal G-1)$ with the whole second derivative, its anomalous vertices included), is exact there: at its stationary
  $\mathcal G[\boldsymbol\Psi]$ it is the 1PI action $\Gamma[\boldsymbol\Psi]$ of
  \fnref{EffectiveAction1PI}, whose components are $\Gamma[\bar\Psi,\Psi]$ with the diagrams of the anomalous vertices.
\item \opt{Source} \pkg{-> U} adds a one-body source to the bare inverse propagator, so that the action gains
  $\int(\text{barred field})\,U\,(\text{field})$: $G_0^{-1}\to G_0^{-1}-U$ for fermions, $G_0^{-1}+U$ for complex bosons,
  $D_0^{-1}+U$ for real bosons ($\frac12\int\phi U\phi$), $\mathcal G_0^{-1}-K$ for a superfield ($\frac12\int\boldsymbol\psi K\boldsymbol\psi$, with the
  one-body source and the pairing sources $\frac12\int\psi K_{\psi\psi}\psi$, $\frac12\int\bar\psi K_{\bar\psi\bar\psi}\bar\psi$, so that the
  stationarity is $\mathcal G^{-1}=\mathcal G_0^{-1}-K-\Sigma[\mathcal G]$). $\Gamma$ depends on $U$ only through the trace term, so that at
  the stationary point the derivative is the explicit one, $\delta\Omega/\delta U(x,y)=G(y,x)=\langle\bar\psi(x)\psi(y)
  \rangle$ ($\frac12G(x,y)$ for real bosons): $N=-\partial\Omega/\partial\mu$ although $\Sigma$ depends on $\mu$. The second
  derivative is the conserving response $\chi=\gamma L\gamma$, $L=L_0+L_0\Gamma^0L$ (\fnref{ExactEquation}
  \opt{Response}, \opt{Susceptibility}), whose static limit is the second derivative of $\Omega$
  (\fnref{EvaluateMatsubara}, option \opt{Limit}).
\item \opt{Method}, passed to \fnref{LuttingerWardFunctional}.
\end{itemize}
\excode{effective_action_source}
\exformula{effective_action_source}
$\phi^4$ theory with its mean field $\bar\phi$ (\pkg{phiMean}) and the propagator $G$ (\pkg{DressedPropagator}), to first
order: the classical action, the trace terms, the mass $\lambda\bar\phi^2/2$ closed by
$G$ in $\frac12{\rm Tr}(S_{\rm int}''[\bar\phi]G)$, and the figure eight:
\excode{effective_action_mean_field}
\exformula{effective_action_mean_field}

\begin{function}{EffectiveAction1PI}
\syntax{EffectiveAction1PI[Z0, interaction, n, phi -> phiMean, D0 -> Dm]}
\syntax{EffectiveAction1PI[Z0, interaction, n, phi -> phiMean, D0 -> D0]}
\syntax{EffectiveAction1PI[Z0, interaction, n, \{psi -> Psi, psibar -> Psibar\}, G0 -> G0]}
\syntax{EffectiveAction1PI[Z0, interaction, n, NambuField -> X, G0 -> G0, "Statistics" -> "Superfield"]}
\end{function}
The one-particle-irreducible effective action $\Gamma[\bar\phi]$ of a real boson to order $n$, the Legendre transform of
$-\log Z$ with respect to the source of $\phi$: the quantum effective action of Jona-Lasinio~\cite{JonaLasinio1964}, the
generating functional of the 1PI vertex functions; at constant $\bar\phi$ it is the volume times the effective potential
(Coleman and Weinberg at one loop~\cite{ColemanWeinberg1973}):
\[
\Gamma[\bar\phi]=S[\bar\phi]+\tfrac12{\rm Tr}\ln S''[\bar\phi]+\Gamma_1[\bar\phi],
\]
with $S''[\bar\phi]=D^{-1}[\bar\phi]=D_0^{-1}+S_{\rm int}''[\bar\phi]$ the fluctuation operator in the background and $\Gamma_1$, defined by
this equation as $\Gamma_2$ is by the 2PI action, minus the connected 1PI vacuum diagrams with two or more loops, with the
vertices of $S_{\rm int}[\bar\phi+\phi]$ of degree three or more and the lines $D=S''[\bar\phi]^{-1}$, the propagator in the background,
written with the head \pkg{Dm}, whose derivative with respect to $\bar\phi$ is declared ($\delta D/\delta\bar\phi=-D\,S_{\rm int}^{(3)}D$). Its derivatives at $\bar\phi=0$ are the
vertex functions of \fnref{ConnectedFunctionFromVertices}: $\Gamma^{(2)}=D_0^{-1}-\Sigma$, and $\Gamma^{(n)}$ is minus the
amputated 1PI $n$-point function, whose skeletons are \fnref{PerturbativeFullVertex}. It is \fnref{EffectiveAction2PI} with
its mean field at the stationary $G[\bar\phi]$; at $G=D$, $\Gamma_2[\bar\phi,D]$ is the part of $\Gamma_1$ 2PI in the lines $D$.
\excode{effective_action_1pi}
\exformula{effective_action_1pi}
The vertex functions are its derivatives at $\bar\phi=0$, where $D=D_0$ (replace \pkg{Dm} by \pkg{D0}, and contract
$D\,D^{-1}$ with \fnref{ExpandTrace}: before it, a picture draws $D^{-1}$ as a line, with the message
\pkg{inversekernel}). With the lines $D$ a derivative brings a coupling, since $\delta D/\delta\bar\phi$
does: the derivatives of the action of order $n$ hold terms of higher orders, not all of them complete, and only the
orders up to $n$ are kept (\fnref{OrderPart}); with bare lines, \pkg{D0 -> D0}, the derivatives keep the order of the action.
To first order, $\Gamma^{(2)}=D_0^{-1}-\Sigma$ with the tadpole:
\excode{effective_action_1pi_vertex}
\exformula{effective_action_1pi_vertex}
With bare lines, \pkg{D0 -> D0}, the vertices of degree two in the fluctuation stay vertices: $\Gamma[\bar\phi]=S[\bar\phi]+
\frac12{\rm Tr}\ln D_0^{-1}$ minus the connected 1PI vacuum diagrams with lines $D_0$ and the vertices of degree two or more, the
same series written in $D_0$: its diagrams with one loop are the rings of $\frac12{\rm Tr}\ln S''[\bar\phi]-\frac12{\rm Tr}\ln D_0^{-1}$,
those with two or more $\Gamma_1$ with its lines $D$ expanded in $D_0$. For fermions, with the Grassmann-odd means
$\Psi=\langle\psi\rangle$, $\bar\Psi=\langle\bar\psi\rangle$ (declared as Grassmann sources), the action has the same form, with the
signs of a Grassmann integral:
\[
\Gamma[\bar\Psi,\Psi]=S[\bar\Psi,\Psi]-\tfrac12{\rm Tr}\ln\bigl(-S''[\bar\Psi,\Psi]\bigr)+\Gamma_1[\bar\Psi,\Psi].
\]
$S''$ is the Nambu matrix of the second derivatives with respect to $(\Psi,\bar\Psi)$, the trace running over its components: its
normal block is $D^{-1}[\bar\Psi,\Psi](x,y)=\delta^2S/\delta\bar\Psi(x)\,\delta\Psi(y)$, and its anomalous blocks,
$\delta^2S/\delta\bar\Psi(x)\,\delta\bar\Psi(y)=V(x,y)\Psi(x)\Psi(y)$ and $\delta^2S/\delta\Psi(x)\,\delta\Psi(y)=V(x,y)\bar\Psi(x)\bar\Psi(y)$,
come from the vertices $V\Psi\Psi\bar\psi\bar\psi$ and $V\bar\Psi\bar\Psi\psi\psi$; the lines of $\Gamma_1$ are $S''[\bar\Psi,\Psi]^{-1}$, the
propagator in the background, with anomalous blocks as in the Nambu--Gor'kov formalism. Such a line is not one directed line,
and for fermions the action is written with bare lines, its only form: $S[\bar\Psi,\Psi]-{\rm Tr}\ln(-G_0^{-1})$ minus the connected
1PI vacuum diagrams with lines $G_0$ and the vertices of $S_{\rm int}[\bar\Psi+\bar\psi,\Psi+\psi]$ of degree two or more, the anomalous
ones included, whose diagrams with one loop are the rings of $-\frac12{\rm Tr}\ln(-S''[\bar\Psi,\Psi])+{\rm Tr}\ln(-G_0^{-1})$ and those
with two or more $\Gamma_1$ with its lines expanded in $G_0$. At zero means
the derivatives, $\Psi$ first and then $\bar\Psi$, give $G_0^{-1}-\Sigma$ with the perturbative self-energy, and
$\delta^4\Gamma/\delta\bar\Psi(x_1)\,\delta\Psi(x_2)\,\delta\bar\Psi(x_3)\,\delta\Psi(x_4)$ is the full vertex $F(x_1,x_2;x_3,x_4)$ with bare
lines, whose skeletons are \fnref{PerturbativeFullVertex}; in code, with the rightmost derivative first,
\pkg{FunctionalDerivative[Gamma, Psibar[a], Psi[b]]} is $G_0^{-1}(a,b)-\Sigma(a,b)$ and
\pkg{FunctionalDerivative[Gamma, Psibar[x1], Psi[x2], Psibar[x3], Psi[x4]]} is $F(x_1,x_2;x_3,x_4)$. \fnref{EffectiveAction2PI} with its means at the stationary $G$ is
the functional of Nambu and Gor'kov at zero anomalous propagator: to third order it differs from $\Gamma[\bar\Psi,\Psi]$ by
the diagrams with an anomalous vertex, which the superfield recovers. For a superfield (\pkg{"Statistics" -> "Superfield"},
the mean $\boldsymbol\Psi=\langle\boldsymbol\psi\rangle$, \pkg{X} in \pkg{NambuField -> X}) the 1PI action is
$\Gamma[\boldsymbol\Psi]=S[\boldsymbol\Psi]-\frac12{\rm Tr}\ln(-S''[\boldsymbol\Psi])+\Gamma_1[\boldsymbol\Psi]$, with
$S[\boldsymbol\Psi]=-\frac12\boldsymbol\Psi\mathcal G_0^{-1}\boldsymbol\Psi+S_{\rm int}[\boldsymbol\Psi]$, written with bare lines,
$S[\boldsymbol\Psi]-\frac12{\rm Tr}\ln(-\mathcal G_0^{-1})$ minus the connected 1PI vacuum diagrams with lines $\mathcal G_0$ and the
vertices of degree two or more; its components (\fnref{SuperfieldComponents}) are
$\Gamma[\bar\Psi,\Psi]$, and it is the superfield \fnref{EffectiveAction2PI} at its stationary $\mathcal G[\boldsymbol\Psi]$. Other lines for fermions
and superfields, a single Grassmann field or complex bosons give the message \pkg{onepi} and \pkg{\$Failed}. The electron gas to first order:
\excode{effective_action_1pi_fermions}
\exformula{effective_action_1pi_fermions}
and its second derivative at zero means, $G_0^{-1}-\Sigma$ with the Hartree and Fock self-energy:
\excode{effective_action_1pi_fermions_derivative}
\exformula{effective_action_1pi_fermions_derivative}

\subsection{Potthoff's functionals}

\begin{function}{LegendreTransform}
\syntax{LegendreTransform[Phi, G -> S]}
\end{function}
Potthoff's universal functional~\cite{Potthoff2003} $\mathcal F[S]=\Phi[G[S]]-{\rm Tr}(S\,G[S])$ for fermions ($+\frac12{\rm Tr}(SG)$
for real bosons, $+{\rm Tr}(SG)$ for complex bosons), the Legendre transform of the Luttinger--Ward functional from $G$ to
the self-energy, $G[S]$ being the propagator at which $\Sigma[G]=S$. It is inert: its derivative with respect to $S(x,y)$
is $-G[S](y,x)$ (\fnref{LegendrePropagator}); with respect to another kernel of $\Phi$, an interaction $V$, it is
$\delta\Phi/\delta V$ at $G=G[S]$, since $\mathcal F$ is stationary in $G$ (in the code $S$ is \pkg{Sg}, written $\Sigma$):
\excode{legendre_interaction}
\exformula{legendre_interaction}
$\mathcal F$ depends on the interaction only: the diagrams of $\Phi$ contain $G$ and $V$, never $G_0$.

\begin{function}{LegendrePropagator}
\syntax{LegendrePropagator[G -> S, \{V, \dots\}][x, y]}
\end{function}
$G[S](x,y)$, the propagator at which the self-energy of $\Phi$ (a functional of $G$ and of the kernels \pkg{V}, \dots) is
$S$; shown as $G[S](x,y)$, and written as $G(x,y)$ by \fnref{ExpandTrace} and \fnref{ZeroDimensionalValue}. Its
derivatives are not available ($\delta G[S]/\delta S$ is the inverse of the irreducible vertex): \fnref{FunctionalDerivative}
gives the message \pkg{LegendrePropagator::derivative} and a \pkg{\$Failed} term.

\begin{function}{SelfEnergyFunctional}
\syntax{SelfEnergyFunctional[Phi, G0 -> S, G]}
\end{function}
Potthoff's self-energy functional $\Omega_{G_0}[S]=\mathcal F[S]-{\rm Tr}\ln\bigl(-(G_0^{-1}-S)\bigr)$ for fermions
($\mathcal F[S]+\frac12{\rm Tr}\ln(D_0^{-1}-S)$ for real bosons, $\mathcal F[S]+{\rm Tr}\ln(G_0^{-1}-S)$ for complex bosons). $G_0$ enters only
through the trace-log: Potthoff writes $\Omega_t[\Sigma]$, his $t$ being the one-body part. Since $\mathcal F$ is the same for every
system with the same interaction, it can be computed exactly on a simpler reference system (a cluster, an impurity): the
idea of the variational cluster approximation and of dynamical mean-field theory. $\Omega_{G_0}$ is stationary where
$G[S]=(G_0^{-1}-S)^{-1}$, the Dyson equation, and equal to $-\log Z$ there. It is a Legendre transform, not the
substitution $G\to(G_0^{-1}-S)^{-1}$ in $\Gamma[G]$, which is stationary at the same point but has $G_0$ also inside
$\Phi$:
\excode{self_energy_functional}
\exformula{self_energy_functional}
\excode{self_energy_functional_derivative}
\exformula{self_energy_functional_derivative}

\subsection{Zero dimensions}

\begin{function}{ZeroDimensionalValue}
\syntax{ZeroDimensionalValue[expr, \{G -> g, V -> u, \dots\}]}
\syntax{ZeroDimensionalValue[expr, rules, "LoopFactor" -> n, "Loops" -> \{G, \dots\}]}
\syntax{ZeroDimensionalValue[expr, \{G -> matrix, V -> array, \dots\}, "Indices" -> \{x -> i, \dots\}]}
\end{function}
The value of \pkg{expr} in zero dimensions (one point, possibly with internal states), where every integral can be done:
the integrals are dropped, $\delta\to1$, every kernel is replaced by its value, and every closed loop of kernels and every
trace gets the factor \opt{LoopFactor} ($2s+1$; default 1); ${\rm Tr}\,A\to n\,a$, ${\rm Tr}\ln A\to\ln a^n$.
\fnref{LegendreTransform} is written out, with $G\to g(S)$ among the rules. \opt{Loops} gives the heads of the lines whose
loops are counted; by default (\pkg{Automatic}) every kernel of two labels that is neither symmetric nor an interaction
line, and the kernels of operators made of them. For one level with spin $\frac12$ ($Z=a^2-u$, $G_0=1/a$) the
self-energy and the effective action to second order:
\excode{zero_dimensional_value}
\extext{zero_dimensional_value}
When the values are arrays (a matrix for a kernel of two labels, an array of rank four for a vertex) the labels run over the
internal states: the integrals are sums, $\delta$ is the identity, operators are matrices, ${\rm Tr}\ln A\to\ln\det A$, and
\opt{Indices} gives the states of the free labels; so the superfield of one site with spin $\frac12$, whose $4\times4$ Nambu
matrix has the anomalous entries $\kappa$ and $\bar\kappa$ (\pkg{pair}, \pkg{pairb}), is checked against the exact $Z=a^2+k\bar k-u$ with the pairing
sources. Its $\Phi$ (\pkg{PhiNambu} of \fnref{SuperfieldComponents}) to second order, a function of $g^2-\kappa\bar\kappa$:
\excode{zero_dimensional_arrays}
\exformula{zero_dimensional_arrays}
The skeleton series of $\Sigma[g]$ has the Catalan numbers as coefficients and converges, for $4ug^2<1$, to one of the two
branches $\Sigma_\pm(g)=(-1\pm\sqrt{1+4ug^2})/(2g)$ of the exact $\Sigma[g]$; at strong repulsion the physical self-energy is
on the other one, and where the branches meet the irreducible vertex
diverges~\cite{KozikFerreroGeorges2015,RossiWerner2015}. The showcase shows this exactly, and its relation to the
divergences of the irreducible vertex found in the Hubbard model~\cite{Schaefer2013}.

% ------------------------------------------------------------------ 11
\section{Two-particle channels, Bethe--Salpeter ladders and the parquet}\label{sec:parquet}

\subsection{Channels}

The four legs of a vertex can be divided into two pairs in three ways~\cite{Bickers2004}; in each channel the
vertex obeys a Bethe--Salpeter equation, and the decomposition follows from the Legendre transforms of the 2PI
formalism~\cite{DeDominicisMartin1964b}. With the fermionic convention (legs $x_1$, $x_3$
outgoing, $x_2$, $x_4$ incoming), and for real bosons:
\begin{center}
\small
\begin{tabular}{@{}lll@{\hspace{3em}}ll@{}}
\toprule
fermions & & & bosons & \\
\midrule
\opt{ParticleHole} & $\ph$ & $(x_1x_2|x_3x_4)$ & \opt{s} & $(x_1x_2|x_3x_4)$\\
\opt{TransverseParticleHole} & $\phbar$ & $(x_1x_4|x_3x_2)$ & \opt{t} & $(x_1x_3|x_2x_4)$\\
\opt{ParticleParticle} & $\pp$ & $(x_1x_3|x_2x_4)$ & \opt{u} & $(x_1x_4|x_2x_3)$\\
\bottomrule
\end{tabular}
\end{center}
A superfield (\pkg{"Statistics" -> "Superfield"}: one Grassmann field with an antisymmetric propagator, such as the
Nambu field $(\psi,\bar\psi)$) has the channels of real bosons, $s$, $t$, $u$. A diagram is reducible in a channel if
cutting two propagators separates its pairs. By the parquet theorem a diagram is
reducible in at most one channel, so that $F=\Gamma_{\rm irr}+\sum_r R_r$, with $\Gamma_{\rm irr}$ fully irreducible (often written
$\Lambda$; here $\Lambda$ is the vertex of the density) and $R_r$ the diagrams
reducible in the channel $r$; $\Gamma_r=F-R_r$ is the vertex irreducible in $r$, and $R_r=\Gamma_r\circ_r F$ (Bethe--Salpeter
equation).

\subsection{Functions}

\begin{function}{ParquetDecomposition}
\syntax{ParquetDecomposition[F, \{x1, x2, x3, x4\}, \{G, \dots\}]}
\syntax{ParquetDecomposition[F, \{x1, x2, x3, x4\}, \{G, \dots\}, "Statistics" -> "Boson"]}
\syntax{ParquetDecomposition[F, \{a, b, c, d\}, \{G\}, "Statistics" -> "Superfield"]}
\end{function}
Classifies the diagrams of the vertex \pkg{F} by two-particle reducibility, cutting the lines \pkg{G}. The result is an
association with the keys \opt{FullyIrreducible}, the three channels of the statistics, and \opt{SeveralChannels} (the
diagrams reducible in more than one channel, which the parquet theorem excludes: always 0 for a vertex built by the
package). The parts it finds by cutting lines are derivatives: \opt{FullyIrreducible} is the fourth derivative of the 2PI
action with respect to the means at fixed $G$ (\fnref{EffectiveAction2PI} with \opt{MeanFields}), and $F-R_r$, the vertex
irreducible in the channel $r$, the second derivative of $\Phi$ (\fnref{IrreducibleVertex}); the classification checks them.
\excode{parquet_phi4}
\extext{parquet_phi4}

\begin{function}{BetheSalpeterLadder}
\syntax{BetheSalpeterLadder[L, R, \{x1, x2, x3, x4\}, channel, G]}
\syntax{BetheSalpeterLadder[L, R, \{x1, x2, x3, x4\}, channel, G, "MaximumOrder" -> \{V, n\}]}
\end{function}
The contraction $L\circ R$ of two four-point vertices, expressions in the labels \pkg{x1, \dots, x4}, with a pair of
propagators \pkg{G} in the given channel, integrated over the internal points $5,6,7,8$ ($G(a,b)$ runs from $b$ to $a$):
\begin{center}
\small
\begin{tabular}{@{}ll@{}}
\toprule
\opt{ParticleHole} & $\phantom{-\tfrac12}\,L(x_1,x_2;5,6)\,G(6,7)\,G(8,5)\,R(7,8;x_3,x_4)$\\
\opt{TransverseParticleHole} & $-\phantom{\tfrac12}\,L(x_1,6;5,x_4)\,G(6,7)\,G(8,5)\,R(7,x_2;x_3,8)$\\
\opt{ParticleParticle} & $-\tfrac12\,L(x_1,5;x_3,6)\,G(5,7)\,G(6,8)\,R(7,x_2;8,x_4)$\\
\opt{s} (bosons) & $-\tfrac12\,L(x_1,x_2,5,6)\,G(5,7)\,G(6,8)\,R(7,8,x_3,x_4)$\\
\opt{t} (bosons) & $-\tfrac12\,L(x_1,5,x_3,6)\,G(5,7)\,G(6,8)\,R(7,x_2,8,x_4)$\\
\opt{u} (bosons) & $-\tfrac12\,L(x_1,5,6,x_4)\,G(5,7)\,G(6,8)\,R(7,x_2,x_3,8)$\\
\opt{s}, \opt{t}, \opt{u} (superfields) & the same with the factors $+\tfrac12$, $-\tfrac12$, $+\tfrac12$\\
\bottomrule
\end{tabular}
\end{center}
With these conventions every channel obeys $R_r=\Gamma_r\circ_r F$ and $F=\Gamma_{\rm irr}+\sum_r R_r$; for fermions
$\Gamma_{\ph}=\Gamma^0$. For $\phi^4$ theory, with $F=\Gamma^{(4)}=\lambda\,\delta\delta\delta+O(\lambda^2)$ (Euclidean), the
parquet equations read $F=\Gamma_{\rm irr}-\frac12\sum_{r=s,t,u}\Gamma_r\,L_0\,F$, with $L_0$ the pair of lines of the channel $r$
and the factor $\frac12$ that of a pair of identical lines; $\Gamma_s=4\,\delta^2\Phi/\delta G\,\delta G$ and
$\Gamma_{\rm irr}=\delta^4\Gamma/\delta\bar\phi^4$ at
fixed $G$ (section~\ref{sec:effective}). With \opt{MaximumOrder} \pkg{-> \{V, n\}} only the products of total order $n$ or less are
built. The direct interaction iterated in the particle--hole channel gives a bubble:
\excode{ladder_ph}
\exboth{ladder_ph}

\begin{function}{CheckParquetEquations}
\syntax{CheckParquetEquations[F, \{x1, x2, x3, x4\}, G, coupling, n]}
\syntax{CheckParquetEquations[F, \{x1, x2, x3, x4\}, G, x, n, "Statistics" -> "Boson"]}
\syntax{CheckParquetEquations[F, \{a, b, c, d\}, G, V, n, "Statistics" -> "Superfield"]}
\end{function}
Checks, order by order up to $n$ in \pkg{coupling} (a coupling symbol or the head of the interaction lines), that every diagram
of \pkg{F} is reducible in at most one channel and that the reducible part of each channel equals the ladder
$\Gamma_r\circ_r F$ with $\Gamma_r=F-R_r$. The result is an association: \opt{Residuals} (for each channel, the list of
the residuals of orders $1,\dots,n$, all 0 when the equations hold), \opt{DiagramCounts} (the numbers of diagrams per
order in each class), \opt{SeveralChannels} and \opt{FullyIrreducible}.
\excode{check_parquet_phi4}
\extext{check_parquet_phi4}

\subsection{Superfields}\label{sec:superfields}

The Nambu field $\boldsymbol\psi=(\psi,\bar\psi)$ is one Grassmann field with an antisymmetric propagator, the Nambu
propagator $\mathcal G$; its functional $\Phi[\mathcal G]$ (\fnref{LuttingerWardFunctional}, \fnref{EffectiveAction2PI}
with \pkg{"Statistics" -> "Superfield"}) gives by its derivatives the normal and anomalous self-energies and the kernels of
all three channels (the conventions are in section~\ref{sec:conventions}). Its components are found as follows.

\begin{function}{SuperfieldComponents}
\syntax{SuperfieldComponents[expr, \{a -> psibar[x1], b -> psi[x2], \dots\}, rules]}
\syntax{rules = \{NambuField -> \{psi, psibar\}, h -> \{G, Gpp, Gbb\}, v -> S, \dots\}}
\end{function}
The components of an expression of a superfield $\boldsymbol\psi_a=(\psi,\bar\psi)$, an ordinary Grassmann field of the
package whose label $a=(x,s)$ stands for a point and a component, with the antisymmetric propagator
$\mathcal G_{ab}=-\langle T\boldsymbol\psi_a\boldsymbol\psi_b\rangle$ (\pkg{"Statistics" -> "Superfield"} in
\fnref{SelfEnergy}, \fnref{IrreducibleVertex} and the parquet functions). The external labels are given with their component
and a new name; the rules give the components of the field (required), of its mean, of the two-label kernels (a propagator
\pkg{h -> \{G, Gpp, Gbb\}}, $\mathcal G=\left(\begin{smallmatrix}\mathcal G_{\psi\psi}&G\\-G^T&\mathcal G_{\bar\psi\bar\psi}\end{smallmatrix}\right)$,
or a matrix of heads, zeros and functions of two labels) and of the vertices (the component interaction, whose derivatives
are the components); the rule of a mean is \pkg{X -> \{Psi, Psibar\}}. Traces and kernels of operators without anomalous
blocks are written with the normal blocks. With anomalous blocks, a trace of products is written as rings
(\fnref{ExpandTrace}), while ${\rm Tr}\ln(-\mathcal G^{-1})$ has no component form: it is kept as it is, with the message
\pkg{trace}, beside the components of the other terms (it does not depend on the means; its derivatives with respect to
$G$ and the anomalous blocks are those of the superfield, $\delta{\rm Tr}\ln(-\mathcal G^{-1})/\delta\mathcal G=\mathcal G^{-1}$,
taken first and then projected at zero anomalous blocks: a derivative of the kept trace with respect to a component gives
the message \pkg{supertrace}), so that the 2PI action
$\Gamma[\boldsymbol\Psi,\mathcal G]$ of \fnref{EffectiveAction2PI} has the components $\Gamma[\bar\Psi,\Psi,G,\mathcal G_{\psi\psi},\mathcal
G_{\bar\psi\bar\psi}]$, whose derivatives with respect to the means and the blocks are the vertices of the table of
\fnref{EffectiveAction2PI}. The order of the labels of a derivative is the order of the legs: with the legs
\pkg{a -> psibar[x1]}, \pkg{b -> psi[x2]}, \pkg{c -> psibar[x3]}, \pkg{d -> psi[x4]}, the derivative with respect to
\pkg{X[a], X[b], X[c], X[d]} projects onto that with respect to \pkg{Psibar[x1], Psi[x2], Psibar[x3], Psi[x4]}, and an odd
permutation of the four Grassmann derivatives changes the sign. With
it the gap function $\Delta=\Sigma_{\bar\psi\bar\psi}$ and the particle--particle kernel (the channel \pkg{"t"} of
\fnref{IrreducibleVertex}) are one call away. The superfield of the electron gas is written as $\phi^4$ theory, with one
source, one antisymmetric propagator $\mathcal G$ and one antisymmetric vertex $\mathcal V$ (\pkg{NambuPropagator},
\pkg{NambuVertex}), drawn as a point standing for $-\mathcal V$; its $\Phi$ has
one diagram at each order:
\excode{superfield_phi}
\exboth{superfield_phi}
Its components at first order are the Hartree and Fock diagrams and the pairing term, an anomalous line of each kind,
$\mathcal G_{\psi\psi}$ and $\mathcal G_{\bar\psi\bar\psi}$ (\pkg{AnomalousGreenFunction}, \pkg{AnomalousGreenFunctionBar}; the
weights read them along the fermion line, section~\ref{sec:weights}):
\excode{superfield_components}
\exboth{superfield_components}
The particle--particle kernel of the electron gas is the second derivative of these components with respect to the
anomalous propagators, at zero anomalous propagators; at first order the bare vertex:
\excode{superfield_pp_derivative}
\exboth{superfield_pp_derivative}
The self-energy of the superfield is $\boldsymbol\Sigma=-2\,\delta\Phi/\delta\mathcal G$, and its $\bar\psi\bar\psi$
component at first order is the gap function, $\Delta(x_1,x_2)=V(x_1,x_2)\,\mathcal G_{\psi\psi}(x_1,x_2)$, the BCS gap
equation (the shortcut is \fnref{SelfEnergy} with the option \pkg{"Statistics" -> "Superfield"}):
\excode{superfield_gap}
\exboth{superfield_gap}
The kernels of the three channels from the same $\Phi$: $\Gamma_s=4\,\delta^2\Phi/\delta\mathcal G_{ab}\,\delta\mathcal G_{cd}$, and
$\Gamma_t$, $\Gamma_u$ the same function with its labels permuted (with the sign of the permutation), their numbers of
diagrams at second order and the residuals of the shortcut \fnref{IrreducibleVertex} with \pkg{"Channel" -> All}; then
the channel \pkg{"t"} projected in the normal state, the particle--particle kernel found above, to second order:
\excode{superfield_channels}
\extext{superfield_channels}
\excode{superfield_pp_kernel}
\extext{superfield_pp_kernel}

% ------------------------------------------------------------------ 12
\section{Exact equations}\label{sec:exact}

The full objects of many-body theory --- $G$, $\Sigma$, the vertices $\Gamma^0$ and $F$, the vertex $\Lambda$ of the
density, the polarization $P$ and the screened interaction $W$ --- obey exact equations. \fnref{ExactEquation} writes them
with the indices contracted and the signs of the conventions of section~\ref{sec:conventions}; the objects are heads (for
example blobs declared by \fnref{DeclareBlob}, to draw the equation) or functions of their labels (\fnref{LabelFunction}
of a perturbative series, to check the equation order by order with \fnref{CheckEquation}).

\begin{function}{PerturbativeDensityVertex}
\syntax{PerturbativeDensityVertex[Z0, V, n, \{psi, psibar\}, \{x1, x2, z\}, G0 -> G]}
\syntax{PerturbativeDensityVertex[Z0, V, n, \{phi\}, \{x1, x2, z\}, D0 -> G]}
\end{function}
The vertex of the density at $z$, $\Lambda(x_1,x_2;z)=-\delta G^{-1}(x_1,x_2)/\delta U(z)$ for a potential $U$ coupled to
$\rho=\bar\psi\psi$, so that $\langle\psi(x_1)\bar\psi(x_2)\rho(z)\rangle_c=L_0\Lambda$, with $L_0$ the pair of propagators
$G(x_1,a)\,G(b,x_2)$ (integrated over $a$, $b$): the bare vertex $\delta(x_1,z)\,\delta(x_2,z)$ and the skeleton diagrams to
order $n$, with \pkg{G0} replaced by \pkg{G}. For a real boson the operator is $\rho=\phi^2/2$, a local mass term. The
vertex is the derivative that defines it: with $G^{-1}=G_0^{-1}-U-\Sigma$, $\Lambda=\gamma+\delta\Sigma/\delta U$, the
perturbative self-energy differentiated through its bare lines (the potential \pkg{LocalPotential} of
\fnref{ConnectedCorrelationFunction}), its skeleton diagrams with $G_0\to G$ ($\Lambda=\gamma-\delta\Sigma/\delta U$ for a real
boson, with $D_0^{-1}\to D_0^{-1}+U$). To first order, for the electron gas:
\excode{density_vertex}
\expictures{density_vertex}
This function computes the same from the three-point function:
\excode{density_vertex_derivative}
\extext{density_vertex_derivative}

\begin{function}{LabelFunction}
\syntax{LabelFunction[e, \{x1, \dots, xn\}]}
\end{function}
The function of $n$ labels that gives \pkg{e} with \pkg{x1}, \dots, \pkg{xn} replaced by its arguments.

\begin{function}{ExactEquation}
\syntax{ExactEquation[name, \{"G" -> G, "Sigma" -> S, \dots\}, labels]}
\syntax{ExactEquation[name, objects, labels, "Statistics" -> "Boson"]}
\syntax{ExactEquation[name, objects, labels, "MaximumOrder" -> \{V, n\}]}
\syntax{ExactEquation[name, \{\dots, "BareVertex" -> gamma\}, labels]}
\end{function}
The exact equation \pkg{lhs == rhs}, both sides in canonical form. Fermions leave a four-point vertex at $x_1$, $x_3$ and
enter at $x_2$, $x_4$; $\gamma(x_1,x_2;z)=\delta(x_1,z)\,\delta(x_2,z)$ is the bare vertex of the density.
\begin{center}
\small
\begin{tabularx}{\textwidth}{@{}l>{\ttfamily\small\raggedright\arraybackslash}p{3.6cm}>{\raggedright\arraybackslash}X@{}}
\toprule
name & \textrm{objects} & equation\\
\midrule
\opt{Dyson} & G, G0, Sigma & $G=G_0+G_0\Sigma G$\\
\opt{BetheSalpeter} & F, Gamma0, G & $F=\Gamma^0+\Gamma^0\circ F$ (option \opt{Channel})\\
\opt{Parquet} & F, Lambda, GammaPh, GammaPhbar, GammaPp, G & $F=\Gamma_{\rm irr}+\sum_r\Gamma_r\circ_rF$\\
\opt{SchwingerDyson} & Sigma, G, F, V & $\Sigma=\Sigma_{\rm H}+\Sigma_{\rm F}-\int V\,G\,L_0\,F$\\
\opt{Vertex} & Lambda, Gamma0, G & $\Lambda=\gamma+\Gamma^0L_0\Lambda$\\
\opt{VertexFromF} & Lambda, F, G & $\Lambda=\gamma+FL_0\gamma$\\
\opt{SelfEnergyFromVertex} & Sigma, G, V, Lambda & $\Sigma=\Sigma_{\rm H}-\int V\,G\,\Lambda$\\
\opt{Hedin} & Sigma, G, W, P, Lambda, GammaXC, V & $\Sigma=\Sigma_{\rm H}-G\,W\tilde\Lambda$, $P=L_0\tilde\Lambda$,
$W=V+VPW$, $\tilde\Lambda=\gamma+\Gamma_{\rm xc}L_0\tilde\Lambda$\\
\opt{Response} & L, Gamma0, G & $L=L_0+L_0\Gamma^0L$\\
\opt{Susceptibility} & Chi, L & $\chi=\gamma\,L\,\gamma$\\
\bottomrule
\end{tabularx}
\end{center}
The labels are \pkg{\{x, y\}} for \opt{Dyson}, \opt{SchwingerDyson} and \opt{SelfEnergyFromVertex},
\pkg{\{x1, x2, x3, x4\}} for \opt{BetheSalpeter}, \opt{Parquet} and \opt{Response}, \pkg{\{x1, x2, z\}} for
\opt{Vertex} and \opt{VertexFromF}, \pkg{\{x, y, z\}} for \opt{Hedin} and \pkg{\{z1, z2\}} for \opt{Susceptibility}. For real bosons the channel vertices of \opt{Parquet} are
\pkg{Gammas}, \pkg{Gammat}, \pkg{Gammau}. In \opt{Parquet}
the object \opt{Lambda} is the fully irreducible four-point vertex $\Gamma_{\rm irr}$ of section~\ref{sec:parquet}; in the other
rows $\Lambda(x_1,x_2;z)$ is the three-point vertex of the density, $L_0$ the pair of propagators of the channel (closed
at $z$ in $P$), and in \opt{Hedin}~\cite{Hedin1965} the object \opt{Lambda} is its proper part
$\tilde\Lambda$ (irreducible with respect to a single interaction line), with
$\Gamma_{\rm xc}=\delta\Sigma_{\rm xc}/\delta G$, $\Gamma^0$ without its Hartree part. For real bosons (\opt{Statistics}
\pkg{-> "Boson"}) the two lines of a pair carry the factor $-\frac12$; \opt{SchwingerDyson} is that of $\phi^4$ theory,
$\Sigma=-\frac\lambda2G_{xx}\delta_{xy}+\frac\lambda6\int G\,L_0\,F$, with the object \opt{Coupling}. With
\opt{MaximumOrder} \pkg{-> \{V, n\}} both sides keep only their parts of order $n$ or less in $V$ (the ladders of
\opt{BetheSalpeter} and \opt{Parquet} only the products of total order $n$ or less): enough for \fnref{CheckEquation} to
order $n$, and much faster, since the products of two series are not brought into canonical form beyond that order. The
Dyson equation drawn with blobs:
\excode{exact_dyson}
\expictures{exact_dyson}
\paragraph{Responses} \opt{Response} is the Bethe--Salpeter equation of the response of $G$ to a one-body source $U$,
$L=L_0+L_0\Gamma^0L$, with $L_0$ the pair of propagators. For fermions, with the source $\int\bar\psi\,U\psi$ and the
legs of a four-point vertex (leaving at $x_1$, $x_3$, entering at $x_2$, $x_4$),
$L(x_1,x_2;x_3,x_4)=\delta G(x_1,x_2)/\delta U(x_4,x_3)=-\langle\psi(x_1)\bar\psi(x_2)\psi(x_3)\bar\psi(x_4)\rangle_c
+G(x_1,x_4)\,G(x_3,x_2)$ and
\[
L(x_1,x_2;x_3,x_4)=G(x_1,x_4)\,G(x_3,x_2)+\int G(x_1,a)\,G(b,x_2)\,\Gamma^0(a,b;c,d)\,L(d,c;x_3,x_4);
\]
for real bosons, with $\frac12\int\phi\,U\phi$, $L=-2\,\delta G/\delta U=\langle\phi_1\phi_2\phi_3\phi_4\rangle_c+L_0$,
with $L_0=G_{13}G_{24}+G_{14}G_{23}$ the two pairs of propagators, and
\[
L(x_1,x_2;x_3,x_4)=L_0(x_1,x_2;x_3,x_4)-\frac12\int G(x_1,a)\,G(x_2,b)\,\Gamma^0(a,b;c,d)\,L(c,d;x_3,x_4).
\]
\opt{Susceptibility} contracts $L$ with $\gamma$, the bare vertex of a one-body operator
$O(z)=\int\bar\psi(a)\,\gamma(a,b;z)\,\psi(b)$ ($\frac12\phi\gamma\phi$ for real bosons, with a factor $-\frac14$),
\[
\chi(z_1,z_2)=\delta\langle O(z_1)\rangle/\delta u(z_2)=\delta^2\Omega/\delta u(z_1)\,\delta u(z_2)=-\langle O(z_1)\,
O(z_2)\rangle_c,
\]
with the sign of the Lindhard function. Without \opt{BareVertex} the operator is the density; \pkg{"BareVertex"
-> gamma}, a function of three labels (a current, a spin density), is also taken by \opt{Vertex} and \opt{VertexFromF}. With
the kernel $\Gamma^0=\delta\Sigma/\delta G$ of a truncated $\Phi$ the response is that of a conserving approximation
(section~\ref{sec:effective}):
\excode{exact_response}
\expictures{exact_response}

\begin{function}{CheckEquation}
\syntax{CheckEquation[lhs == rhs, x, n]}
\syntax{CheckEquation[\{eq1, eq2, \dots\}, x, n]}
\end{function}
The residuals of orders $0,\dots,n$ in \pkg{x} (a coupling symbol or the head of the interaction lines) of an exact
equation, all 0 when it holds to order $n$. The vertex equation of $\phi^4$ theory to second order:
\excode{check_vertex_phi4}
\extext{check_vertex_phi4}

\subsection{Auxiliary fields}\label{sec:auxiliary}

\begin{function}{HubbardStratonovich}
\syntax{HubbardStratonovich[\{Z0, interaction\}, O, \{phi, J\}, D]}
\end{function}
Decouples an interaction $\frac12\int_{x,y}K(x,y)\,O(x)\,O(y)$, with $O$ a local operator of even parity given as a
function of one label (\pkg{psibar[\#] ** psi[\#] \&}, \pkg{phi[\#]\^{}2 \&}), by a real auxiliary field $\phi$ with
$\langle\phi\phi\rangle_0=K$ and the coupling $-i\int\phi\,O$: $\langle e^{i\int\phi O}\rangle_\phi=e^{-\frac12\int OKO}$
(Hubbard--Stratonovich transformation~\cite{Stratonovich1958,Hubbard1959}). The kernel is read off the interaction --- terms with two labels, or local terms
such as $\frac\lambda{4!}\phi^4=\frac12\cdot\frac\lambda{12}(\phi^2)^2$ --- and \pkg{D} is the head of the propagator of
$\phi$, standing for $K$; with \pkg{D = Automatic} and $K=c\,h(x,y)$ the head $h$ itself is the propagator and
$\sqrt c$ goes into the coupling. Terms of another form are kept, with a message. The result is an association with the
\opt{FreeGeneratingFunctional} (the old one times the Gaussian of $\phi$), the \opt{Interaction}, the \opt{Propagator} and
the \opt{Kernel} (a function of two labels, to substitute $D$ back); the field and its source are declared. For the electron
gas:
\excode{hs_electron_gas}
\exformula{hs_electron_gas}

\begin{function}{IntegrateOutField}
\syntax{IntegrateOutField[\{Z0, interaction\}, phi, J]}
\end{function}
The converse: a real Gaussian field with $\langle\phi\phi\rangle_0=D$ (read off \pkg{Z0}) that enters the interaction
linearly, $\int\phi(z)A(z)$, is integrated out, $\langle e^{-\int\phi A}\rangle_\phi=e^{\frac12\int ADA}$, and the
interaction becomes $-\frac12\int A(x)\,D(x,y)\,A(y)$ plus the other terms. Integrating out the phonons gives the retarded
interaction of the electrons:
\excode{integrate_out_phonons}
\exformula{integrate_out_phonons}

% ------------------------------------------------------------------ 13
\section{Connected functions as trees of vertex functions}\label{sec:trees}

\begin{function}{ConnectedFunctionFromVertices}
\syntax{ConnectedFunctionFromVertices[\{x1, \dots, xn\}, "Boson"]}
\syntax{ConnectedFunctionFromVertices[\{x1, \dots, xn\}, "Fermion"]}
\syntax{ConnectedFunctionFromVertices[\{x1, \dots, xn\}, "FermionLines"]}
\syntax{ConnectedFunctionFromVertices[\{x1, \dots, xn\}, "FermionLines", \{"Out", "In", \dots\}]}
\syntax{ConnectedFunctionFromVertices[n, type]}
\end{function}
The connected $n$-point function at zero source as a sum of trees of full propagators \pkg{FullPropagator} and
one-particle-irreducible vertices \pkg{VertexFunction} (the integer form uses the labels \pkg{x1, \dots, xn}). The trees
are generated by differentiating the full propagator with the rules that follow from the Legendre transform (the tree
expansion~\cite{JonaLasinio1964}, \cite[sections~7.5 and 7.8]{ZinnJustin2002}), declared with
\fnref{DeclareFunctionalDerivative} for the source \pkg{ExternalSource}: for bosons
\[
\frac{\delta\bar\phi_a}{\delta J_b}=G_{ab},\qquad \frac{\delta G_{ab}}{\delta J_c}=-G_{ax}G_{by}G_{cz}\Gamma^{(3)}_{xyz},\qquad
\frac{\delta\Gamma^{(n)}_{a\cdots}}{\delta J_b}=\Gamma^{(n+1)}_{a\cdots x}G_{xb},
\]
with $\bar\phi=\langle\phi\rangle$ the mean field (\pkg{MeanField}); for fermions the superfield rules, with the Grassmann grading of
\fnref{DeclareGrassmannGraded} and antisymmetric propagators and vertices. The numbers of bosonic trees are 4, 26 and 236
for $n=4,5,6$. The function resets the declarations of these four objects each time it is called, so that the bosonic
and fermionic rules do not mix.
\excode{tree_four_point}
\exboth{tree_four_point}
Every tree is generated once, with the weight $(-1)^v$ for $v$ vertices. Integrating over the external labels identifies
the trees that differ only by the labels of their legs, and the canonical form counts them. The weights and the numbers
of vertices of the three shapes of the connected function $W^{(5)}$ ($W=\log Z$; $W$ is the screened interaction only in
Hedin's equations, section~\ref{sec:exact}): one tree with $\Gamma^{(5)}$, ten with $\Gamma^{(4)}\Gamma^{(3)}$ and fifteen with
three $\Gamma^{(3)}$,
\excode{tree_shapes}
\extext{tree_shapes}
For fermions only the trees with even vertices survive at zero source, with the sign of the permutation of the legs:
\excode{tree_fermion}
\exboth{tree_fermion}
With \pkg{"FermionLines"} the superfield trees are projected on the components $\psi$, $\bar\psi$: the result is
$\langle\psi(x_1)\bar\psi(x_2)\psi(x_3)\bar\psi(x_4)\cdots\rangle_c$ (or, with a third argument, $\psi$ for \pkg{"Out"} and
$\bar\psi$ for \pkg{"In"} at each label), a sum of trees of directed lines \pkg{FullPropagator[x, y]} $=G(x,y)=-\langle
\psi(x)\bar\psi(y)\rangle_c$ and vertices \pkg{VertexFunction[o1, \dots, on, i1, \dots, in]}, with $n$ outgoing and $n$
incoming legs, antisymmetric within each group. $W^{(6)}$ has ten trees, $\Gamma^{(6)}$ and nine with two $\Gamma^{(4)}$:
\excode{tree_fermion_lines}
\expictures{tree_fermion_lines}

\begin{function}{FullPropagator}
\syntax{FullPropagator[x, y]}
\end{function}
The full propagator (connected two-point function), displayed $G_{xy}$.

\begin{function}{VertexFunction}
\syntax{VertexFunction[x1, \dots, xn]}
\end{function}
The one-particle-irreducible vertex $\Gamma^{(n)}$, displayed with its labels; drawn as a grey blob.

\begin{function}{MeanField}
\syntax{MeanField[x]}
\end{function}
The mean field $\bar\phi_x=\langle\phi_x\rangle$ at the current source (the one-point function, the expectation value of the field), displayed $\bar\phi_x$.

\begin{function}{ExternalSource}
\syntax{ExternalSource[x]}
\end{function}
The source $J$ with respect to which \fnref{ConnectedFunctionFromVertices} differentiates.

% ------------------------------------------------------------------ 14
\section{Pictures}\label{sec:pictures}

\subsection{Functions}

\begin{function}{DrawDiagram}
\syntax{DrawDiagram[term, \{x1, \dots, xn\}, options]}
\end{function}
Draws one term as a Feynman diagram (Mathematica \pkg{Graphics}); with several terms it is \fnref{DiagramPlot}.

\begin{function}{DiagramPlot}
\syntax{DiagramPlot[expr, \{x1, \dots, xn\}, options]}
\end{function}
Draws every canonical term of \pkg{expr} with its weight, in a grid (\pkg{Graphics}), or as a \pkg{Grid} of weights and
pictures with \pkg{"Output" -> "Grid"}. The terms are in canonical order, those with fewer integration variables first.
The weights are those of section~\ref{sec:weights}.

In a notebook the pictures follow its appearance. With Mathematica 14.3 or later, through \pkg{LightDarkSwitched}, the
lines and the text are dark on a light notebook and light on a dark one; what hides the lines under it (the inside of a blob, of an
external point or of a source insertion, the clearing behind the label of a hatched or dotted blob) has the colour of the
background of the notebook, and the grey fills are a darker grey there. Before 14.3 the pictures have the colours of a
light notebook. \pkg{"Appearance" -> "Light"} or \pkg{"Dark"} forces one of them, on a background of its own (white, or
the dark grey of a dark notebook), for a picture exported for a page of the other kind;
\pkg{SetOptions[\{DiagramPlot, DrawDiagram\}, "Appearance" -> "Dark"]} does it for a session. The TikZ code of
\fnref{DiagramToTikZ}, for white pages, ignores the option.

\begin{function}{DiagramToTikZ}
\syntax{DiagramToTikZ[expr, \{x1, \dots, xn\}, options]}
\end{function}
The same pictures as \LaTeX/TikZ code: a string with one \pkg{tikzpicture} per term, each preceded by its weight and
the sign, to be used in math mode; with \pkg{"Output" -> "List"}, the list of the pairs \{weight, picture\} (strings). A term
without lines (a number) has no picture. The code needs
\begin{code}
\usepackage{tikz}
\usetikzlibrary{decorations.pathmorphing}
\end{code}
(the library draws the wavy and zigzag straight lines; arcs of these styles are written as curves through computed
points, and arrow heads as small filled polygons, so that the pictures compile at any scale). The pictures are aligned
on the math axis: the external line of a two-point function lies on it, so that sums of diagrams line up with their
signs. The labels are anchored on the side facing their point, away from the lines that leave it. The sizes of the arrow
heads, dashes, waves and double lines follow \opt{LineWidth} and \opt{Scale}: the waves are 1.6\,mm long at the default
scale and at least 1.3\,mm in small pictures.
\excode{tikz_fock}
\extext{tikz_fock}

\subsection{The weights}\label{sec:weights}

The number written in front of a picture, times the values of its elements, is the term. A propagator stands for itself
($G(x,y)$ for an arrow from $y$ to $x$); an interaction line (\fnref{DeclareInteractionLine}, or the option
\opt{Interactions}) stands for minus its kernel, $-V$; a local vertex (\fnref{DeclareLocalVertex}) for minus its kernel; a
vertex for minus its coupling (\fnref{DeclareCoupling}): the minus signs come from expanding $e^{-S_{\rm int}}$. The
internal points are integrated over. So the Fock self-energy $-G(x,y)V(x,y)$ has the weight $1$, the Hartree
self-energy $\delta(x,y)\int_zV(x,z)G(z,z)$ the weight $-1$ (the sign of its closed fermion loop), and the tadpole of $\phi^4$
theory, $-\frac\lambda2\int_zD_0(a,z)D_0(z,z)D_0(z,b)$, the weight $\frac12$. Only the weights written with the pictures
follow this convention: the expressions, and \fnref{DiagramTerms}, keep their coefficients, and \pkg{"Weights" ->
"Coefficients"} writes the coefficients themselves.

An anomalous line ($-\langle T\psi\psi\rangle$, $-\langle T\bar\psi\bar\psi\rangle$, styles \opt{DivergingArrows} and
\opt{ConvergingArrows}) and an antisymmetric blob of two labels (a gap function) show no orientation: the weight takes them
read along their fermion line, around a loop in either sense (a loop holds an even number of them, so the sense does not
matter) and along an open line from its end at the first external label (the sense matters only when both ends are of one
kind). So $V(x_1,x_2)\,\mathcal G_{\psi\psi}(x_1,x_2)$, the gap function at first order, has the weight $-1$ with the labels
\pkg{\{x1, x2\}} and $+1$ with \pkg{\{x2, x1\}}, and the pairing term $-\frac12\int_{z_1z_2}V(z_1,z_2)\,\mathcal G_{\psi\psi}(z_1,z_2)\,\mathcal
G_{\bar\psi\bar\psi}(z_1,z_2)$ of the Luttinger--Ward functional, a loop read as $\mathcal G_{\psi\psi}(z_1,z_2)\,\mathcal G_{\bar\psi\bar\psi}(z_2,z_1)$, the
weight $-\frac12$. The undirected lines and the vertex of a superfield are antisymmetric too, but on no fermion line: their
pictures fix a term up to the order of its labels, and the weight refers to the canonical labelling.

\subsection{What is drawn}

\begin{itemize}
\item Two-label factors are lines, drawn in the style of their head (\fnref{DeclareLineStyle}); a directed line
  \pkg{G[x, y]} runs from \pkg{y} to \pkg{x}. Factors with three or more labels are grey blobs (vertex functions), with
  the lines that end on them drawn to the blob. Factors with one label are source insertions, a crossed circle.
\item Labels joined by \pkg{DeltaFunction} are one point.
\item An external label alone at the end of a single line (a propagator or an interaction line) is an external point of a
  correlation function, drawn as an open circle. Any other external label sits on a vertex of an amputated function
  (self-energy, vertex) and gets an external leg. For fermions the leg has an arrow: outgoing if one more fermion line
  enters the vertex than leaves it, incoming in the opposite case, none if as many enter as leave. The directions are
  decided from all the terms of the expression together, so that the legs of a sum are consistent; \opt{LegDirections}
  sets them explicitly.
\item An external label at which a density is inserted (two fields at one point, such as $\bar\psi\psi$ or $\phi^2$: in
  every term of the expression the label is alone at its point, which ends exactly two propagators, one entering and one
  leaving or two undirected ones, and nothing else, and it is not an argument of a local vertex) is the insertion of an
  external field $U(z)$: it is drawn as the crossed circle of a source insertion, without a leg.
  \opt{AmputatedLegs}~\pkg{->}~\pkg{True} or \pkg{False} turns this off.
\item Disconnected parts are drawn side by side: vacuum bubbles next to the rest, and a product of two one-point functions
  (such as $\langle\phi(x)\rangle\langle\phi(y)\rangle$) as two pictures.
\end{itemize}

\subsection{Options}\label{sec:options}

\fnref{DiagramPlot}, \fnref{DrawDiagram} and \fnref{DiagramToTikZ} take the options below; \pkg{SetOptions} changes their
defaults, for instance \pkg{SetOptions[DiagramToTikZ, "Scale" -> 0.5]}, and an unknown option name gives a message.

\begin{center}
\small
\begin{tabularx}{\textwidth}{@{}>{\ttfamily\small}l>{\ttfamily\small}l>{\raggedright\arraybackslash}X@{}}
\toprule
\textrm{option} & \textrm{default} & \\
\midrule
"VertexLabels" & Automatic & the labels of the external points; \texttt{All}: also $z_1,z_2,\dots$; \texttt{None}: no labels\\
"Labels" & \{\} & TikZ: rules \pkg{label -> "LaTeX"}; by default \pkg{x12} is written $x_{12}$, \pkg{lambda} $\lambda$, \pkg{etabar}
  $\bar\eta$\\
"ExternalOrder" & Automatic & the order of the external labels, counterclockwise from the top left (for four points: top
  left, bottom left, bottom right, top right): those given to the function, or a list \pkg{\{x1, \dots\}};
  \texttt{"Uncrossed"}, \pkg{\{"Uncrossed", \{x1, \dots\}\}}: in each four-point diagram the legs of one side exchanged
  where that keeps its lines from crossing\\
"LegDirections" & Automatic & arrows on the external legs: inferred from the fermion lines; \texttt{"Alternating"}: out,
  in, out, in, \dots; a list of \texttt{"Out"}, \texttt{"In"}, \texttt{None}; \texttt{None}\\
"AmputatedLegs" & Automatic & \texttt{True}: every external label gets a leg; \texttt{False}: none, all are points\\
"InteractionStyle" & Automatic & \texttt{"Wavy"}, \texttt{"DoubleWavy"}, \texttt{"Curly"}, \texttt{"Dashed"},
  \texttt{"Dotted"}, \texttt{"Zigzag"}: the style of all the interaction lines\\
"Directed", "Interactions" & \{\} & heads drawn as arrows, or as wavy interaction lines, in this picture only (e.g.\ a
  screened $W$ declared \texttt{"DoubleWavy"}, drawn as an ordinary wavy line in a skeleton expansion)\\
"Couplings" & \{\} & couplings absorbed into the vertices in this picture only, as with \fnref{DeclareCoupling}\\
"LineWidth" & Automatic & the width of the lines in pt (TikZ 0.7, Mathematica 1.3); arrow heads, dashes and double lines
  follow it\\
"Scale" & 0.7 & TikZ: centimetres per unit of the drawing\\
"LabelFont" & "\textbackslash\textbackslash small" & TikZ: the font command of the labels\\
"Output" & "Graphics" & \texttt{"Grid"} (\pkg{DiagramPlot}); \texttt{"List"} (\pkg{DiagramToTikZ})\\
"Columns" & Automatic & \pkg{DiagramPlot}: diagrams per row\\
"Weights" & True & \pkg{DiagramPlot}: show the weights (section~\ref{sec:weights}); \texttt{"Coefficients"}: the
  coefficients of the terms, every line and vertex standing for its kernel (also \pkg{DiagramToTikZ})\\
ImageSize & Automatic & Mathematica: the width of the picture in pixels (\pkg{Automatic}: 46 pixels per unit)\\
"Appearance" & Automatic & Mathematica: the colours of a light or a dark notebook, as the notebook is (14.3 or later);
  \texttt{"Light"}, \texttt{"Dark"}: forced, on a background (section~\ref{sec:pictures}); ignored by \pkg{DiagramToTikZ}\\
\bottomrule
\end{tabularx}
\end{center}
The exchange diagram of the Green function at first order with \pkg{"VertexLabels" -> All}, and with
\pkg{"InteractionStyle" -> "Dashed"}:
\excode{vertex_labels}
\expictures{vertex_labels}
\excode{interaction_style}
\expictures{interaction_style}
The four labels of a four-point function are at the corners counterclockwise from the top left, in the order given, in
every diagram; a diagram whose channel joins two opposite corners, as the bubble of the channel $t$ $(ac|bd)$ of
$\Gamma^{(4)}$ in $\phi^4$ theory (the example of \fnref{PerturbativeFullVertex}), is drawn with two of its legs crossing as
an X just outside it. With \pkg{"ExternalOrder" -> "Uncrossed"} the legs of one side of such a diagram are exchanged
instead, the labels moving with them: of the order given, the order with the two labels on the right exchanged and the
order with the two labels on the top exchanged (the three ways of going round four corners), each diagram takes the first
in which its lines, taken straight, cross least:
\excode{external_order}
\expictures{external_order}
The value \pkg{\{"Uncrossed", \{x4, x2, x3, x1\}\}} starts from another order: here the particle--particle ladder of the
bare vertex, the incoming legs $x_4$, $x_2$ on the left, with $x_1$ and $x_3$ exchanged where the particle entering at
$x_2$ leaves at $x_1$:
\excode{external_order_list}
\expictures{external_order_list}

\subsection{Line styles}

\begin{center}
\small
\begin{tabularx}{\textwidth}{@{}>{\ttfamily\small}lX@{}}
\toprule
\textrm{style} & \\
\midrule
"Plain" & a plain line (bosonic propagators, the default)\\
"Arrow" & a propagator with an arrow from the second label to the first\\
"DoubleArrow" & the same, double (a dressed propagator)\\
"DottedArrow" & the same, dotted (a ghost)\\
"Double" & a double line\\
"Wavy" & an interaction line\\
"DoubleWavy" & a double wavy line (a dressed interaction beside the bare one: $W$ next to $V$)\\
"Curly" & a coil (a gluon)\\
"Dashed" & a dashed interaction line (a phonon)\\
"Dotted" & a dotted interaction line\\
"Zigzag" & a zigzag interaction line\\
"Impurity" & a dashed interaction line with a cross (a disorder average)\\
"DivergingArrows" & an anomalous propagator $-\langle T\psi\psi\rangle$, two arrowheads pointing to the ends\\
"ConvergingArrows" & an anomalous propagator $-\langle T\bar\psi\bar\psi\rangle$, two arrowheads pointing to the middle\\
\bottomrule
\end{tabularx}
\end{center}

\subsection{Layout}\label{sec:layout}

No coordinate is stored: the layout is computed from the diagram, deterministically, so that the same diagram is always
drawn the same way.

\paragraph{Fermion lines} In a theory of fermions coupled by density-type vertices --- two-body interactions, phonons,
impurities --- every vertex lies on exactly one fermion line: an open line between external points, or a closed loop (a
vertex with a tadpole $G(z,z)$ is a loop by itself). The fermion lines are found by following the directed propagators,
and the diagram is drawn line by line:
\begin{itemize}
\item a two-point function whose external points are the ends of one open line has it on a horizontal axis, the first
  external label on the left; the lines between vertices of the axis are arcs, above or below;
\item a four-point function whose external points are the ends of two open lines has them as two parallel rails. The
  external labels are at the corners of a rectangle, in the order of \opt{ExternalOrder}; each rail joins two neighbouring
  corners, so that the rails are horizontal or vertical according to the pairing of the external points (if a rail would
  join opposite corners, the rails are laid out with the labels of the two right corners exchanged, and the two legs on
  the right cross back as an X; with \texttt{"Uncrossed"} they stay exchanged). The legs point to the corners;
\item the interaction lines between the rails are perpendicular to them, so that ladders are drawn as ladders; two lines
  that cross make a symmetric cross, with more room around it. Arcs between vertices of one rail go outside, or inside
  when that avoids a conflict;
\item every closed loop is a circle with its vertices evenly spaced, or a bubble of two arcs when it has two vertices.
  A bubble joined to two points of a line that are further apart than its width (a dressed interaction line) keeps its
  width and hangs below them, and its two interaction lines are bent into one round arc with it.
  The loops form blocks hanging from the open lines. A loop joined to one vertex of a line hangs from it, pointing away;
  a loop joined to several vertices of a line sits above them, its joined vertices on the lower half of the circle right
  above their points, so that the lines joining them are perpendicular to the line; loops hanging from loops point away
  from their centre; a loop joined to both rails lies between them; chains of loops between two points are straight;
\item the side of every arc and block is chosen by trying all the assignments and counting conflicts: two arcs that cross,
  an arc over a vertex with something on its side, blocks with interleaved attachments, and (less) an arc nested under a
  block, a block at the end of an arc. The positions along the lines follow from the room that what hangs from them needs,
  and they are widened until no two loops overlap;
\item vacuum diagrams made of three or more loops joined as a tree have the most central loop in the middle, turned so
  that as many lines as possible leave it horizontally or vertically, and the others around it.
\end{itemize}

\paragraph{Trees of vertex functions} A connected function written as a tree of vertex functions joined by propagators,
with the external points at its leaves (the tree expansion, \fnref{ConnectedFunctionFromVertices}), has its central vertex,
or its two central vertices, in the middle. Chains of vertices are straight; the legs of a vertex at the end of a chain
are fanned outwards, those of a vertex along a chain perpendicular to it; a vertex with three or more vertex neighbours
has them all around. The external points follow each other counterclockwise in the order of their labels; a tree of
four points and two vertices whose pairs are the first and the last label is drawn upright, and one that pairs opposite
corners is drawn upright with the two legs on the right crossing, at 60 degrees, clear of the vertices. A tree whose
leaves are free ends (its labels integrated over: the shapes of the trees $W^{(n)}$ under \fnref{ConnectedFunctionFromVertices}) is drawn
alike, the free ends where the external points would be and its centre on the axis; its inner nodes may also be bare
vertices of three or more propagators. Its leaves have no order: the branches of a vertex are arranged by their shapes,
so that a symmetric tree is drawn symmetric about the vertical axis, and the legs along a chain are on its upper side.

\paragraph{Other diagrams} Diagrams of other kinds --- bosonic theories, vertices with more than two fermion lines,
propagators alone --- are drawn by the general layout: the external points of a two-point function on
a horizontal axis with the rest in layers above it; those of an $n$-point function on a circle (at the corners of a
square for four points), counterclockwise from the top left, the open lines between neighbouring external points
straight, and the rest placed by a barycentric embedding (a vertex with two legs at opposite corners, which would be in
the middle with its legs across the picture, is placed with the labels of the right corners exchanged, and the two legs
on the right cross back as an X outside the diagram, at 60 degrees when both leave bubbles, which they flatten into
lenses; a loop through the vertex of such a leg is then drawn with straight lines if its circle would be crossed); vacuum parts as a cycle with the rest inside, or several loops
side by side, around a circle, or one inside the other (two loops joined by three or more lines). Closed fermion loops are
circles in this layout too. When three interaction lines would be diameters of a fermion loop drawn as a circle, meeting at
its centre, one is drawn outside the loop, over the top, and the other two cross at the centre; with four or more the vertices
are moved alternately apart, so that the lines cross in pairs. Two bubbles side by side are joined at the sides facing each other, and where two wavy lines
cross their waves are flattened in the TikZ output.
Closed loops of undirected propagators with three or more vertices are drawn as circles too, and two bubbles that meet
at a vertex as flat lenses. Of two chords of a loop that would cross, one is drawn outside it. A picture without external
points is turned upright along its principal axes.

\paragraph{Symmetry} Among layouts of equal quality the symmetric one is taken, so that symmetric diagrams are drawn
symmetric.

% ------------------------------------------------------------------ 15
\section{Formulas in \LaTeX}\label{sec:latex}

\begin{function}{ExpressionToLaTeX}
\syntax{ExpressionToLaTeX[expr, options]}
\end{function}
The \LaTeX{} code of a formula, to be used in math mode: integrals with their variables as a subscript,
$\int_{z_1\,z_2}$; functions with their arguments in parentheses; \pkg{DeltaFunction} as $\delta(x,y)$; Grassmann products
in their order, after the commuting factors; rational coefficients as fractions, factors with negative powers in a
denominator; derivatives of functions of one argument with primes; \pkg{E\^{}x} as $e^{x}$, square roots as
$\sqrt{\phantom{x}}$, \pkg{Log}, \pkg{Sin}, \pkg{Cos}, \dots{} as $\log$, $\sin$, $\cos$, \dots; the canonical integration
variables as $z_k$, $w_k$, \dots. Options:
\begin{itemize}
\item \opt{Names} \pkg{-> \{h -> "G\_0", \dots\}}: the \LaTeX{} of heads and symbols. By default a head is written with the
  name given by \fnref{DeclareDisplayName}, else with its symbol name: greek names become letters (\pkg{lambda},
  $\lambda$), trailing digits a subscript (\pkg{G0}, $G_{0}$), a final \pkg{bar} a bar (\pkg{etabar}, $\bar\eta$), and other
  names of several letters are set upright. Subexpressions can be named too, for instance
  \pkg{FermiFunction[xi[p1]] -> "n\_1"} (a factor, or a term of a sum);
\item \opt{Labels} \pkg{-> \{x -> "x\_1", \dots\}}: the \LaTeX{} of labels, as in \fnref{DiagramToTikZ};
\item \opt{Subscripts} \pkg{-> \{xi, v, \dots\}}: heads whose arguments are written as a subscript, \pkg{xi[p]} as $\xi_{p}$
  instead of $\xi(p)$, for dispersions and interactions in momentum space (section~\ref{sec:matsubara}).
\end{itemize}
The statistical factors, delta functions and momentum sums of section~\ref{sec:matsubara} are written $n_F(\epsilon)$,
$n_B^{\prime}(\epsilon)$, $\delta(\omega-\epsilon)$, $\sum_{p_1,p_2}$, and \pkg{UnitStep} as $\theta$. The objects of
section~\ref{sec:dimreg}: a loop integral as $\int_{k\,l}$ followed by one fraction with the lines in its denominator,
$1/(k^2+m^2)$, a massive bubble as $B_{a,b}(p^2;m_1,m_2)$, \pkg{RenormalizationScale} as $\mu$, \pkg{EulerGamma} as
$\gamma_E$, \pkg{Zeta} as $\zeta$, \pkg{ArcTanh} as $\operatorname{artanh}$; a power of a power is written with
parentheses, $(p^2)^{1-2\varepsilon}$. The objects of section~\ref{sec:qed}: $\gamma^\mu$, $\slashed p$ (\verb|\slashed|, of the
\LaTeX{} package \pkg{slashed}), the unit matrix $\mathbf 1$, $\operatorname{tr}[\dots]$, $\bar u(p_2)\dots u(p_1)$, $g^{\mu\nu}$, $p^\mu$, $p\cdot
q$, \pkg{SpacetimeDimension} as $d$, the Feynman parameters as $x_i$; the denominators of the lines as one fraction, and
the integrals over the loop momentum and the Feynman parameters as $\int_k$, $\int_0^1dx_1$ and
$\int_0^1dx_1\,dx_2\,dx_3\,\delta(1-x_1-x_2-x_3)$. All the Lorentz indices are written upper.
\excode{latex_names}
\extext{latex_names}

% ------------------------------------------------------------------ 16
\section{Evaluation: Matsubara sums}\label{sec:matsubara}

The diagrams of the package are functions of labels. For a system that is translation invariant in space and in imaginary
time, at a temperature $T$, they become the familiar expressions in momentum and Matsubara
frequency~\cite{Matsubara1955,Mahan,BruusFlensberg}: every independent loop of a diagram (not only the closed fermion
loops) carries a momentum $p$ and a frequency, and contributes $T\sum_n\int\!\frac{d^dp}{(2\pi)^d}$.
\fnref{EvaluateMatsubara} routes the momenta and frequencies through the lines and carries out the frequency sums by
residues; the momentum integrals are left as \pkg{MomentumSum}. The result
is continued to real frequency by \fnref{AnalyticContinuation}, and \fnref{RetardedImaginaryPart} gives its imaginary part,
which the finite-temperature cutting rules (\fnref{CuttingRuleImaginaryPart}, \fnref{CheckCuttingRules}) compute
independently.

\subsection{Conventions}\label{sec:matsubara-conventions}

\begin{itemize}
\item \textbf{Propagators.} $G_0(k,i\omega_n)=1/(i\omega_n-\xi_k)$ and $G(\tau)=T\sum_ne^{-i\omega_n\tau}G(i\omega_n)$, consistent with
  $G_0(x,y)=-\langle T\psi(x)\bar\psi(y)\rangle_0$; a line \pkg{G[x, y]} carries its momentum and frequency from $y$ to $x$.
  The fermionic frequencies are $\omega_n=(2n+1)\pi T$, the bosonic ones $\nu_n=2n\pi T$; a symbol \pkg{w} declared by
  \fnref{DeclareMatsubaraFrequency} stands for one of them and is written \pkg{I w} in expressions.
\item \textbf{Lines.} The rules of \fnref{EvaluateMatsubara} say what each head is: \pkg{G0 -> \{"Fermion", xi\}} for
  $1/(i\omega-\xi(p))$; \pkg{V -> \{"Instantaneous", v\}} for an interaction $V(x,y)=v(\mathbf r_x-\mathbf r_y)\,\delta(\tau_x-\tau_y)$,
  which carries no frequency and becomes $v(q)$; \pkg{D0 -> \{"Boson", Function[\{q, z\}, d]\}} for a boson propagator, a
  rational function $d(q,z)$ of $z=i\nu_n$. Instantaneous and boson lines are even, $v(-q)=v(q)$ and $d(-q,-z)=d(q,z)$.
\item \textbf{Signs.} The package's correlation functions are time-ordered products without the minus sign of the
  response functions. The density correlation $\langle\rho_a\rho_b\rangle_c=-G_0(a,b)\,G_0(b,a)$ is, to lowest order, minus
  the Lindhard function~\cite{Lindhard1954}
  $\chi_0(q,i\nu_n)=(2s+1)\sum_p\bigl(n_F(\xi_p)-n_F(\xi_{p+q})\bigr)/(i\nu_n+\xi_p-\xi_{p+q})$, and the boson propagator
  $D_0=\langle\phi\phi\rangle_0$ of $Z_0=e^{\frac12J\cdot D_0\cdot J}$ is minus the propagator
  $\mathcal D_0=-\langle T\phi\phi\rangle_0$ of the textbooks: for phonons~\cite{Mahan}
  $\mathcal D_0(q,i\nu_n)=2\Omega_q/((i\nu_n)^2-\Omega_q^2)$, so that the line of the package, $D_0=-\mathcal D_0$, is
  \pkg{Function[\{q, z\}, 2 Om[q]/(Om[q]\^{}2 - z\^{}2)]}.
\item \textbf{Convergence factors.} A sum $T\sum_n g(i\omega_n)$ of a summand that decays only like $1/\omega_n$ depends on
  how it is cut off. A propagator at equal labels, $G(z,z)$, and the propagators between labels joined by instantaneous
  interactions are taken at equal times, $G(\tau=0^-)$, with the factor $e^{i\omega_n0^+}$:
  $T\sum_ne^{i\omega_n0^+}/(i\omega_n-\xi)=n_F(\xi)$, the density. \fnref{EvaluateMatsubara} finds these lines from the
  topology of the diagram; \fnref{MatsubaraSum} takes the factor as an option. The factor is the ordering of the
  operators at equal times, and the option \opt{ConvergenceFactor} of \fnref{EvaluateMatsubara} chooses one of three:
  \pkg{1} (the default), $G(0^-)=\langle\psi^\dagger\psi\rangle$, the factor $e^{i\omega_n0^+}$ and $n_F(\xi)$; \pkg{-1},
  $G(0^+)=-\langle\psi\psi^\dagger\rangle$, the factor $e^{-i\omega_n0^+}$ and $n_F(\xi)-1$; \pkg{0}, the mean
  $\frac12\langle[\psi^\dagger,\psi]\rangle$, the symmetric sum without a factor and $n_F(\xi)-\frac12$; a boson line,
  even and continuous at $\tau=0$, needs no factor (\fnref{MatsubaraSum} of a boson pole gives $-n_B$, $-n_B-1$ and
  $-n_B-\frac12$). They differ by a constant per state; the symmetric choice is natural on a lattice with a
  particle--hole symmetric interaction, where the difference is a shift of the chemical potential, while in the continuum
  the sum over states of the constant diverges.
\item \textbf{Statistical factors.} $n_F(\epsilon)=1/(e^{\epsilon/T}+1)$ and $n_B(\epsilon)=1/(e^{\epsilon/T}-1)$ are inert
  (\pkg{FermiFunction}, \pkg{BoseFunction}); a pole of order $k$ gives their derivatives of order $k-1$, as for a self-energy
  insertion on a line. They are periodic in the Matsubara frequencies: $n_F(\epsilon+i\nu_m)=n_F(\epsilon)$,
  $n_F(\epsilon+i\omega_m)=-n_B(\epsilon)$, $n_B(\epsilon+i\omega_m)=-n_F(\epsilon)$; after each sum the declared
  frequencies are removed from the arguments in this way, and every argument is written with its first non-constant term
  positive ($n_F(-\epsilon)=1-n_F(\epsilon)$, $n_B(-\epsilon)=-1-n_B(\epsilon)$); constants such as the chemical
  potential do not count, so that $n_F(\epsilon_p-\mu)$ stays as it is. $n_F(0)=\frac12$ at every temperature.
\item \textbf{Analytic continuation.} Only after the sums, when the statistical factors no longer depend on the external
  frequency, can $i\omega_n$ be continued to $\omega+i0^+$, which gives the retarded function~\cite{FetterWalecka}.
\item \textbf{Form of the results.} Terms with the same denominator are put together; a function of one external
  frequency is written in partial fractions of it (part by part, see \fnref{EvaluateMatsubara}); and the statistical
  factors are written as a sum over processes, every line emitted ($1-n_F$, $1+n_B$) or absorbed ($n_F$, $n_B$), when
  this is shorter, as in the second-order self-energy below.
\item \textbf{Lines.} Fermion lines are directed, $1/(i\omega_n-\xi_k)$, and boson lines even: an antisymmetric kernel
  (an anomalous propagator $-\langle T\psi\psi\rangle$ or a superfield propagator) is refused with the message
  \pkg{EvaluateMatsubara::anomalous}, and an expression of a superfield is evaluated through its components at zero
  anomalous propagators (\fnref{SuperfieldComponents}).
\end{itemize}

\subsection{Frequency sums}

\begin{function}{MatsubaraSum}
\syntax{MatsubaraSum[f, z, "Fermion"]}
\syntax{MatsubaraSum[f, z, "Boson", "ConvergenceFactor" -> eta]}
\end{function}
$T\sum_nf(z=i\omega_n)$ over the fermionic or the bosonic Matsubara frequencies, for $f$ rational in $z$ with a
denominator that factors into linear and quadratic factors (propagators), by residues: the sum over the poles $z_j$ of $f$ of
$\operatorname{Res}[f\,n_F](z_j)$ (fermions) or $-\operatorname{Res}[f\,n_B](z_j)$ (bosons). Denominators of degree two in
$z$ (boson propagators) are split into linear factors; residues are Laurent coefficients, computed without
\pkg{Residue}, \pkg{Limit} or \pkg{Apart}. How $f$ decays at
large $z$ is read from the degrees in $z$ of its terms. The option \opt{ConvergenceFactor} matters when $f$ decays like
$1/z$: \pkg{1} for the factor $e^{z0^+}=e^{i\omega_n0^+}$ (as for a density), \pkg{-1} for $e^{-z0^+}=e^{-i\omega_n0^+}$,
\pkg{0} for the symmetric sum; for $f=1/(z-a)$ they give $n_F(a)$, $n_F(a)-1$ and $n_F(a)-\frac12$ (bosons: $-n_B(a)$,
$-n_B(a)-1$ and $-n_B(a)-\frac12$). With the default \pkg{Automatic} the symmetric sum is taken, with the message
\pkg{MatsubaraSum::conv}. A summand that does not decay (the sum diverges: \pkg{MatsubaraSum::diverge}), or whose
denominator is not split into poles, gives a message and \pkg{\$Failed}.
\excode{matsubara_sum}
\exformula{matsubara_sum}
\excode{matsubara_two_poles}
\exformula{matsubara_two_poles}
A double pole gives the derivative of the Fermi function:
\excode{matsubara_double_pole}
\exformula{matsubara_double_pole}
A pole on the axis of the frequencies summed (an anomalous term, as for $1/z$ in a bosonic sum) gives
\pkg{MatsubaraSum::axis}: the energies are taken to be generic. A pole at a frequency of the other statistics is an
ordinary one: the fermionic sum of $1/(z(z-a))$ is $(n_F(a)-n_F(0))/a$, with $n_F(0)=\frac12$.

\paragraph{Logarithms} A summand may contain logarithms $c\ln r(z)$, with $c$ free of $z$ and $r$ rational in $z$: the sums
of the traces ${\rm Tr}\ln$ of section~\ref{sec:effective}. Each zero $a$ of $r$ counts $+1$ and each pole $-1$, with
\[
T\sum_ne^{i\eta\omega_n0^+}\ln(a-i\omega_n)=-T\ln\bigl(1-n_F(a)\bigr)+\frac{1-\eta}{2}\,a=T\ln\bigl(1+e^{-a/T}\bigr)
+\frac{1-\eta}{2}\,a
\]
for fermions and $-T\ln(1+n_B(a))+\frac{1-\eta}{2}a=T\ln(1-e^{-a/T})+\frac{1-\eta}{2}a$ for bosons, $\eta=1,-1,0$ the
three convergence factors (the derivative in $a$ is minus the sum of $1/(z-a)$, and the constant is fixed by an empty level,
$a\to+\infty$). A constant factor gives nothing: the trace of the identity is zero. The factor of a real boson,
$\ln(W^2-z^2)$, gives $W+2T\ln(1-e^{-W/T})$ with every ordering, the zero-point energy included. The temperature appears
explicitly, as \fnref{Temperature}:
\excode{matsubara_log}
\exformula{matsubara_log}

\begin{function}{DeclareMatsubaraFrequency}
\syntax{DeclareMatsubaraFrequency[w, "Fermion"]}
\syntax{DeclareMatsubaraFrequency[\{w1, w2, \dots\}, "Boson"]}
\end{function}
Declares symbols as Matsubara frequencies of the given statistics, so that statistical factors at energies shifted by
them are simplified by periodicity and that \fnref{EvaluateMatsubara} knows the statistics of the external frequencies.
\fnref{ResetDeclarations} removes the declarations.
\excode{matsubara_periodicity}
\exformula{matsubara_periodicity}

\begin{function}{FermiFunction}
\syntax{FermiFunction[e]}
\end{function}
The Fermi function $n_F(e)=1/(e^{e/T}+1)$, inert, except that $n_F(0)=\frac12$. Its derivative of order $k$,
\pkg{Derivative[k][FermiFunction][e]}, is written $n_F^{\prime}(e)$, $n_F^{\prime\prime}(e)$, \dots.

\begin{function}{BoseFunction}
\syntax{BoseFunction[e]}
\end{function}
The Bose function $n_B(e)=1/(e^{e/T}-1)$, inert.

\begin{function}{Temperature}
\syntax{Temperature}
\end{function}
The temperature $T$ of the statistical factors, written $T$ in \LaTeX. It appears explicitly only in the sums of
logarithms (\fnref{MatsubaraSum}) and in the traces ${\rm Tr}\ln$ evaluated by \fnref{EvaluateMatsubara}:
$-T\ln(1-n_F(a))=T\ln(1+e^{-a/T})$ is minus the grand potential of one level.

\subsection{Diagrams}

\begin{function}{EvaluateMatsubara}
\syntax{EvaluateMatsubara[diagrams, \{x1 -> \{k1, I w1\}, x2 -> \{k2, I w2\}, \dots\}, rules, options]}
\end{function}
Evaluates the diagrams of an expression in momentum and Matsubara frequency. \pkg{ki} and \pkg{I wi} enter the diagram at
the external label \pkg{xi}; they must add up to zero, and the statistics of \pkg{wi} is declared by
\fnref{DeclareMatsubaraFrequency}. For a two-point function $F(x,y)$ the momentum $k$ enters at $y$ and leaves at $x$,
\pkg{\{y -> \{k, I w\}, x -> \{-k, -I w\}\}}. The rules are those of section~\ref{sec:matsubara-conventions}; the factors of
a term must be lines of the rules, delta functions or constants. Labels joined by delta functions are one point. A
spanning tree of the diagram is chosen with the instantaneous and boson lines first, so that the loop momenta and
frequencies run along fermion lines; every other line closes a loop, and gets the loop momentum \pkg{LoopMomentum[i]}
and the frequency \pkg{LoopFrequency[i, statistics]}; the lines of the tree carry what conservation at the points gives
them. The frequencies are then summed one after the other by \fnref{MatsubaraSum}, with the convergence factors of the
lines at equal times. A sum that depends on the convergence factor when no line at equal times gives one is taken
symmetric, with the message \pkg{EvaluateMatsubara::conv} (a boson line that tends to a constant at large frequency
should be split, its constant given as an instantaneous line); a summand that does not decay gives
\pkg{EvaluateMatsubara::diverge}. A function of one external frequency is written in partial fractions of it, the
denominators of the textbooks. The lines that carry the external frequency and no loop frequency --- the external legs
of a function that is not amputated, the lines between its irreducible parts --- stay as factors, and the partial
fractions are taken in each part between them, so that no part is evaluated at the poles of another. Options:
\opt{SpinDegeneracy} \pkg{-> 1}, the factor $2s+1$ of every closed fermion loop (for a spin-independent interaction,
spin is not in the labels); \opt{LoopMomenta} \pkg{-> Automatic} (\pkg{LoopMomentum[i]}, written $p_i$) or a list of
names; \opt{ConvergenceFactor} \pkg{-> 1}, the ordering of the fermion lines at equal times
(section~\ref{sec:matsubara-conventions}): \pkg{1} for $G(\tau=0^-)$ and the factor $e^{i\omega_n0^+}$, \pkg{-1} for
$G(0^+)$ and $e^{-i\omega_n0^+}$, \pkg{0} for their mean, the symmetric sum; another value gives
\pkg{EvaluateMatsubara::convergence}. Invalid arguments, or a term that cannot be evaluated, give a message and
\pkg{\$Failed}. The Hartree and Fock self-energies of the electron gas, from the Luttinger--Ward functional of
section~\ref{sec:functionals}:
\excode{evaluate_hartree_fock}
\exformula{evaluate_hartree_fock}
With the symmetric sum every line at equal times has $n_F-\frac12$ in place of $n_F$, and the self-energy changes by a
constant, $-v(0)\sum_p1+\frac12\sum_pv(k-p)$, independent of $k$ on a lattice:
\excode{evaluate_hartree_fock_symmetric}
\exformula{evaluate_hartree_fock_symmetric}
The bubble, which is minus the Lindhard function:
\excode{evaluate_lindhard}
\exformula{evaluate_lindhard}
The second-order skeleton self-energy, with the direct diagram (a closed loop, the factor 2 of spin) and the exchange
diagram; both have the statistical factors of the decay of a particle into two particles and a hole, and of the inverse
process. The output is written with $\xi_1=\xi(p_1)$, $\xi_2=\xi(p_2)$, $\xi_3=\xi(k-p_1+p_2)$ and $n_i=n_F(\xi_i)$, the
interaction as a subscript (options \opt{Names} and \opt{Subscripts} of \fnref{ExpressionToLaTeX}):
\excode{sigma_second_order}
\exformula{sigma_second_order}
A self-energy insertion on a line gives a double pole, and the derivative of the Fermi function. The exchange diagram
with a Fock insertion on its line, a diagram of the perturbative self-energy (weight $+1$), and its value:
\excode{fock_insertion}
\expictures{fock_insertion}
\excode{double_pole_diagram}
\exformula{double_pole_diagram}
A boson line: the self-energy of electrons of spin $\frac12$ (\opt{SpinDegeneracy} \pkg{-> 2}) coupled to phonons, to
order $g^2$~\cite{Migdal1958}, with the boson rule of section~\ref{sec:matsubara-conventions} and the phonon frequency
$\Omega_q$ (\pkg{Om[q]}). The first term is the
exchange of a phonon, with the factors of the textbooks~\cite{Mahan}, $n_B+1-n_F$ for its emission and $n_B+n_F$ for its
absorption; the second is the tadpole, the phonon at zero momentum, $D_0(0,0)=2/\Omega_0$, attached to the density of
the two spin projections, $2\sum_pn_F(\xi_p)$:
\excode{phonon_self_energy}
\exformula{phonon_self_energy}

\paragraph{Traces} With no external point, \pkg{TraceLog[A]} and \pkg{OperatorTrace[A]} (section~\ref{sec:effective}),
\pkg{A} built from kernels that have line rules (\pkg{"Fermion"}, $1/(z-\xi_k)$; \pkg{"Instantaneous"}, $v_k$, e.g.\ a
static self-energy; \pkg{"Boson"}, $f(k,z)$), are evaluated per unit of $\beta$ times the volume, as the vacuum diagrams:
${\rm Tr}\ln A=(2s+1)\sum_kT\sum_ne^{i\eta\omega_n0^+}\ln A(k,i\omega_n)$, and ${\rm Tr}\,A$ the same with $A$, its constant
part dropped (${\rm Tr}\,1=0$). With \fnref{EffectiveAction2PI} this is the grand potential $\Omega/V$, with the ordering of
\opt{ConvergenceFactor}: for the free electron gas, $-{\rm Tr}\ln(-G_0^{-1})$,
\excode{evaluate_grand_potential}
\exformula{evaluate_grand_potential}
\paragraph{Limits} The option \opt{Limit} \pkg{-> "Static"} sets all external momenta and frequencies to zero: the limit
$q\to0$ at zero frequency, the thermodynamic response, in which double poles give derivatives of $n_F$. \pkg{"Dynamic"} sets
the external momenta to zero and takes the external frequency to zero at the end: zero for the density of a conserving
approximation. The terms are grouped by their loop momenta, and the limit is taken in each group; when the singular terms
do not cancel as they are written (at second order they cancel only after a change of loop momenta) the message
\pkg{EvaluateMatsubara::limitterms} is given and the function at zero momentum is returned. The bubble:
\excode{lindhard_limits}
\extext{lindhard_limits}

\begin{function}{DiagramIntegrand}
\syntax{DiagramIntegrand[diagrams, externals, rules, options]}
\end{function}
The integrand of each term before the frequency sums, as \fnref{EvaluateMatsubara} routes it: a list, one entry per term,
of \pkg{\{"Weight" -> \dots, "Integrand" -> \dots, "LoopMomenta" -> \dots, "LoopFrequencies" -> \dots,
"EqualTimeFrequency" -> \dots\}}. The integrand is a function of \pkg{I LoopFrequency[i, statistics]}, and the weight
includes the factor of \opt{SpinDegeneracy}. \opt{EqualTimeFrequency} is the sum of the frequencies of the fermion lines
at equal times, the exponent $\Omega$ of their convergence factor $e^{i\Omega0^+}$ (with the default
\opt{ConvergenceFactor} of \fnref{EvaluateMatsubara}); it is 0 when there is none, as for the bubble. Options
\opt{SpinDegeneracy} and \opt{LoopMomenta}, as in \fnref{EvaluateMatsubara}. Useful for numerical checks, such as the sums over many frequencies
of the test suite:
\excode{diagram_integrand}
\extext{diagram_integrand}

\begin{function}{MomentumSum}
\syntax{MomentumSum[f, \{p1, p2, \dots\}]}
\end{function}
The inert integral $\int\!\frac{d^dp_1}{(2\pi)^d}\cdots f$ (or $\frac1V\sum_p f$) over the loop momenta, written
$\sum_{p_1,p_2}f$. The factors that do not depend on the loop momenta are taken out of it.

\begin{function}{LoopMomentum}
\syntax{LoopMomentum[i]}
\end{function}
The momentum of the $i$-th loop, written $p_i$, and $\ell_i$ in QED (in a \fnref{MinkowskiLoopIntegral} and in the Dirac and Lorentz
objects of an integrand taken out of it), where $p_1, p_2, \dots$ are the external momenta.

\begin{function}{LoopFrequency}
\syntax{LoopFrequency[i, "Fermion"]}
\syntax{LoopFrequency[i, "Boson"]}
\end{function}
The Matsubara frequency of the $i$-th loop in \fnref{DiagramIntegrand}, written $\omega_i$ or $\nu_i$; it carries its
statistics.

\subsection{Real frequency}

\begin{function}{AnalyticContinuation}
\syntax{AnalyticContinuation[expr, w -> omega]}
\end{function}
Continues $i w\to\omega+i0^+$: the retarded function, understood at $\omega+i0^+$. The message
\pkg{AnalyticContinuation::periodicity} warns when statistical factors still depend on \pkg{w}, which happens when its
statistics was not declared before the sums.

\begin{function}{RetardedImaginaryPart}
\syntax{RetardedImaginaryPart[expr, omega]}
\end{function}
The imaginary part of a retarded function of $\omega+i0^+$ with real coefficients (real dispersions and interactions):
$1/(\omega-E+i0^+)\to-\pi\,\delta(\omega-E)$, and a pole of order $k$ gives
$-\pi\frac{(-1)^{k-1}}{(k-1)!}\delta^{(k-1)}(\omega-E)$. It works inside \pkg{MomentumSum}; the result is a sum of delta
functions, each with its coefficient. For the bubble, continued to $\omega+i0^+$, the imaginary part is $-\operatorname{Im}\chi_0^R$:
\excode{lindhard_imaginary_part}
\exformula{lindhard_imaginary_part}

\begin{function}{EnergyDelta}
\syntax{EnergyDelta[x]}
\end{function}
The inert delta function $\delta(x)$ of an energy; \pkg{Derivative[k][EnergyDelta][x]} is $\delta^{(k)}(x)$.

\begin{function}{ZeroTemperatureLimit}
\syntax{ZeroTemperatureLimit[expr]}
\end{function}
The limit $T\to0$ of the statistical factors: $n_F(e)\to\theta(-e)$ (\pkg{UnitStep[-e]}), $n_B(e)\to-\theta(-e)$, and
$n_F^{(k)}(e)\to-\delta^{(k-1)}(e)$, $n_B^{(k)}(e)\to\delta^{(k-1)}(e)$. Since $n_F(0)=\frac12$ at every temperature
while \pkg{UnitStep[0]} is 1, an energy that may vanish should be substituted before the limit is taken. The terms of
the sums of logarithms go as $T\ln(1-n_F(e))\to e\,\theta(-e)$ and $T\ln(1+n_B(e))\to0$.
\excode{lindhard_zero_temperature}
\exformula{lindhard_zero_temperature}

\subsection{Cutting rules}\label{sec:cutting}

The imaginary part of a retarded function is a sum over the ways of cutting the diagram into two parts, the lines that
are cut being on shell~\cite{Cutkosky1960}, each cut with its factors. At finite temperature every cut line carries the factor of its emission
($1-n_F$, $1+n_B$) or absorption ($n_F$, $n_B$), and the inverse process adds the factor $1\pm e^{-\omega/T}$
(fermions, bosons)~\cite{Weldon1983,KobesSemenoff1985,KobesSemenoff1986,LandsmanVanWeert1987,LeBellac1996,KapustaGale2006}.
The package derives these rules in imaginary time, independently of the residues: the points joined by instantaneous
interactions form the time vertices of the diagram; the lines inside a time vertex are static (a fermion line is
$G(\tau=0^-)$, that is $n_F$); the other lines form chains from the vertex where the external frequency enters to the one
where it leaves, all the lines of a chain carrying the same frequency. Each chain is on shell at the poles of its
product of propagators, a double pole (a static self-energy on the line) giving a derivative of the delta function.

\begin{function}{CuttingRuleImaginaryPart}
\syntax{CuttingRuleImaginaryPart[diagrams, externals, rules, w -> omega, options]}
\end{function}
The imaginary part of the retarded two-point function from the cutting rules. Arguments and options are those of
\fnref{EvaluateMatsubara}; the frequency \pkg{w} must enter the diagram at one external point and leave it at another
(\pkg{I w} and \pkg{-I w} in the frequencies of exactly two external points), else the result is
\pkg{CuttingRuleImaginaryPart::frequency} and \pkg{\$Failed}. A term in which it enters and leaves at the same time
vertex, such as the Hartree and Fock self-energies, does not depend on it and gives 0. A fermion line at equal times
inside a time vertex is $n_F(\xi)-(1-\eta)/2$, with $\eta$ the value of \opt{ConvergenceFactor}. The second-order self-energy (compare
with the residues, \fnref{RetardedImaginaryPart} of \pkg{AnalyticContinuation[sigma2, w -> omega]}, which gives the same
expression):
\excode{sigma_cutting_rules}
\exformula{sigma_cutting_rules}
\textbf{Present limits.} The rules are implemented for two-point functions whose propagating lines form chains between
the two external time vertices: bubbles, sunsets and the exchange of a boson, with static insertions such as Hartree
and Fock on their lines. The external time vertices may have any number of propagating lines, the others must have
two. A diagram with an inner time vertex of more than two propagating lines --- a vertex correction, two bubbles joined
by an interaction, and so, for the electron gas, every skeleton diagram of the self-energy from the third order on ---
gives \pkg{CuttingRuleImaginaryPart::class} and \pkg{\$Failed}: it needs the general circling rules of real-time thermal
field theory, with retarded and advanced uncut lines~\cite{KobesSemenoff1985,KobesSemenoff1986}, which are not
implemented. Of the 42 diagrams of the perturbative self-energy at third order, the 10 skeletons are of this kind; the
other 32 are handled: the 12 sunsets with a Hartree or Fock insertion, and the 20 Hartree and Fock diagrams with their
line dressed to second order, which do not depend on the frequency. The energies are generic (no anomalous terms), and
the coefficients are taken to be real.

\begin{function}{CheckCuttingRules}
\syntax{CheckCuttingRules[diagrams, externals, rules, w -> omega, options]}
\end{function}
Compares, term by term and for each delta function, the imaginary part from the Matsubara sums with the cutting rules;
the statistical factors are compared through $e^{\epsilon/T}$. It gives 0 when they agree, else the list of the
differences (\pkg{\$Failed} for a term to which the cutting rules do not apply), and \pkg{\$Failed} for a frequency
\pkg{w} that does not cross the diagram, as in \fnref{CuttingRuleImaginaryPart}. Options as in
\fnref{EvaluateMatsubara}. For the second-order self-energy and for the bubble:
\excode{check_cutting_rules}
\extext{check_cutting_rules}

% ------------------------------------------------------------------ 17
\section{Evaluation: dimensional regularization}\label{sec:dimreg}

At zero temperature, in Euclidean space, the diagrams of a translation-invariant theory are integrals over loop
momenta. In $d=4-2\varepsilon$ dimensions their ultraviolet divergences become poles in
$\varepsilon$~\cite{tHooftVeltman1972,BolliniGiambiagi1972}, which minimal subtraction removes.
\fnref{DiagramLoopIntegral} routes the momenta through the lines of the diagrams of the package, as
\fnref{EvaluateMatsubara} does, and gives their loop integrals (\fnref{LoopIntegral}); \fnref{IntegrateLoops} evaluates
them, in closed form in $d$, when their loops can be done one after the other as one-loop integrals; \fnref{EpsilonExpand}
and \fnref{PolePart} expand the results in $\varepsilon$. For the counterterms of a massless theory,
\fnref{InfraredRearrangedIntegral} lets one momentum flow through each diagram so that no infrared divergence appears,
and \fnref{BetaFunctionFromZ} and \fnref{AnomalousDimensionFromZ} give the renormalization-group functions from the
renormalization constants. The renormalization of $\phi^4$ theory to two loops (section~\ref{sec:dimreg-phi4}) uses all
of them.

\subsection{Conventions}\label{sec:dimreg-conventions}

\begin{itemize}
\item \textbf{Action.} Euclidean, $S=\int\bigl(\frac12(\partial\phi)^2+\frac12m^2\phi^2\bigr)+S_{\rm int}$, with
  $S_{\rm int}=\frac{\lambda}{4!}\int\phi^4$ for $\phi^4$ theory: the free generating functional \pkg{Z0} and the
  interaction \pkg{Phi4Vertex} of section~\ref{sec:perturbation}. A line is $D_0(q)=1/(q^2+m^2)$,
  \fnref{ScalarPropagator}\pkg{[q, m]} (the rule \pkg{D0 -> \{"Scalar", m\}}). The diagrams have the coefficients of the
  package. The vertex functions are the derivatives of the 1PI action at zero mean field,
  $\Gamma^{(n)}=\delta^n\Gamma/\delta\bar\phi^n$ (\fnref{EffectiveAction1PI}): $\Gamma^{(2)}=G^{-1}=p^2+m^2-\Sigma$, with the
  self-energy $\Sigma$ of $G=G_0+G_0\Sigma G_0+\cdots$, and the four-point vertex $\Gamma^{(4)}=\lambda-\cdots$; the amputated
  1PI parts of the connected functions give the same, $\Sigma$ and minus $\Gamma^{(4)}$ (the check).
\item \textbf{Dimension and measure.} $d=4-2\varepsilon$, and every loop is
  \[
  \int_k\equiv\mu^{2\varepsilon}\Bigl(\frac{e^{\gamma_E}}{4\pi}\Bigr)^{\varepsilon}\int\!\frac{d^dk}{(2\pi)^d},\qquad
  \int_k\frac{1}{(k^2+\Delta)^n}=\frac{e^{\varepsilon\gamma_E}\mu^{2\varepsilon}}{16\pi^2}\,\frac{\Gamma(n-d/2)}{\Gamma(n)}\,
  \Delta^{d/2-n},
  \]
  with $\mu$ = \fnref{RenormalizationScale}. With this measure the minimal subtraction of the poles in $\varepsilon$ is the
  $\overline{\rm MS}$ scheme and $\mu$ its scale: no $\gamma_E$ and no $\log4\pi$ appear in the results, and every loop
  is $\kappa=1/(16\pi^2)$ times a function of $\varepsilon$, of the logarithms $\log(p^2/\mu^2)$, $\log(m^2/\mu^2)$ and of
  ratios of masses and momenta.
\item \textbf{Momenta.} The momentum $p_i$ enters the diagram at the external label $x_i$ (\pkg{\{x1 -> p1, \dots\}}; they
  add up to zero), and a line carries its momentum from its second label to its first, written with a positive first
  term (the lines are even). External momenta are symbols, their squares \pkg{p\^{}2}; the square of a combination is
  \fnref{MomentumSquared}\pkg{[q]}.
\item \textbf{Renormalization.} $\phi_B=Z_\phi^{1/2}\phi$, $m_B^2=Z_{m^2}m^2/Z_\phi$ and
  $\lambda_B=\mu^{2\varepsilon}(e^{\gamma_E}/4\pi)^{\varepsilon}Z_\lambda\lambda$; $\beta=\mu\,d\lambda/d\mu$,
  $\gamma_\phi=\frac12\mu\,d\log Z_\phi/d\mu$ and $\gamma_{m^2}=\mu\,d\log m^2/d\mu$, as
  in~\cite{PeskinSchroeder,ZinnJustin2002,KleinertSchulteFrohlinde2001}.
\end{itemize}

\subsection{Loop integrals}

\begin{function}{DiagramLoopIntegral}
\syntax{DiagramLoopIntegral[diagrams, \{x1 -> p1, x2 -> p2, \dots\}, rules, options]}
\end{function}
The momentum-space form of the diagrams: each canonical term becomes its coefficient times
\pkg{LoopIntegral[}product of its lines\pkg{, }loop momenta\pkg{]} (a tree has no integral). The momenta are routed as in
\fnref{EvaluateMatsubara}: labels joined by delta functions are one point, the lines of a spanning tree carry what
conservation gives them, and every other line closes a loop and carries a loop momentum, \pkg{LoopMomentum[i]} or a name
of the option \opt{LoopMomenta} \pkg{-> \{k1, \dots\}}. The rule \pkg{D0 -> \{"Scalar", m\}} makes \pkg{D0[x, y]} the line
\pkg{ScalarPropagator[q, m]} $=1/(q^2+m^2)$ ($m=0$ for a massless line), \pkg{Q -> \{"Kernel", f\}} gives the line
\pkg{Q[x, y]} the value \pkg{f[q]}, an insertion such as a counterterm $c\,q^2$ (\pkg{f} must be even). The factors of a
term must be lines of the rules, delta functions or constants (\pkg{DiagramLoopIntegral::factor}), and the external
momenta must add up to zero on each connected part (\pkg{::conservation}). The self-energy of $\phi^4$ theory to second
order, $\Sigma=D_0^{-1}-\Gamma^{(2)}$ from the second derivative of the 1PI action at zero mean field
(section~\ref{sec:dimreg-conventions}), and its loop integrals:
\excode{dimreg_sigma}
\expictures{dimreg_sigma}
The amputated one-particle-irreducible part of the propagator gives the same:
\excode{dimreg_sigma_check}
\extext{dimreg_sigma_check}
\excode{dimreg_integrand}
\exformula{dimreg_integrand}
The momentum $p$ does not flow through the tadpoles, and the double tadpole is two loops, one inside the other. A
kernel line in a loop:
\excode{dimreg_kernel}
\exformula{dimreg_kernel}

\begin{function}{LoopIntegral}
\syntax{LoopIntegral[f, \{k1, k2, \dots\}]}
\end{function}
The inert loop integral $\int_{k_1k_2\cdots}f$, each loop with the measure of section~\ref{sec:dimreg-conventions}; $f$ is
a product of \fnref{ScalarPropagator} factors, with numerators that are squares of momenta. It depends on $\varepsilon$
only through $d$, so that \fnref{EpsilonExpand} refuses it (\pkg{EpsilonExpand::loop}): it is evaluated first by
\fnref{IntegrateLoops}. So is a \fnref{MinkowskiLoopIntegral} of section~\ref{sec:qed}, evaluated by \fnref{IntegrateOneLoop}.

\begin{function}{ScalarPropagator}
\syntax{ScalarPropagator[q, m]}
\end{function}
The line $1/(q^2+m^2)$ of momentum $q$ and mass $m$; \pkg{ScalarPropagator[q, m]\^{}a} is $1/(q^2+m^2)^a$, for any power
$a$: the one-loop integrals inside a diagram leave massless lines of powers such as $1+\varepsilon$. A massless line of
zero momentum, \pkg{ScalarPropagator[0, 0]}, is infinite (\pkg{IntegrateLoops::zero}).

\begin{function}{MomentumSquared}
\syntax{MomentumSquared[q]}
\end{function}
The square of a momentum $q$: $p^2$ for a symbol, $c^2p^2$ for a number times a symbol; a combination stays, with the sign
that makes its first term positive.
\excode{dimreg_momentum_squared}
\extext{dimreg_momentum_squared}

\begin{function}{RenormalizationScale}
\syntax{RenormalizationScale}
\end{function}
The scale $\mu$ of the measure of every loop, the $\overline{\rm MS}$ scale; written $\mu$.

\subsection{One-loop integrals and the loops of a diagram}

\begin{function}{IntegrateLoops}
\syntax{IntegrateLoops[expr, eps]}
\end{function}
Evaluates the \fnref{LoopIntegral} objects of \pkg{expr} in $d=4-2\varepsilon$ when their loops can be done one after the
other. At each step a loop momentum that appears in at most two different propagators, with coefficient $\pm1$, is
integrated: one propagator gives a tadpole (\fnref{TadpoleIntegral}); two massless ones a G-function and a new massless
line of the momentum $P$ that flows through them,
\[
\int_k\frac{1}{(k^2)^a\,((k+P)^2)^b}=\frac{e^{\varepsilon\gamma_E}\mu^{2\varepsilon}}{16\pi^2}\,G(a,b)\,(P^2)^{2-\varepsilon-a-b};
\]
two lines of which one is massive a \fnref{BubbleIntegral}, if $P$ is free of the other loop momenta. When no loop
momentum qualifies, a change of loop momenta that makes the momentum of a line a loop momentum is tried. This covers the
massless propagator-type integrals made of nested one-loop integrals (the sunset, bubbles inserted in lines, the
two-loop vertex diagrams with one momentum through two of their legs) and products of massive one-loop factors.
Numerators $q^2$ (from kernel lines) are written $1/D_0(q)-m^2$ with the mass of a line of the same momentum, and a
numerator in a bubble with a massive line is reduced to tadpoles (the odd powers of $k\cdot P$ vanish, the even ones are
averaged over the directions of $k$). Scaleless integrals are 0. Lines outside the loop integrals become
$1/(q^2+m^2)^a$. The massive tadpole and the massless sunset, from the diagrams above:
\excode{dimreg_tadpole}
\exformula{dimreg_tadpole}
\excode{dimreg_sunset}
\exformula{dimreg_sunset}
With massless lines, as in the sunset above, the double tadpole is 0: its inner loop is a massless tadpole. A massive bubble stays as \fnref{BubbleIntegral}, and a
numerator next to a massive line gives tadpoles:
\excode{dimreg_massive_bubble}
\exformula{dimreg_massive_bubble}
\excode{dimreg_numerator}
\exformula{dimreg_numerator}
Messages: \pkg{IntegrateLoops::irreducible}, an integral that does not reduce to nested one-loop ones (the massive sunset,
three massive lines in two loops), and \pkg{::numerator}, a numerator that is not the square of a momentum ($k\cdot p$):
the integral stays unevaluated; \pkg{::ir}, an infrared divergence (a massless line of power 2 or more at $\varepsilon=0$
in a loop), so that the poles of the result are not only ultraviolet ones; \pkg{::scaleless}, a logarithmically
divergent scaleless integral, which is 0 because its ultraviolet and infrared poles cancel: its ultraviolet pole is lost;
\pkg{::zero}, a massless line of zero momentum outside the loops. A massless tadpole, which is only power divergent, is 0
without a message.

\begin{function}{TadpoleIntegral}
\syntax{TadpoleIntegral[a, m, eps]}
\end{function}
The tadpole $\int_k1/(k^2+m^2)^a$, exact in $d$; 0 for $m=0$ and for an integer power $a\le0$.
\excode{dimreg_tadpole_integral}
\exformula{dimreg_tadpole_integral}

\begin{function}{GFunction}
\syntax{GFunction[a, b, eps]}
\end{function}
The coefficient of the massless one-loop integral, for any powers $a$, $b$:
\begin{gather*}
\int\!\frac{d^dk}{(2\pi)^d}\,\frac{1}{(k^2)^a\,((k+p)^2)^b}=\frac{G(a,b)}{(4\pi)^{d/2}}\,(p^2)^{d/2-a-b},\\
G(a,b)=\frac{\Gamma(d/2-a)\,\Gamma(d/2-b)\,\Gamma(a+b-d/2)}{\Gamma(a)\,\Gamma(b)\,\Gamma(d-a-b)}.
\end{gather*}
\excode{dimreg_g_function}
\exformula{dimreg_g_function}

\begin{function}{BubbleIntegral}
\syntax{BubbleIntegral[\{a, b\}, p2, \{m1, m2\}, eps]}
\end{function}
The bubble $\int_k1/\bigl((k^2+m_1^2)^a((k+p)^2+m_2^2)^b\bigr)$ at $p^2=$ \pkg{p2}. It has closed forms in $d$ when both lines
are massless (a G-function), when a power is an integer $\le0$ (a numerator, reduced to tadpoles), and at \pkg{p2 = 0}
with equal masses, with a massless line, or with different masses for $a=b=1$ (partial fractions). Otherwise it stays, written $B_{a,b}(p^2;m_1,m_2)$, as
the Feynman-parameter integral
\[
\frac{e^{\varepsilon\gamma_E}\mu^{2\varepsilon}}{16\pi^2}\,\frac{\Gamma(a+b-d/2)}{\Gamma(a)\,\Gamma(b)}\int_0^1\!dx\,x^{a-1}(1-x)^{b-1}
\Delta^{d/2-a-b},\qquad \Delta=x(1-x)\,p^2+x\,m_1^2+(1-x)\,m_2^2,
\]
which \fnref{EpsilonExpand} expands. The massless bubble:
\excode{dimreg_bubble_integral}
\exformula{dimreg_bubble_integral}

\subsection{Expansion in \texorpdfstring{$\varepsilon$}{epsilon}}

\begin{function}{EpsilonExpand}
\syntax{EpsilonExpand[expr, eps, n]}
\end{function}
The Laurent expansion of \pkg{expr} in $\varepsilon$ up to $\varepsilon^n$ included, computed without \pkg{Series}:
$\Gamma(n_0+c\varepsilon)$ through the polygamma functions (at positive integers through $\zeta(k)$) and the recurrence at
the poles; powers $X^{c\varepsilon}$ through logarithms, written $\log(X/\mu^2)$ when the powers of $\mu$ balance them, as
they do in every dimensionless combination; rational functions and exponentials of $\varepsilon$ term by term; massive
bubbles through their Feynman parameters, in closed form for the finite part of $B_{1,1}$ and of the bubbles with
$a+b=3$ (logarithms, the square root of the K\"all\'en function, $\operatorname{artanh}$ for equal masses), as
\fnref{FeynmanParameterIntegral} beyond. Products and powers of factors with poles are expanded to the order asked. An
unevaluated \fnref{LoopIntegral} or \fnref{MinkowskiLoopIntegral} gives \pkg{EpsilonExpand::loop} and \pkg{\$Failed}; so do a function that is not
expanded (\pkg{::function}) and a bubble with powers that are not positive or that depend on $\varepsilon$, or that is
infrared divergent (\pkg{::bubble}). The massless bubble, the sunset and the bubble with two equal masses, in units of
$\kappa=1/(16\pi^2)$ per loop:
\excode{dimreg_expand_bubble}
\exformula{dimreg_expand_bubble}
\excode{dimreg_expand_sunset}
\exformula{dimreg_expand_sunset}
\excode{dimreg_expand_massive}
\exformula{dimreg_expand_massive}
Three poles in a product:
\excode{dimreg_zeta}
\exformula{dimreg_zeta}

\begin{function}{PolePart}
\syntax{PolePart[expr, eps]}
\end{function}
The part of the Laurent expansion with negative powers of $\varepsilon$, \pkg{EpsilonExpand[expr, eps, -1]}: the
$\overline{\rm MS}$ counterterm is minus it.
\excode{dimreg_pole_part}
\exformula{dimreg_pole_part}

\begin{function}{FeynmanParameterIntegral}
\syntax{FeynmanParameterIntegral[f, FeynmanParameter]}
\end{function}
The inert integral $\int_0^1dx\,f$ over the Feynman parameter $x$, left by \fnref{EpsilonExpand} for the coefficients of a
massive bubble that it does not give in closed form; \pkg{N} evaluates it with \pkg{NIntegrate} (\pkg{N[expr, n]} to
\pkg{n} digits). The order $\varepsilon$ of a bubble with two masses:
\excode{dimreg_feynman_parameter}
\extext{dimreg_feynman_parameter}

\begin{function}{FeynmanParameter}
\syntax{FeynmanParameter}
\end{function}
The variable $x$ of \fnref{FeynmanParameterIntegral}, written $x$. \pkg{FeynmanParameter[i]}, written $x_i$, are the
Feynman parameters of \fnref{IntegrateOneLoop}, in \fnref{ParameterIntegral} and \fnref{SimplexIntegral}
(section~\ref{sec:qed}).

\subsection{Counterterms}

\begin{function}{InfraredRearrangedIntegral}
\syntax{InfraredRearrangedIntegral[diagrams, \{x1, \dots, xn\}, p, rules, eps, options]}
\end{function}
Evaluates each term of the diagrams with one momentum \pkg{p} entering at one external label and leaving at another, the
other external momenta zero: the pairs of labels are tried in order until the term reduces to nested one-loop integrals
without infrared divergences and without logarithmically divergent scaleless integrals. In $\overline{\rm MS}$ the
counterterms of a logarithmically divergent function, once its subdivergences are subtracted, are polynomial in the
external momenta and in the masses, that is constants, and can be computed at any momenta and masses that avoid infrared
divergences: this is the infrared rearrangement of the multiloop calculations of renormalization-group
functions~\cite{ChetyrkinTkachov1981,KleinertSchulteFrohlinde2001}. Each term is routed by itself; the one-loop
subdiagrams of a term and the counterterms that subtract them, which come from the one-loop terms with $\lambda\to\lambda
Z_\lambda$, are then at the same momenta. A part of a term with a massless tadpole, which is only power divergent, is 0:
in a massless theory it has no counterterm. A term without a valid routing gives \pkg{InfraredRearrangedIntegral::norouting},
with the problems found, and \pkg{\$Failed}. Only the counterterms left after the subtraction of the subdivergences are
physical: the finite parts, and the pole parts of the single rearranged diagrams with their $\log(p^2/\mu^2)/\varepsilon$,
depend on the routing. With the option \opt{Routing} \pkg{-> True} the result is the list of \pkg{\{term, \{label in, label
out\}, value\}}. The 1PI four-point vertex of $\phi^4$ theory to third order, the fourth derivative of the 1PI action
(\fnref{EffectiveAction1PI}) at $\bar\phi=0$, its diagrams at order $\lambda^3$, and the three bubbles at order $\lambda^2$:
\excode{dimreg_gamma4}
\expictures{dimreg_gamma4}
\excode{dimreg_rearranged_one_loop}
\exformula{dimreg_rearranged_one_loop}
At order $\lambda^3$ the pole parts, in units of $\lambda^3\kappa^2$, with $L=\log(p^2/\mu^2)$: the six diagrams with a bubble
in a vertex, routed from the vertex with two external legs to one with one; the three chains of two bubbles; and the
three bubbles with a tadpole on a line, which are 0:
\excode{dimreg_rearranged_routing}
\extext{dimreg_rearranged_routing}

\begin{function}{MassInsertion}
\syntax{MassInsertion[diagrams, G0, z]}
\end{function}
$-\partial/\partial m^2$ of the diagrams through their lines $G_0=1/(q^2+m^2)$: the sum, over the lines \pkg{G0[x, y]}, of
the diagram with that line replaced by \pkg{G0[x, z] G0[z, y]}, the insertion of $\phi^2/2$ at the new label \pkg{z} (a
symbol that does not appear in the diagrams, otherwise \pkg{MassInsertion::label} and \pkg{\$Failed}). So
$\partial\Sigma/\partial m^2$ is minus the insertion in $\Sigma$, logarithmically divergent; at $m=0$ the momentum entering at
\pkg{z} makes its infrared rearrangement possible, and it gives $Z_{m^2}$:
\excode{dimreg_mass_insertion}
\expictures{dimreg_mass_insertion}

\begin{function}{BetaFunctionFromZ}
\syntax{BetaFunctionFromZ[Z, g, eps, n, options]}
\end{function}
The beta function $\beta(g)=\mu\,dg/d\mu$ at $\varepsilon=0$, to order $g^n$, of a coupling renormalized as
$g_B=\mu^{2\varepsilon}Zg$ (in $\overline{\rm MS}$, with the factor $(e^{\gamma_E}/4\pi)^\varepsilon$ of the measure): from
$\mu\,dg_B/d\mu=0$, the beta function in $d$ dimensions is $\beta_\varepsilon=-2\varepsilon g/(1+g\,\partial_g\log Z)$, and
$\beta=\beta_\varepsilon+2\varepsilon g$ is expanded in $g$. The limit $\varepsilon\to0$ must exist, which checks the higher
poles of $Z$ against the single ones; otherwise \pkg{BetaFunctionFromZ::pole} and \pkg{\$Failed}.
Option \opt{ScalePower} \pkg{-> s} for $g_B=\mu^{s\varepsilon}Zg$: $s=2$ for $\lambda\phi^4$ (the default), $s=1$ for a gauge or
a Yukawa coupling in $d=4-2\varepsilon$.

\begin{function}{AnomalousDimensionFromZ}
\syntax{AnomalousDimensionFromZ[Z, g, eps, betaEpsilon, n]}
\end{function}
$\mu\,d\log Z/d\mu=\beta_\varepsilon\,\partial_g\log Z$ at $\varepsilon=0$, to order $g^n$, with \pkg{betaEpsilon}
$=-s\varepsilon g+\beta(g)$ the beta function in $d=4-2\varepsilon$; \pkg{AnomalousDimensionFromZ::pole} and \pkg{\$Failed} if
the limit does not exist. For a field $\phi_B=Z_\phi^{1/2}\phi$, $\gamma_\phi=\frac12\mu\,d\log Z_\phi/d\mu$.

\subsection{The renormalization of \texorpdfstring{\smash{$\phi^4$}}{phi4} theory to two loops}\label{sec:dimreg-phi4}

The bare functions are the package's diagrams of the examples above, massless and with the infrared rearrangement: the
self-energy with $p$ through $a$ and $b$, the vertex $\Gamma^{(4)}$ to order $\lambda^3$, and $\partial\Sigma/\partial m^2$
with the momentum entering at the insertion $u$. The pole part of $\partial\Sigma/\partial m^2$ (here with $\lambda$ in units
of $16\pi^2$) has a non-local term, $\log(p^2/\mu^2)/\varepsilon$, which the counterterm diagrams cancel:
\excode{dimreg_bare}
\exformula{dimreg_bare}
The renormalized functions
\[
\Gamma_R^{(2)}=Z_\phi\bigl(p^2-\Sigma_B\bigr),\qquad \Gamma_R^{(4)}=Z_\phi^2\,\Gamma_B^{(4)},\qquad
\frac{\partial\Gamma_R^{(2)}}{\partial m^2}=Z_{m^2}\Bigl(1-\frac{\partial\Sigma_B}{\partial m_B^2}\Bigr)
\]
must be finite, order by order, with the bare functions taken at $\lambda_B$. Every power of $\lambda_B$ comes with a loop
more than the tree, so that with $\lambda_B=\mu^{2\varepsilon}(e^{\gamma_E}/4\pi)^\varepsilon Z_\lambda\lambda$ every loop has the
measure of the package and $\Gamma^{(4)}$ an overall factor $\mu^{2\varepsilon}(e^{\gamma_E}/4\pi)^\varepsilon$, which is divided
out: $\lambda\to\lambda Z_\lambda$ in the bare functions gives exactly the counterterm diagrams. With unknown coefficients of
the poles (\pkg{c[1]} the single pole at one loop, \pkg{c[2]} and \pkg{c[3]} the double and the single pole at two loops,
in units of $\kappa$ per loop: the symbol \pkg{kappa} stays in the $Z$'s and is given its value $1/(16\pi^2)$ in the
functions that must be finite), the poles must vanish for every value of $L=\log(p^2/\mu^2)$: the non-local ones cancel by
themselves once the one-loop counterterms are in, and the others fix the coefficients:
\excode{dimreg_z_solution}
\extext{dimreg_z_solution}
that is, $Z_\phi=1-\frac{\lambda^2\kappa^2}{24\varepsilon}$, $Z_{m^2}=1+\frac{\lambda\kappa}{2\varepsilon}+\lambda^2\kappa^2
\bigl(\frac{1}{2\varepsilon^2}-\frac{1}{4\varepsilon}\bigr)$ and
\excode{dimreg_z_lambda}
\exformula{dimreg_z_lambda}
The double poles follow from the single ones (\pkg{zl[2] = zl[1]\^{}2}), as the renormalization group requires, and
the renormalization-group functions are
\excode{dimreg_beta}
\exformula{dimreg_beta}
\excode{dimreg_gamma_phi}
\exformula{dimreg_gamma_phi}
\excode{dimreg_gamma_m2}
\exformula{dimreg_gamma_m2}
These are the values of the
textbooks~\cite{PeskinSchroeder,ZinnJustin2002,KleinertSchulteFrohlinde2001}, the case $N=1$ of the $O(N)$ model
(Kleinert and Schulte-Frohlinde write $d=4-\varepsilon$, so that their poles $1/\varepsilon$ are the $1/(2\varepsilon)$ here).
In $d=4-\epsilon$ ($\epsilon=2\varepsilon$), at the fixed point of $-\epsilon\lambda+\beta(\lambda)$, they give the
exponents of the Ising universality class to order $\epsilon^2$~\cite{WilsonFisher1972}, $\eta=\epsilon^2/54$ and
$\nu=\frac12+\frac{\epsilon}{12}+\frac{7\epsilon^2}{162}$; the showcase computes them.
The test suite checks the same constants, the absence of poles in the three renormalized functions, the consistency of
the double poles, and the one-loop mass counterterm against the massive tadpole.

\subsection{Limits, and what would come next}

\begin{itemize}
\item \textbf{Nested one-loop integrals only.} At each step a loop momentum must appear in at most two different
  propagators, after at most one change of loop momenta. Not covered: the massive sunset and the massive two-loop vacuum
  integrals; genuine two-loop topologies with three propagators in a loop (the kite of the three-loop $\beta$ function of
  $\phi^4$); one-loop triangles and boxes (needed for $\Gamma^{(4)}$ at generic momenta; in Minkowski space,
  \fnref{IntegrateOneLoop} of section~\ref{sec:qed} writes any one-loop integral as Feynman-parameter integrals, which
  \pkg{N} evaluates); a massive bubble inside a further
  loop (its momentum must be external).
\item \textbf{Numerators} are squares of momenta; there are no tensor numerators $k\cdot p$, $k^\mu k^\nu$ (at one loop, in
  Minkowski space, \fnref{IntegrateOneLoop} of section~\ref{sec:qed} does them).
\item \textbf{Massive bubbles} are in closed form to $\varepsilon^0$ for $a=b=1$ and $a+b=3$; beyond, their coefficients are
  \fnref{FeynmanParameterIntegral}, evaluated numerically.
\item \textbf{Infrared rearrangement.} Only the $Z$ factors, the counterterms left after the subtraction of the
  subdivergences, are physical. The first routing without infrared problems is taken, which is legitimate for the
  ultraviolet counterterms of logarithmically divergent functions; quadratically divergent ones are first made
  logarithmic by a derivative ($\partial/\partial p^2$, or $\partial/\partial m^2$ by \fnref{MassInsertion}).
\end{itemize}
The natural next steps:
\begin{itemize}
\item \textbf{Integration by parts}~\cite{ChetyrkinTkachov1981}: identities between integrals of the same topology. One of
  them, the triangle rule, writes the massless two-loop master diagram (the kite, finite in four dimensions, where it is
  $6\zeta(3)/p^2$ in units of $\kappa^2$) as a combination of G-functions; with it the massless two-loop propagator
  integrals, and the three-loop $\beta$ function of $\phi^4$ by infrared rearrangement, are within reach. General two-loop
  integrals with masses need a systematic reduction to master integrals (linear systems over rational functions of $d$).
\item \textbf{Massive two-loop integrals} (the sunset with masses, the vacuum integrals with three masses): closed forms
  in hypergeometric functions, expansions in ratios of masses and momenta, or numerical integration over Feynman
  parameters. The $\overline{\rm MS}$ constants do not need them, the finite parts of massive functions do.
\item \textbf{Tensor numerators}, for gauge theories and fermions, in nested integrals: G-functions with numerators. At one
  loop, the Lorentz indices, the Dirac algebra in $d$ dimensions and the tensor integrals are in section~\ref{sec:qed}.
\end{itemize}

% ------------------------------------------------------------------ 18
\section{Evaluation: QED at one loop}\label{sec:qed}

Quantum electrodynamics is a theory of the package like any other: the electron and the photon are the fields of a
Gaussian generating functional, and the vertex is a local kernel. A label stands for a point together with an index ---
the spinor index of an end of an electron line, the Lorentz index of an end of a photon line --- so that the diagrams of
the package, with their weights and the signs of their closed fermion loops, are the Feynman diagrams of QED in
condensed notation. \fnref{QEDIntegrand} writes them in momentum space: it routes the momenta through the lines, as
\fnref{EvaluateMatsubara} and \fnref{DiagramLoopIntegral} do, and follows the spinor indices along the fermion lines, so
that an open line becomes a product of Dirac matrices and a closed loop a trace. \fnref{DiracSimplify} does the Dirac
algebra in $d$ dimensions, \fnref{IntegrateOneLoop} the loop integral by Feynman parameters, exactly in $d$;
\fnref{ReduceParameterIntegrals} and \fnref{ExpandParameterIntegrals} do the integrals over the Feynman parameters in
closed form where they can and expand in $\varepsilon$, and \fnref{VertexFormFactors} splits a vertex into its form
factors. The electron self-energy, the vacuum polarization and the vertex at one loop (section~\ref{sec:qed-results})
use all of them; \pkg{Tests/QEDOneLoop.wl} is the same computation as a script, which takes about ten seconds.

\subsection{Conventions}\label{sec:qed-conventions}

\begin{itemize}
\item \textbf{Dirac algebra.} Minkowski metric $(+,-,-,-)$, $\{\gamma^\mu,\gamma^\nu\}=2g^{\mu\nu}$, $\operatorname{tr}\mathbf 1=4$,
  in $d$ dimensions: $g^\mu{}_\mu=d$. The Lorentz indices are abstract: an index that appears twice in a term is
  contracted (one upper, one lower; in the formulas of the package both are written upper), and there are no explicit
  components. There is no $\gamma_5$.
\item \textbf{Action and Feynman rules}, as in~\cite{PeskinSchroeder} but with $e>0$: the electron has the charge $-e$,
  $\mathcal L=\bar\psi(i\slashed\partial-m)\psi-\frac14F_{\mu\nu}F^{\mu\nu}+e\bar\psi\gamma^\mu\psi A_\mu$ (Peskin and Schroeder
  write $-e\bar\psi\gamma^\mu\psi A_\mu$ with $e<0$ for the electron), so that the vertex is $+ie\gamma^\mu$, the electron line $i(\slashed p+m)/(p^2-m^2+i0)$ and the photon line, in a covariant gauge,
  \[
  \frac{-i}{q^2-\lambda^2+i0}\Bigl(g^{\mu\nu}-(1-\xi)\,\frac{q^\mu q^\nu}{q^2-\lambda^2}\Bigr),
  \]
  with $\xi=1$ in the Feynman gauge and $\xi=0$ in the Landau gauge; a photon mass $\lambda$ is an infrared regulator
  ($\lambda=0$ by default).
\item \textbf{The theory in condensed notation.} With the propagators $S(x,y)=\langle T\psi(x)\bar\psi(y)\rangle_0$ (its row
  is the spinor index of $x$, its column that of $y$) and $D(x,y)=\langle TA(x)A(y)\rangle_0$, the free generating
  functional is $Z_0=\exp\bigl(\frac12\int_{xy}J(x)D(x,y)J(y)+\int_{xy}S(x,y)\,\bar\eta(x)\eta(y)\bigr)$, the electron field
  $\psi$ is $\delta/\delta\bar\eta$ and $\bar\psi$ is $-\delta/\delta\eta$. The package expands $e^{-S_{\rm int}}$; for the
  $e^{iS}$ of Minkowski space, with $S_{\rm int}=e\int\bar\psi\gamma^\mu\psi A_\mu$, the interaction is
  $-iS_{\rm int}=e\int_{xyz}V(x,y,z)\,\bar\psi(x)\psi(y)A(z)$ with the local
  kernel $V(x,y,z)=-i\gamma^{\mu_z}_{a_xa_y}\delta(x-y)\delta(x-z)$ (the rule \pkg{V -> \{"Vertex", -I\}}), so that each vertex
  is $(-1)\,e\,(-i\gamma^\mu)=ie\gamma^\mu$. The 1PI amputated functions of the package are then $-i\Sigma(p)$,
  $i\Pi^{\mu\nu}(q)$ and $ie\,\delta\Gamma^\mu(p_2,p_1)$, the diagrams of~\cite{PeskinSchroeder} with $e\to-e$; $\Sigma$,
  $\Pi$, the form factors and $\beta(e)$, even in $e$ or of one sign for both, are those of~\cite{PeskinSchroeder}.
\item \textbf{Dimension and measure.} $d=4-2\varepsilon$, and every loop has the measure of
  section~\ref{sec:dimreg-conventions}, in Minkowski space: $\int_k\equiv\mu^{2\varepsilon}(e^{\gamma_E}/4\pi)^\varepsilon\int
  d^dk/(2\pi)^d$, so that minimal subtraction is the $\overline{\rm MS}$ scheme. After the Wick rotation,
  \[
  \int_l\frac{(l^2)^J}{(l^2-\Delta)^A}=i\,(-1)^{A+J}\,\frac{e^{\varepsilon\gamma_E}\mu^{2\varepsilon}}{16\pi^2}\,
  \frac{\Gamma(J+d/2)\,\Gamma(A-J-d/2)}{\Gamma(d/2)\,\Gamma(A)}\,\Delta^{d/2+J-A},
  \]
  with $\Delta-i0$. Peskin and Schroeder write $d=4-\epsilon$: their $2/\epsilon-\gamma_E+\log4\pi+\log\mu^2$ is
  $1/\varepsilon+\log\mu^2$ here.
\item \textbf{Momenta.} The momentum $p_i$ enters the diagram at the external label $x_i$ (\pkg{\{x1 -> p1, \dots\}}; they
  add up to zero), and a line \pkg{S[x, y]} carries its momentum from $y$ to $x$, along the arrow of the electron. A
  momentum is a linear combination of momentum names --- symbols, or objects such as \pkg{p[1]} and \pkg{LoopMomentum[i]}
  --- with coefficients that are numbers or polynomials in the Feynman parameters.
\item \textbf{Renormalization}, as in~\cite[section~10.3]{PeskinSchroeder}: $\psi_0=Z_2^{1/2}\psi$, $A_0=Z_3^{1/2}A$,
  $m_0=Z_mm$, $e_0=\mu^\varepsilon Z_1Z_2^{-1}Z_3^{-1/2}e$ (with the factor $(e^{\gamma_E}/4\pi)^{\varepsilon/2}$ of the measure),
  $\delta_i=Z_i-1$. The self-energy is $\Sigma=m\,A(p^2)+\slashed p\,B(p^2)$, and the counterterm adds
  $i(\slashed p\,\delta_2-\delta_m)$ to $-i\Sigma$, with $\delta_m=Z_2m_0-m$: in $\overline{\rm MS}$, $\delta_2$ is the pole
  of $B$, $\delta_m$ minus $m$ times the pole of $A$, and $Z_m-1=\delta_m/m-\delta_2$. The vacuum polarization is
  $\Pi^{\mu\nu}=(q^2g^{\mu\nu}-q^\mu q^\nu)\,\Pi(q^2)$ and $\delta_3$ the pole of $\Pi$. Between on-shell spinors,
  $\bar u(p_2)\Gamma^\mu u(p_1)=\bar u(p_2)\bigl(\gamma^\mu F_1+\frac{i\sigma^{\mu\nu}q_\nu}{2m}F_2+q^\mu F_3\bigr)u(p_1)$,
  with $q=p_2-p_1$ and $\sigma^{\mu\nu}=\frac i2[\gamma^\mu,\gamma^\nu]$.
\end{itemize}

\subsection{Dirac algebra}

\begin{function}{GammaMatrix}
\syntax{GammaMatrix[mu]}
\end{function}
The Dirac matrix $\gamma^\mu$ in $d$ dimensions, written $\gamma^\mu$. The index is abstract: an index that appears twice
in a term is contracted, more than twice is an error (\pkg{DiracSimplify::index}), and explicit components,
\pkg{GammaMatrix[0]}, are not supported (\pkg{DiracSimplify::component}). There is no $\gamma_5$: \pkg{GammaMatrix[5]}
would be an ordinary index called 5. \fnref{DiracSimplify} and \fnref{DiracOrder} refuse it (\pkg{GammaMatrix::gamma5}), as
they refuse explicit components, with the result \pkg{\$Failed}; \fnref{IntegrateOneLoop} gives the same messages and
leaves the integral unevaluated.

\begin{function}{SlashedMomentum}
\syntax{SlashedMomentum[p]}
\end{function}
$\slashed p=\gamma^\mu p_\mu$, linear in $p$, written $\slashed p$ (\verb|\slashed| of the \LaTeX{} package \pkg{slashed}).
The momentum is a linear combination of momenta (symbols, or objects such as \pkg{p[1]} and \pkg{LoopMomentum[i]}) with
coefficients that are numbers or polynomials in \pkg{FeynmanParameter[i]}. A symbolic coefficient, as in $x\,p$, $m\,p$ or
$p/m$, is ambiguous: the object is left as it is, with the message \pkg{DiracSimplify::momentum}.

\begin{function}{DiracProduct}
\syntax{DiracProduct[a, b, \dots]}
\end{function}
The product $a\,b\cdots$ of Dirac matrices: \fnref{GammaMatrix}, \fnref{SlashedMomentum}, sums of them and scalars;
\pkg{DiracProduct[]} is the unit matrix, written $\mathbf 1$. \fnref{DiracSimplify} expands it, takes the scalars out and
contracts its indices. A product of matrices written with \pkg{Times} or \pkg{Power}, which would lose their order, is
refused (\pkg{DiracSimplify::product}, \pkg{\$Failed}).

\begin{function}{DiracTrace}
\syntax{DiracTrace[X]}
\end{function}
The trace of the Dirac matrix $X$, $\operatorname{tr}[X]$, with $\operatorname{tr}\mathbf 1=4$ and no $\gamma_5$;
\fnref{DiracSimplify} evaluates it, by
$\operatorname{tr}(a_1\cdots a_n)=\sum_{k=2}^n(-1)^k\,(a_1\cdot a_k)\operatorname{tr}(a_2\cdots\widehat{a_k}\cdots a_n)$.

\begin{function}{SpinorChain}
\syntax{SpinorChain[\{p2, m2\}, X, \{p1, m1\}]}
\end{function}
$\bar u(p_2)\,X\,u(p_1)$, between on-shell spinors: $p_2^2=m_2^2$, $p_1^2=m_1^2$, $\bar u(p_2)\slashed p{}_2=m_2\bar u(p_2)$ and
$\slashed p{}_1u(p_1)=m_1u(p_1)$. \fnref{DiracSimplify} moves $\slashed p{}_1$ to the right and $\slashed p{}_2$ to the left, by
$\{a,b\}=2a\cdot b$, and applies the Dirac equation; the other momenta of $X$ should be written in terms of $p_1$ and
$p_2$ ($q\to p_2-p_1$).

\begin{function}{LorentzMetric}
\syntax{LorentzMetric[mu, nu]}
\end{function}
The metric $g^{\mu\nu}$: symmetric, with $g^\mu{}_\mu=d$ (\fnref{SpacetimeDimension}); abstract indices.

\begin{function}{FourVector}
\syntax{FourVector[p, mu]}
\end{function}
The component $p^\mu$ of the momentum $p$, $\mu$ an abstract index.

\begin{function}{ScalarProduct}
\syntax{ScalarProduct[p, q]}
\end{function}
The Minkowski scalar product $p\cdot q$, symmetric and bilinear; \pkg{ScalarProduct[p, p]} is $p^2$.

\begin{function}{SpacetimeDimension}
\syntax{SpacetimeDimension}
\end{function}
The dimension $d$ of the Dirac algebra, written $d$. In the results of \fnref{IntegrateOneLoop}, and in
\fnref{ExpandParameterIntegrals}, it is $d=4-2\varepsilon$.

\begin{function}{DiracSimplify}
\syntax{DiracSimplify[expr]}
\end{function}
Expands the Dirac matrices of \pkg{expr} (\fnref{DiracProduct}, \fnref{DiracTrace}, \fnref{SpinorChain}), makes the
slashed momenta, four-vectors and scalar products linear, and contracts the repeated indices:
$\gamma^\mu\gamma_\mu=d$, $\gamma^\mu\slashed a\gamma_\mu=(2-d)\,\slashed a$, $\gamma^\mu\slashed a\slashed b\gamma_\mu=4\,a\cdot b-(4-d)\,\slashed
a\slashed b$, \dots, in general $\gamma^\mu X\slashed a\gamma_\mu=2\,\slashed aX-(\gamma^\mu X\gamma_\mu)\,\slashed a$; then
$\slashed p\slashed p=p^2$, the traces, and the Dirac equation in the spinor chains. An unevaluated
\fnref{MinkowskiLoopIntegral} is simplified inside (one in a spinor chain or a trace is first taken out of it: the
spinors do not depend on the loop momentum, which is renamed if a spinor has a momentum of the same name, and the trace is
linear; an integrand that is 0 gives 0): its integrand is expanded and its indices are counted term by term,
and in the term outside it counts with its free indices. It gives a message and \pkg{\$Failed} for \pkg{GammaMatrix[5]},
explicit components, a product of matrices with \pkg{Times} and an index that appears more than twice in a term (also in
a term of the integrand of a loop integral); a momentum that is not a linear combination of momenta with coefficients that
are numbers or polynomials in the Feynman parameters is left as it is (\pkg{DiracSimplify::momentum}). A contraction, a trace, a spinor chain, and four-vectors:
\excode{qed_contraction}
\exformula{qed_contraction}
\excode{qed_trace}
\exformula{qed_trace}
\excode{qed_spinor_chain}
\exformula{qed_spinor_chain}
\excode{qed_vectors}
\exformula{qed_vectors}

\begin{function}{DiracOrder}
\syntax{DiracOrder[expr]}
\end{function}
\fnref{DiracSimplify}\pkg{[expr]} with the factors of every product put in a canonical order by $\{a,b\}=2a\cdot b$:
two expressions are equal when \pkg{DiracOrder} of their difference is 0. The contraction
$\gamma^\mu\slashed a\slashed b\gamma_\mu$, which \fnref{DiracSimplify} leaves as $2\slashed b\slashed a+(d-2)\slashed
a\slashed b$:
\excode{qed_order}
\exformula{qed_order}

\subsection{The diagrams in momentum space}

The theory, with the conventions of section~\ref{sec:qed-conventions}, and its 1PI functions at one loop. They are the
derivatives of the 1PI action $\Gamma[\bar\Psi,\Psi,\bar A]$ at zero means ($\delta^2\Gamma/\delta\bar\Psi\,\delta\Psi$,
$\delta^2\Gamma/\delta\bar A\,\delta\bar A$, $\delta^3\Gamma/\delta\bar\Psi\,\delta\Psi\,\delta\bar A$); \fnref{EffectiveAction1PI} takes one
field only, so the package takes the other side of the Legendre transform, the
one-particle-irreducible parts of the connected functions with their legs removed: the self-energy $-i\Sigma$, of order $e^2$ (in the formulas the electron
line is written $S$, the photon $D$ and the vertex $V$),
\excode{qed_theory}
\exboth{qed_theory}
the vacuum polarization $i\Pi^{\mu\nu}$, whose closed fermion loop gives it the sign $-1$, and the vertex correction
$ie\,\delta\Gamma^\mu$, of order $e^3$:
\excode{qed_polarization}
\exboth{qed_polarization}
\excode{qed_vertex}
\exboth{qed_vertex}
The tadpoles, a closed loop at the end of a photon line, are not one-particle irreducible (and they vanish, by Furry's
theorem).

\begin{function}{QEDIntegrand}
\syntax{QEDIntegrand[diagrams, \{x1 -> p1, \dots\}, rules, options]}
\end{function}
The momentum-space form of diagrams of the package with spinor and Lorentz indices, in condensed notation. The rules
are \pkg{S -> \{"Fermion", m\}}, the electron line \pkg{S[x, y]} $=\langle T\psi(x)\bar\psi(y)\rangle$, which is
$i(\slashed q+m)/(q^2-m^2)$ with $q$ flowing from $y$ to $x$; \pkg{Ph -> \{"Photon", xi\}}, the photon \pkg{Ph[z, w]} of
section~\ref{sec:qed-conventions}, or \pkg{Ph -> \{"Photon", \{xi, lambda\}\}} for a photon mass $\lambda$; and
\pkg{V~->~\{"Vertex",~c\}}, the local vertex \pkg{V[x, y, z]} $=c\,\gamma^\mu$, with the spinor indices of $x$ and $y$ and the
Lorentz index $\mu$ of $z$. The momentum $p_i$ enters at the external label $x_i$, and the momenta must add up to zero.
The momenta are routed as by \fnref{DiagramLoopIntegral}, the local vertices identifying their labels; the spinor
indices are followed along the fermion lines: an open line is a \fnref{DiracProduct}, from its $\psi$ end to its
$\bar\psi$ end, and a closed loop a \fnref{DiracTrace}, whose minus sign is in the coefficient computed by the package. Each term is
its coefficient times \fnref{MinkowskiLoopIntegral}\pkg{[}numerator times \fnref{FeynmanDenominator} factors\pkg{, \{k\}]}; a
tree has no integral. Options: \opt{LoopMomentum} \pkg{-> k}, the loop momentum (default \pkg{LoopMomentum[1]}, written $\ell_1$ by
\fnref{ExpressionToLaTeX}), and
\opt{Indices} \pkg{-> \{x -> mu, \dots\}}, the Lorentz index of a label (by default the label itself: the internal photon
of the self-energy below carries the indices $z_3$, $z_4$ of its ends). At most one loop and one open fermion line.
Messages, with the result \pkg{\$Failed}: \pkg{QEDIntegrand::rule}, a rule that is not one of these; \pkg{::external}, the
external momenta not given as \pkg{label -> momentum}; \pkg{::factor}, \pkg{::grassmann}, \pkg{::conservation}, as for
\fnref{DiagramLoopIntegral}; \pkg{::loops}, more than one loop; \pkg{::spinor}, spinor indices that do not form fermion
lines, or more than one open line (the four-fermion amplitudes, M{\o}ller and Bhabha scattering).

The tree-level vertex, $ie\gamma^\mu$, the self-energy in a general gauge, and the vacuum polarization:
\excode{qed_tree_vertex}
\exformula{qed_tree_vertex}
\excode{qed_sigma_integrand}
\exformula{qed_sigma_integrand}
\excode{qed_polarization_integrand}
\exformula{qed_polarization_integrand}

\begin{function}{FeynmanDenominator}
\syntax{FeynmanDenominator[q, m]}
\end{function}
The Minkowski denominator $1/(q^2-m^2+i0)$ of a line of momentum $q$ and mass $m$.

\begin{function}{MinkowskiLoopIntegral}
\syntax{MinkowskiLoopIntegral[f, \{k\}]}
\end{function}
The inert one-loop integral $\int_kf$ over the loop momentum $k$, with the measure of
section~\ref{sec:qed-conventions} in Minkowski space, $d=4-2\varepsilon$; written $\int_k$. The integrand $f$ is a product of
\fnref{FeynmanDenominator} factors, with integer powers (a negative power is a numerator), and of a numerator polynomial
in $k$: slashed momenta, four-vectors and scalar products, also in Dirac products, traces and spinor chains.
\fnref{DiracSimplify} simplifies its integrand, \fnref{IntegrateOneLoop} evaluates it.

\subsection{One-loop integrals}

\begin{function}{IntegrateOneLoop}
\syntax{IntegrateOneLoop[expr, eps]}
\end{function}
Evaluates the \fnref{MinkowskiLoopIntegral} objects of \pkg{expr} in $d=4-2\varepsilon$, exactly in $d$: Feynman parameters
\fnref{FeynmanParameter}\pkg{[i]} for any number of denominators and any integer powers,
\[
\prod_i\frac{1}{D_i^{a_i}}=\frac{\Gamma(A)}{\prod_i\Gamma(a_i)}\int\!dx_1\cdots dx_n\,\delta\Bigl(1-\sum_ix_i\Bigr)
\frac{\prod_ix_i^{a_i-1}}{\bigl(\sum_ix_iD_i\bigr)^A},\qquad A=\sum_ia_i,
\]
in the order of the sorted denominators $D_i=(k+r_i)^2-m_i^2$; the shift $k=l-\sum_ix_ir_i$, which leaves $(l^2-\Delta)^A$;
the tensor reduction, with the sum over all the pairings of the indices (odd ranks vanish),
\[
l^{\mu_1}\cdots l^{\mu_{2n}}\to\frac{(l^2)^n\,(g^{\mu_1\mu_2}\cdots g^{\mu_{2n-1}\mu_{2n}}+\dots)}{d(d+2)\cdots(d+2n-2)};
\]
and the master formula of section~\ref{sec:qed-conventions}, with $\Delta-i0$. The result is a sum of Dirac and tensor structures, each times an exact
integral over the Feynman parameters: none for one denominator, \fnref{ParameterIntegral} for two ($x_2=1-x_1$ is
used), \fnref{SimplexIntegral} for three or more. Messages, with the integral left unevaluated:
\pkg{IntegrateOneLoop::loops}, more than one loop momentum; \pkg{::momentum}, a denominator whose momentum is not $\pm
k$ plus external momenta; \pkg{::power}, a non-integer power of a denominator; \pkg{::numerator}, a numerator that is not
polynomial in $k$. Without a denominator that depends on $k$ the integral is 0 (scaleless). A tadpole, a tensor integral,
whose metric comes from the reduction ($\Gamma(3-\varepsilon)/((4-2\varepsilon)\Gamma(2-\varepsilon))=\frac12$ is simplified by
\pkg{FullSimplify}), and the bubble:
\excode{qed_tadpole}
\exformula{qed_tadpole}
\excode{qed_tensor}
\exformula{qed_tensor}
\excode{qed_bubble}
\exformula{qed_bubble}

\begin{function}{ParameterIntegral}
\syntax{ParameterIntegral[f, x]}
\end{function}
The inert integral $\int_0^1dx\,f$, exact in $d$ ($f$ may contain $\varepsilon$). \fnref{DifferentiateUnderIntegral}
differentiates under the integral sign, \fnref{ReduceParameterIntegrals} does it in closed form where possible, and
\pkg{N} evaluates it by \pkg{NIntegrate} (\pkg{N[expr, n]} to $n$ digits), with $\Delta-i0$: the logarithms and the
non-integer powers of a base that is negative somewhere in the interval, above a threshold, are $\log|b|-i\pi$ and
$|b|^ae^{-i\pi a}$ there, and the interval is split at the real roots of the bases. \pkg{N} does not evaluate a power $\le-1$ of a
base with a root inside the interval (with $\Delta-i0$ the integral is finite, but \pkg{NIntegrate} cannot integrate the
pole): \pkg{ParameterIntegral::threshold} and \pkg{Indeterminate}. The variable is
\pkg{FeynmanParameter[1]} in the results of \fnref{IntegrateOneLoop}, but it can be any symbol.

\begin{function}{SimplexIntegral}
\syntax{SimplexIntegral[f, \{x1, \dots, xn\}]}
\end{function}
The inert integral $\int_0^1dx_1\cdots dx_n\,\delta(1-x_1-\dots-x_n)\,f$ over the Feynman parameters of three or more
denominators. \fnref{DifferentiateUnderIntegral}, \fnref{ReduceParameterIntegrals} and \pkg{N} apply to it as to a
\fnref{ParameterIntegral}. \pkg{N} integrates over the unit cube, $x_1=y_1$, $x_2=(1-y_1)y_2$, \dots, with $\Delta-i0$; for a
power $\le-1$ of a base that changes sign in the region (finite with $\Delta-i0$, but a pole that \pkg{NIntegrate} does
not integrate) it gives the message \pkg{SimplexIntegral::threshold} and \pkg{Indeterminate}.
The triangle with a massless line, on shell at zero momentum transfer, $p_1^2=p_2^2=p_1\cdot p_2=m^2$, in units of $m$:
\excode{qed_triangle}
\exformula{qed_triangle}

\subsection{Feynman-parameter integrals}

\begin{function}{ReduceParameterIntegrals}
\syntax{ReduceParameterIntegrals[expr]}
\end{function}
Does the Feynman-parameter integrals of \pkg{expr} where possible, exactly in $d$: a \fnref{SimplexIntegral} whose
non-polynomial factors depend on one parameter only, after the delta function is used, and polynomially on the others
(the vertex at $q^2=0$), is integrated over the others; a \fnref{ParameterIntegral} of powers of monomials
$c\,x^a(1-x)^b$, with exponents that do not depend on $x$ (on shell), is a product of Gamma functions, so that infrared divergences appear as poles in $\varepsilon$:
the dimensional regularization of the infrared. A power $c^a$ of a coefficient is taken at $c-i0$ when $c$ is a negative
number; for a symbolic $c$, $c^a$ is the principal power, which is right for $c>0$, below thresholds. The other integrals
stay, and those with the same variables and the same factor in front are combined (\fnref{CombineParameterIntegrals}).
\excode{qed_reduce}
\exformula{qed_reduce}

\begin{function}{CombineParameterIntegrals}
\syntax{CombineParameterIntegrals[expr]}
\end{function}
Writes the Feynman-parameter integrals of \pkg{expr} that have the same variables and the same factor in front as one
integral of the sum of their integrands. The factor in front is the Dirac or tensor structure and the factors that do not
appear in the integrand, such as the coupling; the numbers and the other factors go into the integrand.
\excode{qed_combine}
\exformula{qed_combine}

\begin{function}{ExpandParameterIntegrals}
\syntax{ExpandParameterIntegrals[expr, eps, n]}
\end{function}
The Laurent expansion of \pkg{expr} in $\varepsilon$ up to $\varepsilon^n$ included: \fnref{ReduceParameterIntegrals}, then
\fnref{EpsilonExpand}, with the integrands of the remaining Feynman-parameter integrals expanded under the integral sign
to the order that the factor in front of each needs; the polynomial coefficients are integrated exactly, the others
stay as \fnref{ParameterIntegral} and \fnref{SimplexIntegral} objects, which \pkg{N} evaluates. The expansion under the
integral sign needs coefficients that are integrable at $\varepsilon=0$. An integrand that is not --- an infrared divergence
at a face of the region, where a base vanishes --- has its singular part subtracted, exactly in $d$: Beta functions
times integrals over the face, which carry the infrared poles, and an integrable remainder. This is done when there is
one singular face (in one dimension, each endpoint) and every factor of the singular terms that depends on the parameters
is a power of a polynomial (as in the results of \fnref{IntegrateOneLoop}), the exponents of the vanishing ones linear in
$\varepsilon$; otherwise \pkg{ExpandParameterIntegrals::ir} and \pkg{\$Failed} (\pkg{ExpandParameterIntegrals::failed} for an
integrand that is not expanded). Other factors (logarithms, Gamma functions of the parameters) are assumed regular and
are not checked. A list is expanded element by element. The triangle above, whose infrared pole
comes from the corner $x_1=1$, where $\Delta=(x_2+x_3)^2$ vanishes:
\excode{qed_triangle_expand}
\exformula{qed_triangle_expand}

\begin{function}{DifferentiateUnderIntegral}
\syntax{DifferentiateUnderIntegral[expr, v]}
\end{function}
The derivative of \pkg{expr} with respect to $v$, with the integrands of its \fnref{ParameterIntegral} and
\fnref{SimplexIntegral} objects differentiated under the integral sign (\pkg{D} itself would apply the chain rule to the
inert heads); for instance $d\Sigma/dp^2$, for the wave-function renormalization (section~\ref{sec:qed-results}).

\begin{function}{VertexFormFactors}
\syntax{VertexFormFactors[vertex, mu, \{p1, p2, m\}]}
\end{function}
The form factors \pkg{\{F1, F2, F3\}} of a vertex $\bar u(p_2)\Gamma^\mu u(p_1)$ between spinors of mass $m$
(section~\ref{sec:qed-conventions}): \pkg{vertex} is a sum of \pkg{SpinorChain[\{p2, m\}, \dots, \{p1, m\}]} terms, which
\fnref{DiracSimplify} reduces to $\gamma^\mu A+p_1^\mu B_1+p_2^\mu B_2$ between the spinors; by the Gordon identity
$\bar u(p_2)(p_1+p_2)^\mu u(p_1)=\bar u(p_2)(2m\gamma^\mu-i\sigma^{\mu\nu}q_\nu)u(p_1)$,
$F_1=A+m(B_1+B_2)$, $F_2=-m(B_1+B_2)$ and $F_3=(B_2-B_1)/2$. A part of another form gives \pkg{VertexFormFactors::rest} and
\pkg{\$Failed}.

\subsection{The electron self-energy, the vacuum polarization and the vertex}\label{sec:qed-results}

\paragraph{Self-energy} The integral of the self-energy above, $\Sigma=i\times(-i\Sigma)$, has the pole, in units of
$e^2/(16\pi^2\varepsilon)=\alpha/(4\pi\varepsilon)$,
\excode{qed_sigma_pole}
\exformula{qed_sigma_pole}
that is $\Sigma_{\rm pole}=\frac{\alpha}{4\pi\varepsilon}\bigl(m(3+\xi)-\xi\slashed p\bigr)$ (the Feynman gauge of
\cite[section~10.3]{PeskinSchroeder} is $\xi=1$). The $\overline{\rm MS}$ counterterms $\delta_2$, gauge dependent and 0 in the
Landau gauge, and $Z_m-1=\delta_m/m-\delta_2=-3\alpha/(4\pi\varepsilon)$, which is not:
\excode{qed_counterterms}
\extext{qed_counterterms}
In the Feynman gauge, to $\varepsilon^0$, the finite parts are integrals over the Feynman parameter of the photon line, with
$\Delta=(1-x_1)(m^2-x_1p^2)$:
\excode{qed_sigma_finite}
\exformula{qed_sigma_finite}
On shell, with $\Sigma=m\,A+\slashed p\,B$, the mass shift is $m-m_0=\Sigma(\slashed p=m)=m(A+B)$ and
$\delta Z_2=d\Sigma/d\slashed p=B+2m^2(A'+B')$ at $p^2=m^2$~\cite[section~7.1]{PeskinSchroeder}; the derivative is taken under
the integral sign, and the infrared divergence on shell is regulated dimensionally (here the poles are ultraviolet and
infrared). In units of $\alpha/(4\pi)$, in every gauge:
\excode{qed_delta_z2}
\exformula{qed_delta_z2}
that is $\delta Z_2=-\frac{\alpha}{4\pi}\bigl(\frac3\varepsilon+4+3\log\frac{\mu^2}{m^2}\bigr)$, and $(m-m_0)/m=-\delta Z_2$. With a photon
mass $\lambda$ (the rule \pkg{PhotonPropagator -> \{"Photon", \{1, lambda\}\}}) the infrared divergence becomes
$\log\lambda^2$:
$\delta Z_2=-\frac{\alpha}{4\pi}\bigl(\frac1\varepsilon+4+\log\frac{\mu^2}{m^2}+2\log\frac{\lambda^2}{m^2}\bigr)$, up to terms that vanish
with $\lambda$.

\paragraph{Vacuum polarization} $\Pi^{\mu\nu}=g^{\mu\nu}\Pi_g+q^\mu q^\nu\Pi_{qq}$ does not depend on $\xi$ (there is no internal
photon), and it is transverse, $\Pi_g+q^2\Pi_{qq}=0$, exactly in $d$: the integrand of $\Pi_g+q^2\Pi_{qq}$ vanishes
identically, as \pkg{FullSimplify} finds.
\excode{qed_transverse}
\exformula{qed_transverse}
With $\Pi^{\mu\nu}=(q^2g^{\mu\nu}-q^\mu q^\nu)\,\Pi(q^2)$, the pole of $\Pi$ is
$\delta_3=-\alpha/(3\pi\varepsilon)$~\cite[section~10.3]{PeskinSchroeder},
\excode{qed_delta3}
\exformula{qed_delta3}
and, with $Z_1=Z_2$ (checked below), $e_0=\mu^\varepsilon Z_3^{-1/2}e$ gives the beta function, $\beta(e)=e^3/(12\pi^2)$, by
\fnref{BetaFunctionFromZ} with \opt{ScalePower} \pkg{-> 1}:
\excode{qed_beta}
\exformula{qed_beta}
The renormalized $\hat\Pi(q^2)=\Pi(q^2)-\Pi(0)$ of~\cite[section~7.5]{PeskinSchroeder} is finite and independent of $\mu$ (the two $\log\mu^2$ cancel, since
$\int_0^1x(1-x)\,dx=\frac16$),
\excode{qed_pi_hat}
\exformula{qed_pi_hat}
and \pkg{N} evaluates it, below the threshold $q^2=4m^2$ and above it, where the prescription $\Delta-i0$ gives the
imaginary part of $\hat\Pi(q^2+i0)$ (here at $q^2=-3m^2$ and $5m^2$, with $m=\mu=1$):
\excode{qed_pi_hat_values}
\extext{qed_pi_hat_values}

\paragraph{Vertex} The integrand of $ie\,\delta\Gamma^\mu(p_2,p_1)$, with the photon momentum $q=p_2-p_1$ entering at $x_3$:
\excode{qed_vertex_integrand}
\exformula{qed_vertex_integrand}
Between the spinors, with $p_1\cdot p_2=m^2-q^2/2$ (\pkg{q2} is $q^2$), \fnref{VertexFormFactors} gives the three form factors,
as Feynman-parameter integrals in a general gauge. $F_2$ is finite in the infrared, and at $q^2=0$ it is Schwinger's
$F_2(0)=e^2/(8\pi^2)=\alpha/(2\pi)$~\cite{Schwinger1948}, in every gauge: the anomalous magnetic moment $(g-2)/2$ of the
electron at one loop.
\excode{qed_form_factors}
\exformula{qed_form_factors}
$F_3$ vanishes, as the Ward identity requires (at $q^2=0$ in closed form, at $q^2\ne0$ its integrand is odd under the
exchange of the parameters of the two electron lines), and $F_1(0)-1=-\delta Z_2$, the infrared divergences again
regulated dimensionally, so that $F_1(0)=1$ after renormalization, exactly~\cite{Ward1950}:
\excode{qed_ward_identity}
\extext{qed_ward_identity}
With a photon mass, which regulates the infrared, the pole of $F_1(0)$ is the ultraviolet one, $\xi\alpha/(4\pi\varepsilon)=-\delta_2$:
$Z_1=Z_2$ in $\overline{\rm MS}$. The test suite checks these results, and also $F_2(-m^2)$ and $\hat\Pi(q^2)$ at $q^2=-3m^2$ and
$5m^2$ against the integrals of~\cite[sections~6.3 and~7.5]{PeskinSchroeder}, the Ward identity with a photon mass,
$F_1(-m^2)$ with its infrared pole in two gauges, the Ward identity of the Compton amplitude at tree level, and the
messages.

\subsection{Limits, and what would come next}

\begin{itemize}
\item \textbf{One loop.} \fnref{QEDIntegrand} and \fnref{IntegrateOneLoop} take a single loop momentum; diagrams and
  integrals with more loops give the messages \pkg{QEDIntegrand::loops} and \pkg{IntegrateOneLoop::loops}.
\item \textbf{No $\gamma_5$}: no traces with $\gamma_5$, no anomalies, no weak interactions (\pkg{GammaMatrix[5]} is refused).
\item \textbf{At most one open fermion line} per term: no four-fermion amplitudes (M{\o}ller and Bhabha scattering at tree
  level or at one loop, \pkg{QEDIntegrand::spinor}).
\item \textbf{Closed forms} only where $\Delta$ is a monomial in the Feynman parameters (on shell, and the vertex at $q^2=0$);
  otherwise the integrals stay inert and \pkg{N} evaluates them: $\Pi(q^2)$, $\Sigma(p^2)$ off shell and $F_{1,2}(q^2)$ have no
  closed form yet. Symbolic closed forms use the principal branch: $c^a$ and the logarithms that \fnref{EpsilonExpand} writes
  outside an integral are right below thresholds, for positive coefficients. Above a threshold, only \pkg{N} applies the
  prescription $\Delta-i0$; a power $\le-1$ of a base that vanishes inside the region (the vertex at timelike $q^2>4m^2$, finite and
  complex there) is not evaluated by \pkg{N} (\pkg{::threshold}, \pkg{Indeterminate}).
\item \textbf{Infrared divergences} under the integral sign are subtracted at one face of the simplex only
  (\pkg{ExpandParameterIntegrals::ir} otherwise).
\end{itemize}
The natural next steps: the closed forms of the one-dimensional integrals of $\Pi(q^2)$ and $\Sigma(p^2)$ (logarithms and the
square root of the K\"all\'en function, as for the massive bubble of section~\ref{sec:dimreg}), with their continuation above
threshold; $\gamma_5$, naively anticommuting or in the scheme of 't~Hooft and Veltman~\cite{tHooftVeltman1972}; several open
fermion lines, for the boxes of M{\o}ller and Bhabha scattering; and two loops, by the nested integration of
section~\ref{sec:dimreg} with tensor numerators.

% ------------------------------------------------------------------ 19
\section{Algorithms and performance}\label{sec:algorithms}

\paragraph{Representation} Internally a term is \pkg{diagramTerm[c, odd, dummies]}: the commuting part, the list of
Grassmann-odd factors in order, and the integration variables. An expression is a list of terms; \pkg{Plus},
\pkg{Times}, \pkg{**}, powers and integrals are expanded into such lists, the integration variables being renamed with
fresh symbols.

\paragraph{Differentiation} The derivative of a term is the sum over its commuting factors (product rule, with the chain
rule inside them) and over its odd factors (with the Grassmann sign). Derivatives of source atoms give delta functions,
integrated at once against the integration variables; delta functions between external labels are normalized so that
every group of labels joined by delta functions is represented by its first label.

\paragraph{Canonical form} A canonical renaming of the integration variables is computed by colour refinement: every
variable starts with the same colour, and is then characterized by the multiset of factors it appears in, the other
variables being replaced by their colours, iteratively until the partition is stable. Ties are broken by
individualization, branching over the members of the first non-trivial colour class, as in graph-isomorphism
algorithms~\cite{McKayPiperno2014}.
Before branching, the transpositions within the class are tested: if all of them map the term to itself (with sign
$+1$) a single branch suffices; if one maps it to minus itself, the term vanishes. Among the renamings found, the one
giving the smallest expression in Mathematica's canonical order is chosen. Nested integrals are canonicalised
recursively, with \pkg{DummyVariable[depth, k]}.

\paragraph{Wick's theorem} The pairings are enumerated recursively, from the first unpaired field. When the partner of a
field belongs to an interaction vertex not yet touched, only the first such vertex of each kind and one field per orbit of
the vertex's symmetries (permutations of its integration variables that leave it unchanged, and identical commuting
fields) are tried, with the corresponding multiplicity; identical commuting fields of a vertex already touched are
interchangeable. The connectivity of the partial pairing is tracked with a union--find structure, and components that
close without external points are discarded when only diagrams without vacuum bubbles, or connected ones, are wanted.
Every contraction $\langle f_if_j\rangle_0$ is computed once from $Z_0$, by differentiation, and cached.

\paragraph{Cuts} Irreducibility, skeletons, channels and loops are decided on the graph of section~\ref{sec:graph}, by
trying all the cuts of one or two lines.

\paragraph{Frequency sums} Momenta and frequencies are routed along a spanning tree of the diagram, chosen with the
instantaneous and boson lines first, so that every loop runs along fermion lines; the points joined by instantaneous
lines form the time vertices, and the fermion lines inside a time vertex are at equal times. The poles of a summand are
read from the linear factors of its denominators (a boson propagator is split into two), and its residues, of any
order, are Laurent coefficients computed by differentiation; whether the sum converges, and whether it depends on the
convergence factor, is read from the degrees in $z$ of its terms. After each sum the statistical factors are freed of
the other frequencies by periodicity. At the end the terms are put over common denominators, in partial fractions in the
external frequency (part by part, between the lines that carry it alone), and the statistical factors are written as a
sum over processes when this is shorter. The cutting rules are derived independently, in imaginary time: the
propagating lines between the two external time vertices are grouped into chains, the product of the propagators of
each chain is expanded in partial fractions of its frequency, and each pole puts the chain on shell with its factor of
emission or absorption (section~\ref{sec:cutting}).

\paragraph{Dimensional regularization} The momenta are routed as for the Matsubara sums, every line becoming a
\pkg{ScalarPropagator}. The integrand of a loop integral is expanded into products of propagators, and in each product
the loops are integrated one at a time, by a loop momentum that appears in at most two different propagators (after a
change of variables, if needed); each step is a closed form in $d$ (tadpoles, G-functions with a new massless line of
non-integer power, bubbles), and each product records what it met: scaleless integrals, power or logarithmically
divergent, infrared divergences, irreducible integrals, numerators. These records give the messages of
\fnref{IntegrateLoops} and the choice of a routing in \fnref{InfraredRearrangedIntegral}. The expansion in $\varepsilon$ uses
truncated Laurent series of its own, $\{v,\{c_0,c_1,\dots\}\}$ for $\sum_ic_i\varepsilon^{v+i}$, multiplied, inverted and
exponentiated term by term, the order of every factor being chosen from the poles of the others, so that the result is
exact to the order asked without \pkg{Series}; the Gamma functions are expanded through their logarithmic derivatives,
and the logarithms of a massive bubble are integrated over its Feynman parameter in closed form.

\paragraph{QED at one loop} \fnref{QEDIntegrand} routes the momenta as \fnref{DiagramLoopIntegral} does, with the local
vertices replaced by delta functions between their labels, and builds the fermion lines from pieces \{row, column,
matrix\}, one for every electron line and vertex: an open line is followed from the label that is not a column, a
closed loop from any of its pieces. The Dirac algebra rewrites products by linearity and contracts an index by
$\gamma^\mu X\slashed a\gamma_\mu=2\slashed aX-(\gamma^\mu X\gamma_\mu)\slashed a$; traces are expanded recursively along their
first element and memoized during one call; in a spinor chain $\slashed p{}_1$ is moved to the right and $\slashed p{}_2$ to
the left by anticommutation. Compound momenta, the loop momentum and the Feynman parameters are replaced by symbols
during the algebra. \fnref{IntegrateOneLoop} groups the terms of a numerator by their denominators, introduces the Feynman
parameters, shifts the loop momentum, reduces the tensor numerators by the sum over the pairings of their indices and
applies the master formula term by term; the coefficients of each Dirac and tensor structure are collected into one
Feynman-parameter integral. Before \fnref{ExpandParameterIntegrals} expands under the integral sign, the behaviour of the
integrand near every face of the simplex is read from the orders of its bases there; a singular face is subtracted with
the leading terms of the integrand in the distance from it, integrated exactly (Beta functions times an integral over
the face). \pkg{N} maps the simplex to the unit cube, takes $\Delta-i0$ for the bases that change sign, and splits a
one-dimensional interval at their real roots.

\paragraph{Performance} The showcase --- the electron gas and $\phi^4$ theory to third order ($G$: 87
diagrams, $\Sigma$, $\log Z$, $\Phi$, $\Sigma[G]$, $\Gamma^0$: 58, the full vertex $F$: 94, the parquet checks), the exact
equations, Hubbard--Stratonovich, the Matsubara sums, $\phi^4$ theory to two loops in dimensional regularization, QCD and
the trees of vertex functions up to six points, with several hundred pictures --- runs in about four minutes.
The QED computation of section~\ref{sec:qed-results} (\pkg{Tests/QEDOneLoop.wl}) takes about ten seconds.
The superfield costs more per order than the electron gas: its Luttinger--Ward functional takes about 3~s at third
order and 2~minutes at fourth, \fnref{EffectiveAction2PI} with its mean about 3~minutes at fourth order (16~s for the
fermionic $\Gamma[\bar\Psi,\Psi,G]$), and \fnref{SuperfieldComponents} of that action about 8~minutes.

% ------------------------------------------------------------------ 20
\section{Compatibility, tests and troubleshooting}\label{sec:troubleshooting}

\paragraph{Compatibility} The package runs in Mathematica~12.1 or later (for the paclet manager that loads it). The
test suite is \pkg{Tests/RunTests.wl}:
\begin{tcblisting}{listing only, colback=codebg, colframe=codeframe, boxrule=0.5pt, arc=2pt, left=6pt, right=6pt,
  top=3pt, bottom=3pt, listing options={style=plainoutput}}
wolframscript -file Tests/RunTests.wl
wolfram -script Tests/RunTests.wl         # the same, without wolframscript
\end{tcblisting}
It prints PASS or FAIL for every check and a summary with the number of checks passed, about 800 in two minutes;
\pkg{Tests/RunTests.out} is the output of the last run. The checks cover the functional derivative (sources, chain rule,
renaming of integration variables, delta functions), Wick's theorem for bosons and fermions (signs, numbers of pairings,
spin indices), the numbers of diagrams of $\phi^4$ theory and of the electron gas (vacuum, Green function, 1PI,
skeleton), the Luttinger--Ward functional and its derivatives, the full vertex and its channels, the parquet equations
for fermions and bosons, the tree expansion, the closed loops, the pictures (arrows, points, layout rules) and the
\LaTeX{} output, and the Matsubara sums: closed forms, sums over many frequencies (the evaluated diagrams against the
brute-force sums of their integrands, with poles of first to third order), the Lindhard function at finite and zero
temperature, the second-order self-energy and the self-energy of electrons and phonons in the form of the textbooks, the
agreement of the residues with the cutting rules, and the messages and \pkg{\$Failed} of divergent sums and invalid
arguments; and dimensional regularization: the loop integrals of the diagrams, tadpoles, bubbles and G-functions against
numerical integrals in non-integer dimensions, the Laurent expansion against \pkg{Series}, the massive
bubbles against their Feynman-parameter integrals, the massless sunset against its computation in position space, the
two-loop four-point diagrams and their routings, the messages of infrared and scaleless divergences, and the
$\overline{\rm MS}$ renormalization of $\phi^4$ theory to two loops ($Z$ factors, $\beta$, $\gamma_\phi$, $\gamma_{m^2}$, the
Wilson--Fisher exponents); and QED at one loop: the Dirac algebra, the traces and the spinor chains against explicit
Dirac matrices in $d=4$ and $d=6$ and explicit spinors, the one-loop integrals against direct numerical integration and
against independent references, the infrared subtraction and the prescription $\Delta-i0$ above threshold, the refused
input, and the self-energy, the vacuum polarization and the vertex against Peskin and Schroeder (section~\ref{sec:qed-results});
and the Hubbard--Stratonovich transformation and the integration of a Gaussian field; and, from
\pkg{Tests/EffectiveActions.wl} and \pkg{Tests/Superfields.wl}, the algebra of operators and their traces, the 2PI and 1PI
effective actions (their pieces, their stationary points and values, their derivatives with respect to the propagator, the
means and a one-body source: $\Sigma$, $\Gamma^0$, $\Gamma_{\rm irr}=\delta^4\Gamma/\delta\bar\phi^4$, the full vertex, the
responses and their static and dynamic limits, conserving approximations), Potthoff's functionals, the zero-dimensional
values (one site, scalars and arrays, against exact integrals), the grand potential in momentum and frequency, and the
superfields (the channels, the parquet equations, \fnref{SuperfieldComponents}, the functional of Nambu and Gor'kov and its
gap equation). The
showcase checks, for the electron gas and for $\phi^4$ theory to third order, the Dyson equation, the skeleton expansion of
$\Sigma[G]$, the homogeneity relations $\int\Sigma^{(m)}G=2m\,\Phi^{(m)}$ and $\int\Gamma^{0(m)}G=(2m-1)\,\Sigma^{(m)}$ (for real
bosons, with $\Sigma=-2\,\delta\Phi/\delta G$ and $\Gamma^0=-2\,\delta\Sigma/\delta G$, $\int\Sigma^{(m)}G=-4m\,\Phi^{(m)}$ and
$\int\Gamma^{0(m)}(x_1,x_2,x_3,x_4)\,G(x_3,x_4)=-2(2m-1)\,\Sigma^{(m)}(x_1,x_2)$), the
symmetry of $\Gamma^0$, the Ward identity, the Bethe--Salpeter equation, the agreement of $\Gamma^0=\delta^2\Phi/\delta G\delta
G$ with the particle--hole irreducible part of $F$, the three Bethe--Salpeter equations of the parquet decomposition,
the Schwinger--Dyson equation, the equations of the vertex of the density and Hedin's equations, and the
Hubbard--Stratonovich theory against the original one; and the tree expansion up to six points against the
zero-dimensional theory.

\paragraph{Messages}
\begin{itemize}
\item \pkg{FunctionalDiagrams::builtin}: a declared name is a built-in symbol of the Wolfram Language; use another name.
\item \pkg{FunctionalDiagrams::clash}: public names of the package are symbols of the context \pkg{System`} (built-in
  functions of a later version of Mathematica), which are used in place of the package's functions; the definitions of
  the package for these names are not made. A name also defined by another package gives \pkg{::shdw} instead
  (section~\ref{sec:install}).
\item \pkg{FunctionalDerivative::times}: a product of two Grassmann-odd factors was written with \pkg{Times}, which
  loses their order and the sign; the evaluation is aborted. Use \pkg{**}.
\item \pkg{FunctionalDiagrams::label}: an external label does not occur in the expression --- a misprint, or the
  arguments in the order of version 1 (line heads before labels).
\item \pkg{FunctionalDiagrams::labelsymbol}: a label is not a symbol (a number, for instance).
\item \pkg{FunctionalDiagrams::optx}: an unknown option of a drawing function or of a function of
  section~\ref{sec:matsubara}.
\item \pkg{FunctionalDiagrams::realboson}: \pkg{"Statistics" -> "Boson"} with a propagator that is not declared
  symmetric (a complex field: use \pkg{"Statistics" -> "ComplexBoson"}).
\item \pkg{FunctionalDiagrams::statisticsdefault}: \fnref{SelfEnergy} or \fnref{IrreducibleVertex} on a propagator declared
  symmetric or antisymmetric, without \opt{Statistics}: the fermionic default is used; give \pkg{"Statistics" -> "Boson"}
  or \pkg{"Superfield"}.
\item \pkg{FunctionalDiagrams::vertexlabels}: a list of labels of the wrong length in \fnref{SelfEnergy},
  \fnref{IrreducibleVertex} or \fnref{PerturbativeFullVertex}; the result is \pkg{\$Failed}.
\item \pkg{FunctionalDiagrams::rulesymmetry}: a rule \pkg{G0 -> G} with \pkg{G0} symmetric (antisymmetric) and \pkg{G} not:
  declare \pkg{G} too, or its terms with the labels exchanged are not recognized as equal.
\item \pkg{FunctionalDiagrams::fieldarity}: a field with another number of arguments than its source in the generating
  functional (\pkg{psi[x]} when $Z_0$ holds \pkg{eta[x, s]}): its expectation values vanish.
\item \pkg{FunctionalDiagrams::opdotpower}: \pkg{A\^{}-1 . B}, which Mathematica reads as \pkg{A\^{}(-(1 . B))}; it is taken as
  \pkg{A\^{}(-1) . B}. \pkg{FunctionalDiagrams::opsource}: a source in an operator outside any kernel; the derivative of
  the trace is taken as 0. The messages \pkg{opvariation}, \pkg{opfunction}, \pkg{optimes} and \pkg{opscalar}:
  operators whose variation is not known (a non-integer power), or that cannot be told from their coefficients
  (section~\ref{sec:effective}).
\item \pkg{FunctionalDiagrams::inversekernel}: \fnref{DiagramPlot} or \fnref{DrawDiagram} of a term with a kernel
  \pkg{OperatorKernel[D\^{}-1]} and lines \pkg{D} (a derivative of a trace or of the 1PI action before \fnref{ExpandTrace}):
  the inverse kernel is drawn as a line, like \pkg{D}, and the products $D\,D^{-1}$ look like chains of lines; contract them
  with \fnref{ExpandTrace} first.
\item \pkg{SuperfieldComponents::trace}: a trace with no component form, ${\rm Tr}\ln$ of a Nambu matrix with anomalous
  blocks; it is kept as it is, beside the components of the other terms.
\item \pkg{FunctionalDerivative::supertrace}: a derivative, with respect to $G$ or an anomalous block, of a trace that
  \fnref{SuperfieldComponents} kept as it is; it is taken as 0. Differentiate the superfield functional with respect to
  $\mathcal G$ first, then take the components at zero anomalous blocks.
\item \pkg{FunctionalDiagrams::statistics}, \pkg{FunctionalDiagrams::method}: an unknown value of \opt{Statistics} or of
  \opt{Method}; the default is used.
\item \pkg{DysonSeries::zeroth}: with \opt{MaximumOrder}, a part of order 0 of the self-energy is left out.
\item \pkg{ConnectedFunctionFromVertices::labels}: fewer than two labels.
\item \pkg{FieldDerivative::nofield}: a field without a source (\fnref{DeclareField}).
\item \pkg{FieldDerivative::times}: the Grassmann fields of a composite operator were multiplied with \pkg{Times}, which
  loses their order; write them as separate entries of the list of fields, or with \pkg{**}. The result is \pkg{\$Failed}.
\item \pkg{FunctionalDerivative::nosource}: a derivative with respect to something that is not a declared source; it
  is 0.
\item \pkg{FunctionalDiagrams::nosources}: a generating functional without any declared source --- the sources were not
  declared, or a propagator named with a built-in symbol such as \pkg{D} evaluated to 0.
\item \pkg{FunctionalDiagrams::lineargs}: a line head with more than two arguments, for instance a spin-dependent
  propagator \pkg{G[x, s, y, t]}, which the cuts treat as a vertex (see \fnref{DeclareDiscreteIndexValues}).
\item \pkg{PerturbativeCorrelationFunction::order}: the order is not a non-negative integer, \pkg{\{k\}} or
  \pkg{\{kmin, kmax\}}.
\item \pkg{GaussianExpectationValue::method}: \opt{WickPairings} was asked for a generating functional that is not Gaussian;
  differentiation is used instead.
\item \pkg{DeclareLineStyle::style}: an unknown line style.
\item \pkg{MatsubaraSum::conv}: the summand decays like $1/z$ and no convergence factor was given; the symmetric sum is
  returned. \pkg{MatsubaraSum::axis}: a pole on the axis of the frequencies summed (an anomalous term).
  \pkg{MatsubaraSum::shift}: a shift that is not a Matsubara frequency. The following give \pkg{\$Failed}:
  \pkg{MatsubaraSum::diverge}, \pkg{::free}, a summand that does not decay, or that does not depend on the variable (the
  sum diverges); \pkg{MatsubaraSum::degree}, \pkg{::pole}, a denominator of degree more than two that does not factor, a
  factor that is not a power of a linear function; \pkg{MatsubaraSum::statistics}, \pkg{::factor}, \pkg{::variable}, an
  unknown statistics or convergence factor, a variable that is not a symbol.
\item \pkg{DeclareMatsubaraFrequency::symbol}: a frequency that is not a symbol.
\item Invalid input of \fnref{EvaluateMatsubara}, \fnref{DiagramIntegrand} and the cutting rules, with the result
  \pkg{\$Failed}: \pkg{EvaluateMatsubara::rule}, a rule that is not a line rule; \pkg{::factor}, a factor that is
  neither a line, a delta function nor a constant; \pkg{::grassmann}, a term with Grassmann sources;
  \pkg{::conservation}, external momenta and frequencies that do not add up to zero; \pkg{::frequency}, an external
  frequency without statistics; \pkg{::external}, an external point that is not \pkg{label -> \{momentum, I w\}};
  \pkg{::loopnames}, \pkg{::loops}, a value of \opt{LoopMomenta} that is not a list, or that gives fewer names than
  loops; \pkg{::convergence}, a value of \opt{ConvergenceFactor} other than 1, $-1$ and 0.
\item \pkg{EvaluateMatsubara::conv}: a frequency sum that depends on the convergence factor, with no line at equal times
  to give one; the symmetric sum is taken. \pkg{EvaluateMatsubara::diverge}: a frequency sum that diverges (a boson line
  that does not vanish at large frequency); the result is \pkg{\$Failed}.
\item \pkg{AnalyticContinuation::periodicity}: statistical factors that still depend on the frequency being continued.
\item \pkg{CuttingRuleImaginaryPart::class}: a diagram to which the cutting rules of section~\ref{sec:cutting} do not
  apply. \pkg{CuttingRuleImaginaryPart::frequency}: a frequency that does not enter the diagram at one external point and
  leave it at another. The result is \pkg{\$Failed}.
\item In dimensional regularization (section~\ref{sec:dimreg}), the messages of \fnref{EvaluateMatsubara} for invalid
  input, with the result \pkg{\$Failed}: \pkg{DiagramLoopIntegral::rule}, \pkg{::external}, \pkg{::factor},
  \pkg{::grassmann}, \pkg{::conservation}, \pkg{::loopnames} and \pkg{::loops}; \pkg{IntegrateLoops::irreducible} and \pkg{::numerator}, an integral that is
  left unevaluated; \pkg{IntegrateLoops::ir}, \pkg{::scaleless}, \pkg{::zero}, an infrared divergence, a logarithmically
  divergent scaleless integral set to 0, a massless line of zero momentum; \pkg{InfraredRearrangedIntegral::norouting}, a
  term without a valid routing (\pkg{\$Failed}); \pkg{EpsilonExpand::loop}, \pkg{::function}, \pkg{::bubble},
  \pkg{::variable}, \pkg{::pole}, an expression that is not expanded (\pkg{\$Failed}); \pkg{MassInsertion::label}, a
  label of the insertion that is not a new symbol; \pkg{BetaFunctionFromZ::pole}, \pkg{AnomalousDimensionFromZ::pole},
  poles of $Z$ that are not consistent (\pkg{\$Failed}).
\item In QED (section~\ref{sec:qed}): \pkg{GammaMatrix::gamma5}, \pkg{DiracSimplify::component}, \pkg{::product},
  \pkg{::index}, input of the Dirac algebra that is refused (\pkg{\$Failed}; in \fnref{IntegrateOneLoop} the integral is left
  unevaluated); \pkg{DiracSimplify::momentum}, a momentum that is not a linear combination of momenta with coefficients
  that are numbers or polynomials in the Feynman parameters, left as it is; \pkg{QEDIntegrand::rule}, \pkg{::external},
  \pkg{::factor}, \pkg{::grassmann}, \pkg{::conservation}, \pkg{::loops}, \pkg{::spinor}, invalid diagrams or rules
  (\pkg{\$Failed}); \pkg{IntegrateOneLoop::loops}, \pkg{::momentum}, \pkg{::numerator}, \pkg{::power}, an integral left
  unevaluated; \pkg{ExpandParameterIntegrals::ir}, \pkg{::failed}, an integral that is not expanded (\pkg{\$Failed});
  \pkg{ParameterIntegral::threshold} and \pkg{SimplexIntegral::threshold}, an integrand with a power $\le-1$ of a base
  that vanishes in the region, above a threshold, which \pkg{N} does not integrate (the result is \pkg{Indeterminate}); \pkg{VertexFormFactors::rest}, a vertex that is not of the form of the form factors
  (\pkg{\$Failed}).
\item Channels, exact equations and trees: \pkg{FunctionalDiagrams::channel}, a value of \opt{Channel} that does not exist
  for the statistics; \pkg{FunctionalDiagrams::ppkernel}, the particle--particle kernel of fermions asked as a derivative
  with respect to $G$ (it is one with respect to the anomalous propagators: use the superfield);
  \pkg{ParquetDecomposition::args}, wrong arguments; \pkg{FunctionalDiagrams::exactname}, \pkg{::exactobject}, an
  unknown equation or a missing object of \fnref{ExactEquation}; \pkg{ConnectedFunctionFromVertices::fermionlines},
  directions that are not \pkg{"Out"} and \pkg{"In"}, as many of each. The result is \pkg{\$Failed}.
\item Auxiliary fields (section~\ref{sec:auxiliary}): \pkg{FunctionalDiagrams::hsform}, a term of the interaction that is
  not of the form $\frac12\int K\,O\,O$ (it is kept); \pkg{::hsnone}, no such term; \pkg{::hshead}, a kernel that is not a
  number times a propagator; \pkg{::outsource}, no Gaussian factor of the field in the free functional; \pkg{::outlinear},
  a field that does not enter the interaction linearly (\pkg{\$Failed}).
\item Operators and zero dimensions (section~\ref{sec:effective}): \pkg{FunctionalDiagrams::opzero}, the inverse of a
  zero operator; \pkg{SuperfieldComponents::head}, a factor of the superfield without a rule;
  \pkg{ZeroDimensionalValue::indices}, free labels without an internal state, \pkg{::rank}, an array of the wrong rank,
  \pkg{::operator}, an operator without a matrix value, \pkg{::loopfactor}, \opt{LoopFactor} ignored with array values.
\item In the Matsubara evaluation: \pkg{EvaluateMatsubara::anomalous}, antisymmetric kernels (anomalous or superfield
  propagators), which cannot be evaluated in momentum and frequency yet: the expression must be written with the normal
  propagators (\fnref{SuperfieldComponents} at zero anomalous propagators); \pkg{EvaluateMatsubara::kernel}, kernels of an
  operator without line rules; \pkg{EvaluateMatsubara::tracegrows}, a trace whose summand grows at large frequency;
  \pkg{MatsubaraSum::log}, a logarithm that is not of the form $c\log r$ with $r$ rational in the frequency. The result is
  \pkg{\$Failed}.
\end{itemize}

\paragraph{Common pitfalls}
\begin{itemize}
\item Labels, integration variables and the heads of functions, sources and fields must be symbols, and not built-in
  ones: \pkg{C}, \pkg{D}, \pkg{E}, \pkg{I}, \pkg{K}, \pkg{N}, \pkg{O}, \pkg{Gamma}, \pkg{Beta}, \pkg{GreenFunction} are built in
  (\pkg{Context[name]} is \pkg{"System`"} for a built-in symbol).
\item Grassmann-odd factors must be multiplied with \pkg{**}, and their order matters.
\item The order of the declarations does not matter: a field may be declared before its source.
\item In a terminal session results print in output form, for instance \pkg{IntegralOver[\dots, \{z1, z2\}]} and
  \pkg{DeltaFunction[a, b]}; the compact forms, $\int_{z_1z_2}$ and $\delta_{ab}$, are shown in notebooks. In a script
  (\pkg{wolfram -script}) \pkg{Print} writes the input form, with the integration variables as \pkg{DummyVariable[1, 1]}:
  \pkg{Print[ToString[expr]]} writes \pkg{z1}.
\item In a script a statement ends at the end of a line where it is complete: \pkg{a = b} on one line and
  \pkg{+ c} on the next are two statements. Break long lines after a binary operator, not before it.
\item A product with an inverse is \pkg{G0\^{}(-1) . G}: \pkg{G0\^{}-1 . G} is read as \pkg{G0\^{}(-(1 . G))} (the package
  corrects it, with the message \pkg{opdotpower}).
\item Long computations can exceed the default \pkg{\$IterationLimit} of Mathematica (4096); set
  \pkg{\$IterationLimit = 10\^{}8} before them, as the scripts of the package do.
\item Results are in canonical form, with the integration variables renamed: compare expressions with
  \pkg{SimplifyDiagrams[a - b] === 0}, not by their appearance.
\item A symmetric function must be declared symmetric before it is used: \pkg{DeclareSymmetric} gives it the attribute
  \pkg{Orderless}, which only affects expressions evaluated afterwards.
\end{itemize}

\section*{References}\label{sec:references}
\addcontentsline{toc}{section}{References}

\begingroup
\small
\renewcommand{\section}[2]{}
\begin{thebibliography}{99}
\setlength{\itemsep}{1pt}
\bibitem{LuttingerWard1960} J.~M. Luttinger and J.~C. Ward, Ground-state energy of a many-fermion system.~II, Phys.\ Rev.\
  \textbf{118}, 1417 (1960).
\bibitem{Baym1962} G. Baym, Self-consistent approximations in many-body systems, Phys.\ Rev.\ \textbf{127}, 1391 (1962).
\bibitem{Bickers2004} N.~E. Bickers, Self-consistent many-body theory for condensed matter systems, in \emph{Theoretical
  Methods for Strongly Correlated Electrons}, edited by D. S\'en\'echal, A.-M. Tremblay and C. Bourbonnais (Springer,
  New York, 2004), p.~237.
\bibitem{DeDominicisMartin1964b} C. De Dominicis and P.~C. Martin, Stationary entropy principle and renormalization in normal and superfluid systems. II. Diagrammatic formulation, J.\ Math.\ Phys.\ \textbf{5}, 31 (1964).
\bibitem{Hedin1965} L. Hedin, New method for calculating the one-particle Green's function with application to the electron-gas problem, Phys.\ Rev.\ \textbf{139}, A796 (1965).
\bibitem{Stratonovich1958} R.~L. Stratonovich, On a method of calculating quantum distribution functions, Sov.\ Phys.\ Dokl.\ \textbf{2}, 416 (1958).
\bibitem{Hubbard1959} J. Hubbard, Calculation of partition functions, Phys.\ Rev.\ Lett.\ \textbf{3}, 77 (1959).
\bibitem{JonaLasinio1964} G. Jona-Lasinio, Relativistic field theories with symmetry-breaking solutions, Nuovo Cimento
  \textbf{34}, 1790 (1964).
\bibitem{ZinnJustin2002} J. Zinn-Justin, \emph{Quantum Field Theory and Critical Phenomena}, 4th ed.\ (Clarendon Press,
  Oxford, 2002).
\bibitem{Matsubara1955} T. Matsubara, A new approach to quantum-statistical mechanics, Prog.\ Theor.\ Phys.\ \textbf{14}, 351 (1955).
\bibitem{Mahan} G.~D. Mahan, \emph{Many-Particle Physics}, 3rd ed.\ (Kluwer Academic/Plenum Publishers, New York, 2000).
\bibitem{BruusFlensberg} H. Bruus and K. Flensberg, \emph{Many-Body Quantum Theory in Condensed Matter Physics: An
  Introduction} (Oxford University Press, Oxford, 2004).
\bibitem{Lindhard1954} J. Lindhard, On the properties of a gas of charged particles, Kgl.\ Danske Videnskab.\ Selskab,
  Mat.-Fys.\ Medd.\ \textbf{28}, no.~8 (1954).
\bibitem{FetterWalecka} A.~L. Fetter and J.~D. Walecka, \emph{Quantum Theory of Many-Particle Systems} (McGraw-Hill, New
  York, 1971; Dover, Mineola, 2003).
\bibitem{Migdal1958} A.~B. Migdal, Interaction between electrons and lattice vibrations in a normal metal, Sov.\ Phys.\
  JETP \textbf{7}, 996 (1958).
\bibitem{Cutkosky1960} R.~E. Cutkosky, Singularities and discontinuities of Feynman amplitudes, J.\ Math.\ Phys.\
  \textbf{1}, 429 (1960).
\bibitem{Weldon1983} H.~A. Weldon, Simple rules for discontinuities in finite-temperature field theory, Phys.\ Rev.\ D
  \textbf{28}, 2007 (1983).
\bibitem{KobesSemenoff1985} R.~L. Kobes and G.~W. Semenoff, Discontinuities of Green functions in field theory at finite
  temperature and density, Nucl.\ Phys.\ B \textbf{260}, 714 (1985).
\bibitem{KobesSemenoff1986} R.~L. Kobes and G.~W. Semenoff, Discontinuities of Green functions in field theory at finite
  temperature and density (II), Nucl.\ Phys.\ B \textbf{272}, 329 (1986).
\bibitem{LandsmanVanWeert1987} N.~P. Landsman and Ch.~G. van Weert, Real- and imaginary-time field theory at finite
  temperature and density, Phys.\ Rep.\ \textbf{145}, 141 (1987).
\bibitem{LeBellac1996} M. Le Bellac, \emph{Thermal Field Theory} (Cambridge University Press, Cambridge, 1996).
\bibitem{KapustaGale2006} J.~I. Kapusta and C. Gale, \emph{Finite-Temperature Field Theory: Principles and Applications},
  2nd ed.\ (Cambridge University Press, Cambridge, 2006).
\bibitem{tHooftVeltman1972} G. 't~Hooft and M. Veltman, Regularization and renormalization of gauge fields, Nucl.\
  Phys.\ B \textbf{44}, 189 (1972).
\bibitem{BolliniGiambiagi1972} C.~G. Bollini and J.~J. Giambiagi, Dimensional renormalization: the number of dimensions as
  a regularizing parameter, Nuovo Cimento B \textbf{12}, 20 (1972).
\bibitem{PeskinSchroeder} M.~E. Peskin and D.~V. Schroeder, \emph{An Introduction to Quantum Field Theory}
  (Addison-Wesley, Reading, 1995).
\bibitem{KleinertSchulteFrohlinde2001} H. Kleinert and V. Schulte-Frohlinde, \emph{Critical Properties of
  $\phi^4$-Theories} (World Scientific, Singapore, 2001).
\bibitem{ChetyrkinTkachov1981} K.~G. Chetyrkin and F.~V. Tkachov, Integration by parts: the algorithm to calculate
  $\beta$-functions in 4 loops, Nucl.\ Phys.\ B \textbf{192}, 159 (1981).
\bibitem{WilsonFisher1972} K.~G. Wilson and M.~E. Fisher, Critical exponents in 3.99 dimensions, Phys.\ Rev.\ Lett.\
  \textbf{28}, 240 (1972).
\bibitem{Schwinger1948} J. Schwinger, On quantum-electrodynamics and the magnetic moment of the electron, Phys.\ Rev.\
  \textbf{73}, 416 (1948).
\bibitem{Ward1950} J.~C. Ward, An identity in quantum electrodynamics, Phys.\ Rev.\ \textbf{78}, 182 (1950).
\bibitem{McKayPiperno2014} B.~D. McKay and A. Piperno, Practical graph isomorphism,~II, J.\ Symbolic Comput.\
  \textbf{60}, 94 (2014).
\bibitem{CornwallJackiwTomboulis1974} J.~M. Cornwall, R. Jackiw and E. Tomboulis, Effective action for composite
  operators, Phys.\ Rev.\ D \textbf{10}, 2428 (1974).
\bibitem{AlmbladhVonBarthVanLeeuwen1999} C.-O. Almbladh, U. von Barth and R. van Leeuwen, Variational total energies from
  $\Phi$- and $\Psi$-derivable theories, Int.\ J.\ Mod.\ Phys.\ B \textbf{13}, 535 (1999).
\bibitem{ColemanWeinberg1973} S. Coleman and E. Weinberg, Radiative corrections as the origin of spontaneous symmetry
  breaking, Phys.\ Rev.\ D \textbf{7}, 1888 (1973).
\bibitem{Potthoff2003} M. Potthoff, Self-energy-functional approach to systems of correlated electrons, Eur.\ Phys.\ J.\ B
  \textbf{32}, 429 (2003).
\bibitem{KozikFerreroGeorges2015} E. Kozik, M. Ferrero and A. Georges, Nonexistence of the Luttinger-Ward functional and
  misleading convergence of skeleton diagrammatic series for Hubbard-like models, Phys.\ Rev.\ Lett.\ \textbf{114}, 156402
  (2015).
\bibitem{RossiWerner2015} R. Rossi and F. Werner, Skeleton series and multivaluedness of the self-energy functional in zero
  space-time dimensions, J.\ Phys.\ A \textbf{48}, 485202 (2015), arXiv:1504.08320.
\bibitem{Schaefer2013} T. Sch\"afer, G. Rohringer, O. Gunnarsson, S. Ciuchi, G. Sangiovanni and A. Toschi, Divergent
  precursors of the Mott-Hubbard transition at the two-particle level, Phys.\ Rev.\ Lett.\ \textbf{110}, 246405 (2013).
\end{thebibliography}
\endgroup

\appendix
% ------------------------------------------------------------------ A
\section{Names in version 1}\label{sec:names}

Version 2 gave all the public functions explicit names:
\begin{center}
\scriptsize
\begin{tabular}{@{}ll@{\quad}ll@{}}
\toprule
version 1 & version 2 & version 1 & version 2\\
\midrule
\pkg{FD} & \pkg{FunctionalDerivative} & \pkg{IrreducibleQ} & \pkg{ParticleIrreducibleQ}\\
\pkg{Int} & \pkg{IntegralOver} & \pkg{IrreduciblePart} & \pkg{ParticleIrreducibleDiagrams}\\
\pkg{Dlt} & \pkg{DeltaFunction} & \pkg{SkeletonPart} & \pkg{SkeletonDiagrams}\\
\pkg{GProd} & \pkg{GrassmannProduct} & \pkg{ConnectedFrom1PI} & \pkg{ConnectedFunctionFromVertices}\\
\pkg{FDSimplify} & \pkg{SimplifyDiagrams} & \pkg{KG, KGamma} & \pkg{FullPropagator, VertexFunction}\\
\pkg{Terms} & \pkg{DiagramTerms} & \pkg{KPhi, KJ} & \pkg{MeanField, ExternalSource}\\
\pkg{ZeroSources} & \pkg{SetSourcesToZero} & \pkg{DiagramGraph} & \pkg{DrawDiagram}\\
\pkg{FieldExpectation} & \pkg{FieldExpectationValue} & \pkg{DeclareOdd} & \pkg{DeclareGrassmannOdd}\\
\pkg{WickExpectation} & \pkg{GaussianExpectationValue} & \pkg{DeclareGraded} & \pkg{DeclareGrassmannGraded}\\
\pkg{PerturbativeCorrelator} & \pkg{PerturbativeCorrelationFunction} & \pkg{DeclareDerivative} & \pkg{DeclareFunctionalDerivative}\\
\pkg{VacuumDiagrams} & \pkg{ConnectedVacuumDiagrams} & \pkg{DeclareIndexValues} & \pkg{DeclareDiscreteIndexValues}\\
\pkg{ConnectedPart} & \pkg{ConnectedDiagrams} & \pkg{DeclareDirected} & \pkg{DeclareDirectedPropagator}\\
\pkg{DropVacuum} & \pkg{DropVacuumBubbles} & \pkg{DeclareInteraction} & \pkg{DeclareInteractionLine}\\
\pkg{ZZ[d, k]} & \pkg{DummyVariable[d, k]} &  & \\
\bottomrule
\end{tabular}
\end{center}
In \pkg{ParticleIrreducibleQ} and \pkg{ParticleIrreducibleDiagrams} the external labels now come before the line heads.

% ------------------------------------------------------------------ B
\section{Index of functions}\label{sec:index}

\begin{center}
\footnotesize
\begin{tabular}{@{}l@{\hspace{0.4em}}r@{\hspace{0.75em}}l@{\hspace{0.4em}}r@{\hspace{0.75em}}l@{\hspace{0.4em}}r@{}}
\fnref{AmputateExternalLegs} & \pageref{fn:AmputateExternalLegs} & \fnref{DrawDiagram} & \pageref{fn:DrawDiagram} & \fnref{MassInsertion} & \pageref{fn:MassInsertion}\\
\fnref{AnalyticContinuation} & \pageref{fn:AnalyticContinuation} & \fnref{DropClosedLoops} & \pageref{fn:DropClosedLoops} & \fnref{MatsubaraSum} & \pageref{fn:MatsubaraSum}\\
\fnref{AnomalousDimensionFromZ} & \pageref{fn:AnomalousDimensionFromZ} & \fnref{DropVacuumBubbles} & \pageref{fn:DropVacuumBubbles} & \fnref{MeanField} & \pageref{fn:MeanField}\\
\fnref{BetaFunctionFromZ} & \pageref{fn:BetaFunctionFromZ} & \fnref{DummyVariable} & \pageref{fn:DummyVariable} & \fnref{MinkowskiLoopIntegral} & \pageref{fn:MinkowskiLoopIntegral}\\
\fnref{BetheSalpeterLadder} & \pageref{fn:BetheSalpeterLadder} & \fnref{DysonSeries} & \pageref{fn:DysonSeries} & \fnref{MomentumSquared} & \pageref{fn:MomentumSquared}\\
\fnref{BoseFunction} & \pageref{fn:BoseFunction} & \fnref{EffectiveAction1PI} & \pageref{fn:EffectiveAction1PI} & \fnref{MomentumSum} & \pageref{fn:MomentumSum}\\
\fnref{BubbleIntegral} & \pageref{fn:BubbleIntegral} & \fnref{EffectiveAction2PI} & \pageref{fn:EffectiveAction2PI} & \fnref{OperatorKernel} & \pageref{fn:OperatorKernel}\\
\fnref{CheckCuttingRules} & \pageref{fn:CheckCuttingRules} & \fnref{EnergyDelta} & \pageref{fn:EnergyDelta} & \fnref{OperatorTrace} & \pageref{fn:OperatorTrace}\\
\fnref{CheckEquation} & \pageref{fn:CheckEquation} & \fnref{EpsilonExpand} & \pageref{fn:EpsilonExpand} & \fnref{OrderPart} & \pageref{fn:OrderPart}\\
\fnref{CheckParquetEquations} & \pageref{fn:CheckParquetEquations} & \fnref{EvaluateMatsubara} & \pageref{fn:EvaluateMatsubara} & \fnref{ParameterIntegral} & \pageref{fn:ParameterIntegral}\\
\fnref{ClosedLoopCount} & \pageref{fn:ClosedLoopCount} & \fnref{ExactEquation} & \pageref{fn:ExactEquation} & \fnref{ParquetDecomposition} & \pageref{fn:ParquetDecomposition}\\
\fnref{CombineParameterIntegrals} & \pageref{fn:CombineParameterIntegrals} & \fnref{ExpandParameterIntegrals} & \pageref{fn:ExpandParameterIntegrals} & \fnref{ParticleIrreducibleDiagrams} & \pageref{fn:ParticleIrreducibleDiagrams}\\
\fnref{ConnectedCorrelationFunction} & \pageref{fn:ConnectedCorrelationFunction} & \fnref{ExpandTrace} & \pageref{fn:ExpandTrace} & \fnref{ParticleIrreducibleQ} & \pageref{fn:ParticleIrreducibleQ}\\
\fnref{ConnectedDiagrams} & \pageref{fn:ConnectedDiagrams} & \fnref{ExpressionToLaTeX} & \pageref{fn:ExpressionToLaTeX} & \fnref{PerturbativeCorrelationFunction} & \pageref{fn:PerturbativeCorrelationFunction}\\
\fnref{ConnectedFunctionFromVertices} & \pageref{fn:ConnectedFunctionFromVertices} & \fnref{ExternalSource} & \pageref{fn:ExternalSource} & \fnref{PerturbativeDensityVertex} & \pageref{fn:PerturbativeDensityVertex}\\
\fnref{ConnectedVacuumDiagrams} & \pageref{fn:ConnectedVacuumDiagrams} & \fnref{FermiFunction} & \pageref{fn:FermiFunction} & \fnref{PerturbativeFullVertex} & \pageref{fn:PerturbativeFullVertex}\\
\fnref{CuttingRuleImaginaryPart} & \pageref{fn:CuttingRuleImaginaryPart} & \fnref{FeynmanDenominator} & \pageref{fn:FeynmanDenominator} & \fnref{PolePart} & \pageref{fn:PolePart}\\
\fnref{DeclareAntisymmetric} & \pageref{fn:DeclareAntisymmetric} & \fnref{FeynmanParameter} & \pageref{fn:FeynmanParameter} & \fnref{QEDIntegrand} & \pageref{fn:QEDIntegrand}\\
\fnref{DeclareBlob} & \pageref{fn:DeclareBlob} & \fnref{FeynmanParameterIntegral} & \pageref{fn:FeynmanParameterIntegral} & \fnref{ReduceParameterIntegrals} & \pageref{fn:ReduceParameterIntegrals}\\
\fnref{DeclareCoupling} & \pageref{fn:DeclareCoupling} & \fnref{FieldDerivative} & \pageref{fn:FieldDerivative} & \fnref{ReducibleDiagrams} & \pageref{fn:ReducibleDiagrams}\\
\fnref{DeclareDirectedPropagator} & \pageref{fn:DeclareDirectedPropagator} & \fnref{FieldExpectationValue} & \pageref{fn:FieldExpectationValue} & \fnref{RenormalizationScale} & \pageref{fn:RenormalizationScale}\\
\fnref{DeclareDiscreteIndexValues} & \pageref{fn:DeclareDiscreteIndexValues} & \fnref{FourVector} & \pageref{fn:FourVector} & \fnref{ResetDeclarations} & \pageref{fn:ResetDeclarations}\\
\fnref{DeclareDisplayName} & \pageref{fn:DeclareDisplayName} & \fnref{FullPropagator} & \pageref{fn:FullPropagator} & \fnref{RetardedImaginaryPart} & \pageref{fn:RetardedImaginaryPart}\\
\fnref{DeclareField} & \pageref{fn:DeclareField} & \fnref{FunctionalDerivative} & \pageref{fn:FunctionalDerivative} & \fnref{ScalarProduct} & \pageref{fn:ScalarProduct}\\
\fnref{DeclareFunctionalDerivative} & \pageref{fn:DeclareFunctionalDerivative} & \fnref{FunctionalDiagramsUpdateCheck} & \pageref{fn:FunctionalDiagramsUpdateCheck} & \fnref{ScalarPropagator} & \pageref{fn:ScalarPropagator}\\
\fnref{DeclareGrassmannGraded} & \pageref{fn:DeclareGrassmannGraded} & \fnref{GammaMatrix} & \pageref{fn:GammaMatrix} & \fnref{SelfEnergy} & \pageref{fn:SelfEnergy}\\
\fnref{DeclareGrassmannOdd} & \pageref{fn:DeclareGrassmannOdd} & \fnref{GaussianExpectationValue} & \pageref{fn:GaussianExpectationValue} & \fnref{SelfEnergyFunctional} & \pageref{fn:SelfEnergyFunctional}\\
\fnref{DeclareGrassmannSource} & \pageref{fn:DeclareGrassmannSource} & \fnref{GFunction} & \pageref{fn:GFunction} & \fnref{SetSourcesToZero} & \pageref{fn:SetSourcesToZero}\\
\fnref{DeclareInteractionLine} & \pageref{fn:DeclareInteractionLine} & \fnref{GrassmannProduct} & \pageref{fn:GrassmannProduct} & \fnref{SimplexIntegral} & \pageref{fn:SimplexIntegral}\\
\fnref{DeclareLineStyle} & \pageref{fn:DeclareLineStyle} & \fnref{HubbardStratonovich} & \pageref{fn:HubbardStratonovich} & \fnref{SimplifyDiagrams} & \pageref{fn:SimplifyDiagrams}\\
\fnref{DeclareLocalVertex} & \pageref{fn:DeclareLocalVertex} & \fnref{InfraredRearrangedIntegral} & \pageref{fn:InfraredRearrangedIntegral} & \fnref{SkeletonDiagrams} & \pageref{fn:SkeletonDiagrams}\\
\fnref{DeclareMatsubaraFrequency} & \pageref{fn:DeclareMatsubaraFrequency} & \fnref{IntegralOver} & \pageref{fn:IntegralOver} & \fnref{SlashedMomentum} & \pageref{fn:SlashedMomentum}\\
\fnref{DeclareSource} & \pageref{fn:DeclareSource} & \fnref{IntegrateLoops} & \pageref{fn:IntegrateLoops} & \fnref{SpacetimeDimension} & \pageref{fn:SpacetimeDimension}\\
\fnref{DeclareSymmetric} & \pageref{fn:DeclareSymmetric} & \fnref{IntegrateOneLoop} & \pageref{fn:IntegrateOneLoop} & \fnref{SpinorChain} & \pageref{fn:SpinorChain}\\
\fnref{DeltaFunction} & \pageref{fn:DeltaFunction} & \fnref{IntegrateOutField} & \pageref{fn:IntegrateOutField} & \fnref{SuperfieldComponents} & \pageref{fn:SuperfieldComponents}\\
\fnref{DiagramIntegrand} & \pageref{fn:DiagramIntegrand} & \fnref{IrreducibleDiagrams} & \pageref{fn:IrreducibleDiagrams} & \fnref{TadpoleIntegral} & \pageref{fn:TadpoleIntegral}\\
\fnref{DiagramLoopIntegral} & \pageref{fn:DiagramLoopIntegral} & \fnref{IrreducibleVertex} & \pageref{fn:IrreducibleVertex} & \fnref{Temperature} & \pageref{fn:Temperature}\\
\fnref{DiagramPlot} & \pageref{fn:DiagramPlot} & \fnref{LabelFunction} & \pageref{fn:LabelFunction} & \fnref{TraceLog} & \pageref{fn:TraceLog}\\
\fnref{DiagramTerms} & \pageref{fn:DiagramTerms} & \fnref{LegendrePropagator} & \pageref{fn:LegendrePropagator} & \fnref{TwoParticleReducibleQ} & \pageref{fn:TwoParticleReducibleQ}\\
\fnref{DiagramToTikZ} & \pageref{fn:DiagramToTikZ} & \fnref{LegendreTransform} & \pageref{fn:LegendreTransform} & \fnref{VertexFormFactors} & \pageref{fn:VertexFormFactors}\\
\fnref{DifferentiateUnderIntegral} & \pageref{fn:DifferentiateUnderIntegral} & \fnref{LoopFrequency} & \pageref{fn:LoopFrequency} & \fnref{VertexFunction} & \pageref{fn:VertexFunction}\\
\fnref{DiracOrder} & \pageref{fn:DiracOrder} & \fnref{LoopIntegral} & \pageref{fn:LoopIntegral} & \fnref{ZeroDimensionalValue} & \pageref{fn:ZeroDimensionalValue}\\
\fnref{DiracProduct} & \pageref{fn:DiracProduct} & \fnref{LoopMomentum} & \pageref{fn:LoopMomentum} & \fnref{ZeroTemperatureLimit} & \pageref{fn:ZeroTemperatureLimit}\\
\fnref{DiracSimplify} & \pageref{fn:DiracSimplify} & \fnref{LorentzMetric} & \pageref{fn:LorentzMetric} & &\\
\fnref{DiracTrace} & \pageref{fn:DiracTrace} & \fnref{LuttingerWardFunctional} & \pageref{fn:LuttingerWardFunctional} & &\\

\end{tabular}
\end{center}

\end{document}
